\chapter{Linear algebra, reprise}
\label{chap:linrep}

We come back to linear algebra, with the task of studying slightly
more sophisticated constructions than those reviewed in
\Cref{chap:linalg}. My main goals are to discuss some aspects of
multilinear algebra, tensor products, and Hom, with special attention
devoted to dual modules. Along the way we will meet several other
interesting notions, including a first taste of projective and
injective modules and of Tor and Ext.

As in \Cref{chap:linalg}, I will attempt to deal with the more general
case of $R$-modules, specializing to vector spaces only when this
provides a drastic simplification or conforms to inveterate
habits. Covering the material at this level of generality gives me an
excellent excuse to complement the preliminary notions described in
the very first chapter of this book.

\section{Preliminaries, reprise}
\label{sec:linrep:preliminaries-reprise}


\subsection{Functors}
\label{subsec:linrep:functors}
In this book I have given an unusually early introduction to the
notion of category, motivated by the near omnipresence of universal
problems at the very foundation of algebra. The hope was that the
unifying viewpoint offered by categories would help the reader in the
task of organizing the information carried by the basic constructions
we have encountered.  However, so far we have never seriously had to
`go from one category to another', and therefore I have not felt
compelled to complete the elementary presentation of categories by
discussing more formally how this is done. The careful reader has
likely seen it happen between the lines, at least in noninteresting
ways: for example, if we strip a ring of its multiplicative structure,
we are left with an abelian group; and if a field $F$ is an extension
of a field $k$, then we may view $F$ as a $k$-vector space: thus,
there are evident ways to go from $\categ{Ring}$ to $\categ{Ab}$ or
from the category $\categ{Fld}_k$ of extensions of $k$ to
$k$-$\categ{Vect}$. More interesting examples are the group of units
of a ring, the homology of a complex, the automorphism group of a
Galois field extension (and note that this latter `turns arrows
around'; see below), etc. Implicitly, we have in fact encountered many
instances of such `functions between categories'.  The notion of
functor formalizes these operations. A functor between two categories
$\categ{C}$, $\categ{D}$ will `send objects of $\categ{C}$ to objects
of $\categ{D}$', and since categories carry the information of
morphisms, functors will have to deal with this information as well.
\begin{defn}{}{defn:linrep:covariant-functor}
  Let $\categ{C}$, $\categ{D}$ be two categories. A \emph{covariant
    functor} $F : \categ{C} \to \categ{D}$ is an assignment of an
  object $F(A) \in \mathrm{Obj}(\categ{D})$ for every
  $A \in \mathrm{Obj}(\categ{C})$ and of a function
  $\Hom_{\categ{C}}(A, B) \to \Hom_{\categ{D}}(F(A), F(B))$ for every
  pair of objects $A$, $B$ in $\categ{C}$; this function is also
  denoted\footnote{This minor abuse of notation is convenient and
    standard.} $F$ and must preserve identities and compositions. That
  is,
  \[
    F(1_A) = 1_{F(A)} \quad \forall A \in \mathrm{Obj}(\categ{C}),
  \]
  and
  \[
    F(\beta \circ \alpha) = F(\beta) \circ F(\alpha) \quad \forall A,
    B, C \in \mathrm{Obj}(\categ{C}),\; \forall \alpha \in
    \Hom_{\categ{C}}(A, B),\; \forall \beta \in \Hom_{\categ{C}}(B,
    C).
  \]
  A \emph{contravariant functor} $\categ{C} \to \categ{D}$ is a
  covariant functor $\categ{C}^{\mathrm{op}} \to \categ{D}$ from the
  opposite category.
\end{defn}
The opposite category $\categ{C}^{\mathrm{op}}$ is obtained from
$\categ{C}$ by `reversing the arrows'; cf.\
\Cref{exer:prelim:op-categ-def}. The fact that contravariant functors
$G : \categ{C} \to \categ{D}$ (that is, covariant functors
$\categ{C}^{\mathrm{op}} \to \categ{D}$) preserve compositions means
that $\forall A, B, C \in \mathrm{Obj}(\categ{C})$,
$\forall \alpha \in \Hom_{\categ{C}}(A, B)$,
$\forall \beta \in \Hom_{\categ{C}}(B, C)$:
\[
  G(\beta \circ \alpha) = G(\alpha) \circ G(\beta).
\]
Note the switch in the order of composition. A covariant functor
`preserves' arrows, in the sense that every diagram in $\categ{C}$,
say a commutative diagram in $\categ{C}$ with vertices $A, B, C, D$
and arrows $\alpha_1, \alpha_2, \alpha_3, \beta, \gamma$, is mapped by
a covariant functor $F : \categ{C} \to \categ{D}$ to a commutative
diagram with arrows in the same directions. A contravariant functor
$G : \categ{C} \to \categ{D}$ `reverses' the arrows. In both cases,
commutative diagrams are sent to commutative diagrams.  In many
contexts, contravariant functors are called \emph{presheaves}.  Thus,
a presheaf of sets on a category $\categ{C}$ is simply a contravariant
functor $\categ{C} \to \categ{Set}$.  If $\Hom_{\categ{C}}(A, B)$
carries more structure, we may require functors to preserve this
structure. For example, we have seen that for
$\categ{C} = R\text{-}\categ{Mod}$, the category of left-modules over
a ring $R$, $\Hom_{R\text{-}\mathrm{Mod}}(A, B)$ is itself an abelian
group (end of \Cref{subsec:rings:the-category-r-mod}); a functor
$F : R\text{-}\categ{Mod} \to S\text{-}\categ{Mod}$ is \emph{additive}
if it preserves this structure, that is, if the corresponding function
\[
  \Hom_{R\text{-}\mathrm{Mod}}(A, B) \to
  \Hom_{S\text{-}\mathrm{Mod}}(F(A), F(B))
\]
is a homomorphism of abelian groups (for all $R$-modules $A$, $B$).

\subsection{Examples of functors}
\label{subsec:linrep:examples-of-functors}
In practice, an assignment of objects of a category to objects of
another category often comes with a (psychologically) `natural'
companion assignment for morphisms, which is evidently covariant or
contravariant. We say that this assignment is `functorial' if it
preserves compositions. All the simple-minded operations mentioned at
the beginning of \Cref{subsec:linrep:functors} (such as viewing a
ring as an abelian group under addition) are covariant functors: for
example, the fact that a homomorphism of rings is in particular a
homomorphism of the underlying abelian groups amounts to the statement
that labeling each ring with its underlying abelian group defines a
covariant functor $\categ{Ring} \to \categ{Ab}$. Since such functors
are obtained by `forgetting' part of the structure of a given object,
they are memorably called \emph{forgetful functors}. I have already
mentioned briefly a few slightly more imaginative examples; here they
are again.
\begin{exam}{}{exam:linrep:units-functor}
  If $R$ is a ring, we have denoted by $R^\ast$ the group of units in
  $R$; every ring homomorphism $R \to S$ induces a group homomorphism
  $R^\ast \to S^\ast$, and this assignment is compatible with
  compositions; therefore this operation defines a covariant functor
  $\categ{Ring} \to \categ{Grp}$.
\end{exam}
\begin{exam}{}{exam:linrep:automorphism-group-contravariant-functor}
  The operation $\Aut_k(-)$ from Galois field extensions to groups is
  contravariantly functorial. Indeed, if $k \subseteq E \subseteq F$
  is viewed as a morphism of two Galois extensions ($k \subseteq E$ to
  $k \subseteq F$), we have a corresponding group homomorphism
  $\Aut_k(F) \to \Aut_k(E)$ defined by restriction:
  $\phi \mapsto \phi|_E$; the latter maps $E$ to $E$ by virtue of the
  Galois condition. The composition of two restrictions is a
  restriction, so the assignment is indeed functorial.
\end{exam}
\begin{exam}{}{exam:linrep:spectrum-contravariant-functor}
  In \Cref{subsec:rings:prime-and-maximal-ideals} I have defined
  the spectrum of a commutative ring $R$, $\mathrm{Spec}\,R$, as the
  set of prime ideals of $R$. If $R$, $S$ are commutative rings and
  $\phi : R \to S$ is a ring homomorphism, then the inverse image
  $\phi^{-1}(\mathfrak{p})$ of a prime ideal $\mathfrak{p}$ is a prime
  ideal; thus, $\phi$ induces a set-function
  $\phi^\ast : \mathrm{Spec}(S) \to \mathrm{Spec}(R)$. This assignment
  is clearly functorial, so we can view $\mathrm{Spec}$ as a
  contravariant functor from the category of commutative rings to
  $\categ{Set}$. The reader will assign a topology to
  $\mathrm{Spec}(R)$
  (\Cref{exer:linrep:spec-zariski-topology-is-functor}), making this
  example (even) more interesting.
\end{exam}
\begin{exam}{}{exam:linrep:presheaf-of-sets}
  If $S$ is a set, recall (\Cref{exam:prelims:power-set-category})
  that one can construct a category $\hat{S}$ whose objects are
  subsets of $S$ and where morphisms correspond to inclusions of
  subsets. A presheaf of sets on $S$ is a contravariant functor
  $\hat{S} \to \categ{Set}$. The prototypical example consists of the
  functor which associates to $U \subseteq S$ the set of functions
  from $U$ to a fixed set; the inclusion $U \subseteq V$ is mapped to
  the restriction $\rho \mapsto \rho|_U$. More interesting (and very
  useful) examples may be obtained by considering richer
  structures. For instance, if $T$ is a topological space, one can
  define a category whose objects are the open sets in $T$ and whose
  morphisms are inclusions; this leads to the notion of presheaf on a
  topological space. Standard notions such as `(the groups of)
  continuous complex-valued functions on open sets of $T$' are most
  naturally viewed as the datum of a presheaf of abelian groups on
  $T$. Such concrete examples often satisfy additional hypotheses,
  which make them \emph{sheaves}, but that is a topic for another
  book. In a sense, studying a branch of geometry (e.g., complex
  differential geometry) amounts to studying spaces endowed with a
  suitable sheaf (e.g., the sheaf of complex differentiable
  functions).
\end{exam}
An important class of functors that are available on any category may
be obtained as follows, and this will be a source of examples in the
next few sections. Let $\categ{C}$ be a category, and let $X$ be an
object of $\categ{C}$. Then the assignments
\[
  A \mapsto \Hom_{\categ{C}}(X, A), \qquad A \mapsto
  \Hom_{\categ{C}}(A, X)
\]
define, respectively, covariant and contravariant functors
$\categ{C} \to \categ{Set}$. These are denoted
$\Hom_{\categ{C}}(X, -)$ and $\Hom_{\categ{C}}(-, X)$,
respectively. For example, if
$A \xrightarrow{\alpha} B \xrightarrow{\beta} C$ is a diagram in
$\categ{C}$, every $\alpha : A \to B$ determines a function
$\Hom_{\categ{C}}(B, X) \to \Hom_{\categ{C}}(A, X)$ defined by mapping
$\xi_B$ to $\xi_A \coloneq \xi_B \circ \alpha$, that is, by requiring
the obvious triangle to be commutative. The associativity of
composition,
$\xi_C \circ (\beta \circ \alpha) = (\xi_C \circ \beta) \circ \alpha$,
expresses precisely the contravariant functoriality of this
assignment.  The functor $\Hom_{\categ{C}}(-, X)$ is often denoted
$h^X$ for short, if $\categ{C}$ is understood.  The brave reader may
contemplate the fact that this technique associates to each object $X$
of a category $\categ{C}$ a presheaf of sets $h^X$ on
$\categ{C}$. This raises natural and important questions, such as how
to tell whether a given presheaf of sets $F$ in fact equals $h^X$ for
some object $X$ of $\categ{C}$ (this makes $F$ \emph{representable})
or whether an object $X$ may be reconstructed from the corresponding
presheaf $h^X$. For example, if $\categ{Top}$ is the category of
topological spaces and $p$ denotes the final object of $\categ{Top}$
(that is, the one-point topological space), then $h^X(p)$ is nothing
but $X$, viewed as a set. For this reason, in such `geometric'
contexts (for example, in algebraic geometry) $h^X(A)$ is called the
`set of $A$-points of $X$'; $h^X$ is the \emph{functor of points} of
$X$.

\subsection{When are two categories `equivalent'?}
\label{subsec:linrep:equivalence-of-categories}
Essentially every notion of `isomorphism' encountered so far has
boiled down to `structure-preserving bijection', and the reader may
well expect that the natural notion identifying two categories should
be drawn from the same model: a functor matching objects of one
category precisely with objects of another, preserving the structure
of morphisms.  One could certainly introduce such a notion, but it
would be exceedingly restrictive. Recall that solutions to universal
problems, that is, just about every construction we have run across,
are only defined up to isomorphism
(\Cref{prop:prelims:initial-final-unique}). Requiring solutions in one
context to match exactly with solutions in another would be
problematic. The structure of an object in a category is adequately
carried by its isomorphism class\footnote{One should not take this
  viewpoint too far.  It is tempting to think of isomorphic objects in
  a category as `the same', but this is also problematic and does not
  lead (as far as I know) to a workable alternative. For example, some
  objects in a category may have `more automorphisms' than others, and
  this information would be discarded if we simply chopped up the
  category into isomorphism classes, draconically promoting all
  isomorphisms to identity morphisms.}, and a natural notion of
`equivalence' of categories should aim at matching isomorphism
classes, rather than individual objects. The \emph{morphisms} are a
more essential piece of information; the quality of a functor is first
of all measured on how it acts on morphisms.
\begin{defn}{}{defn:linrep:faithful-full-functor}
  Let $\categ{C}$, $\categ{D}$ be two categories. A covariant functor
  $F : \categ{C} \to \categ{D}$ is \emph{faithful} if for all objects
  $A$, $B$ of $\categ{C}$, the induced function
  \[
    \Hom_{\categ{C}}(A, B) \to \Hom_{\categ{D}}(F(A), F(B))
  \]
  is injective; it is \emph{full} if this function is surjective, for
  all objects $A$, $B$.
\end{defn}
`Fully faithful' functors $F : \categ{C} \to \categ{D}$ are bijective
on morphisms, and it follows easily that they preserve isomorphism
classes: $A \cong B$ in $\categ{C}$ if and only if $F(A) \cong F(B)$
in $\categ{D}$ (\Cref{exer:linrep:fully-faithful-reflects-iso}). For
all intents and purposes, a fully faithful functor
$F : \categ{C} \to \categ{D}$ identifies $\categ{C}$ with a
full\footnote{See \Cref{exer:prelim:subcat} for the definition of
  `full' subcategory.} subcategory of $\categ{D}$, even if `on
objects' $F$ may be neither injective nor surjective onto that
subcategory. This idea is captured by the following definition:
\begin{defn}{}{defn:linrep:equivalence-of-categories}
  A covariant functor $F : \categ{C} \to \categ{D}$ is an
  \emph{equivalence of categories} if it is fully faithful and
  `essentially surjective'; that is, $F$ induces bijections on
  Hom-sets, and for every object $Y$ of $\categ{D}$ there is an object
  $X$ of $\categ{C}$ such that $F(X) \cong Y$ in $\categ{D}$.
\end{defn}
Two categories $\categ{C}$, $\categ{D}$ are \emph{equivalent} if there
is a functor $F : \categ{C} \to \categ{D}$ satisfying the requirements
in \Cref{defn:linrep:equivalence-of-categories}. It is not at all immediate that this is indeed
an equivalence relation (why is it symmetric?), but this turns out to
be the case. In fact, an equivalence of categories $F$ has an
`inverse', in a suitably weakened sense: a functor
$G : \categ{D} \to \categ{C}$ such that there are `natural
isomorphisms' (see \Cref{subsec:linrep:comparing-functors}) from
$F \circ G$ and $G \circ F$ to the identity. Since we will not use
these notions too extensively, I will let the interested reader
research the issue and find proofs.
\begin{exam}{}{exam:linrep:matrix-category-equivalent-to-vector-spaces}
  (Cf.\ \Cref{exer:prelim:categ-matrix}.) Construct a category by
  taking the objects to be nonneg\-ative integers and $\Hom(m, n)$ to
  be the set of $n \times m$ matrices with entries in a field $k$,
  with composition defined by product of matrices (and suitable care
  concerning matrices with no rows or columns). The resulting category
  is equivalent to the category of finite-dimensional $k$-vector
  spaces. Indeed, we obtain a functor from the former to the latter by
  sending $n$ to the vector space $k^n$ (endowed with the standard
  basis) and each matrix to the corresponding linear map. This functor
  is clearly fully faithful, and it is essentially surjective because
  vector spaces are classified by their dimension. This example makes
  (more) precise the heuristic considerations at the end of
  \Cref{subsec:linalg:matrices}.
\end{exam}
\begin{exam}{}{exam:linrep:category-of-affine-algebraic-sets}
  Recall (\Cref{defn:fields:coordinate-ring}) that the coordinate ring of an
  affine algebraic set over a field $K$ is a reduced, commutative,
  finite-type $K$-algebra. We can define the category
  $K\text{-}\categ{Aff}$ of affine $K$-algebraic sets by prescribing
  that the objects be algebraic subsets of some affine $K$-space and
  defining\footnote{The switch in the order of $S$ and $T$ is
    reasonable: as the reader has verified in
    \Cref{exer:fields:poly-fns-iso-to-coordinate-ring}, $K[S]$ may be
    interpreted as the $K$-algebra of `polynomial functions' on $S$;
    any function $\phi : S \to T$ determines a pull-back of functions,
    $K[T] \to K[S]$.}
  $\Hom_{K\text{-}\mathrm{Aff}}(S, T) =
  \Hom_{K\text{-}\mathrm{Alg}}(K[T], K[S])$.

  Thus, $K\text{-}\categ{Aff}$ is defined in such a way that the
  functor
  $K\text{-}\categ{Aff}^{\mathrm{op}} \to K\text{-}\categ{Alg}$ that
  maps an affine algebraic set $S$ to its coordinate ring $K[S]$ is an
  equivalence of the opposite category of $K\text{-}\categ{Aff}$ with
  the subcategory of reduced, commutative, finite-type $K$-algebra.
\end{exam}

\subsection{Limits and colimits}
\label{subsec:linrep:limits-and-colimits}
The various universal properties encountered along the way in this
book are all particular cases of the notion of categorical
\emph{limit}, which is worth mentioning explicitly. Let
$\mathscr{F} : \categ{I} \to \categ{C}$ be a covariant functor, where
one thinks of $\categ{I}$ as a category `of indices'. The \emph{limit}
of $\mathscr{F}$ is (if it exists) an object $L$ of $\categ{C}$,
endowed with morphisms $\lambda_I : L \to \mathscr{F}(I)$ for all
objects $I$ of $\categ{I}$, satisfying the following properties.
\begin{itemize}
\item If $\alpha : I \to J$ is a morphism in $\categ{I}$, then
  $\lambda_J = \mathscr{F}(\alpha) \circ \lambda_I$.
\item $L$ is final with respect to this property: that is, if $M$ is
  another object, endowed with morphisms $\mu_I$, also satisfying the
  previous requirement, then there exists a unique morphism $M \to L$
  making all relevant diagrams commute\footnote{Any such datum of $M$
    and morphisms is called a \emph{cone} for
    $\mathscr{F}$. Therefore, the limit of $\mathscr{F}$ is a
    `universal cone': it is a final object in the (easily defined)
    category of cones of $\mathscr{F}$.}.
\end{itemize}
The limit $L$ (if it exists) is unique up to isomorphism, as is every
notion defined by a universal property. It is denoted
$\varprojlim \mathscr{F}$. The `left-pointing' arrow reminds us that
$L$ stands `before' all objects of $\categ{C}$ indexed by $\categ{I}$
via $\mathscr{F}$ (but it is---up to isomorphisms---the `last' object
that may do so). This notion is also called \emph{inverse}, or
\emph{projective}, limit.
\begin{exam}{(Products)}{exam:linrep:product-as-limit-over-discrete-category}
  Let $\categ{I}$ be the `discrete category' consisting of two objects
  $\mathbf{1}$, $\mathbf{2}$, with only identity
  morphisms\footnote{This latter prescription is what makes the
    category `discrete'; cf.\ \Cref{exam:prelims:category-leq}.}, and
  let $\mathscr{A}$ be a functor from $\categ{I}$ to any category
  $\categ{C}$; let $A_1 = \mathscr{A}(\mathbf{1})$,
  $A_2 = \mathscr{A}(\mathbf{2})$ be the two objects of $\categ{C}$
  `indexed' by $\categ{I}$. Then $\varprojlim \mathscr{A}$ is nothing
  but the product of $A_1$ and $A_2$ in $\categ{C}$, as defined in
  \Cref{subsec:prelims:products}: a limit exists if and only if a
  product of $A_1$ and $A_2$ exists in $\categ{C}$.

  We can similarly define the product of any (possibly infinite)
  family of objects in a category as the limit over the corresponding
  discrete indexing category, provided of course that this limit
  exists.
\end{exam}
The limit notion is a little more interesting if the indexing category
$\categ{I}$ carries more structure.
\begin{exam}{(Equalizers and kernels)}{exam:linrep:equalizer-as-limit}
  Let $\categ{I}$ again be a category with two objects $\mathbf{1}$,
  $\mathbf{2}$, but assume that morphisms look like this: add to the
  discrete category two `parallel' morphisms $\alpha$, $\beta$ from
  $\mathbf{2}$ to $\mathbf{1}$. A functor
  $\mathscr{K} : \categ{I} \to \categ{C}$ amounts to the choice of two
  objects $A_1$, $A_2$ in $\categ{C}$ and two parallel morphisms
  between them. Limits of such functors are called
  \emph{equalizers}. For a concrete example, assume
  $\categ{C} = R\text{-}\categ{Mod}$ is the category of $R$-modules
  for some ring $R$; let $\phi : A_2 \to A_1$ be a homomorphism, and
  choose $\mathscr{K}$ as above, with $\mathscr{K}(\alpha) = \phi$ and
  $\mathscr{K}(\beta) = $ the zero-morphism. Then
  $\varprojlim \mathscr{K}$ is nothing but the kernel of $\phi$, as
  the reader will verify
  (\Cref{exer:linrep:construction-recovers-kernel}).
\end{exam}
\begin{exam}{(Limits over chains)}{exam:linrep:limit-over-a-chain}
  In another typical situation, $\categ{I}$ may consist of a totally
  ordered set, for example:
  \[
    \cdots \to \mathbf{4} \to \mathbf{3} \to \mathbf{2} \to \mathbf{1}
  \]
  (that is, the objects are $\mathbf{i}$, for all positive integers
  $i$, and there is a unique morphism $\mathbf{i} \to \mathbf{j}$
  whenever $i \geq j$; I am only drawing the morphisms
  $\mathbf{j+1} \to \mathbf{j}$). Choosing
  $\mathscr{F} : \categ{I} \to \categ{C}$ is equivalent to choosing
  objects $A_i$ of $\categ{C}$ for all positive integers $i$ and
  morphisms $\varphi_{ji} : A_i \to A_j$ for all $i \geq j$, with the
  requirement that $\varphi_{ii} = 1_{A_i}$ and
  $\varphi_{kj} \circ \varphi_{ji} = \varphi_{ki}$ for all
  $i \geq j \geq k$. That is, the choice of $\mathscr{F}$ amounts to
  the choice of a diagram
  \[
    \cdots \xrightarrow{\varphi_{45}} A_4 \xrightarrow{\varphi_{34}}
    A_3 \xrightarrow{\varphi_{23}} A_2 \xrightarrow{\varphi_{12}} A_1
  \]
  in $\categ{C}$. An inverse limit $\varprojlim \mathscr{F}$ (which
  may also be denoted $\varprojlim_i A_i$, when the morphisms
  $\varphi_{ji}$ are evident from the context) is then an object $A$
  endowed with morphisms $\varphi_i : A \to A_i$ such that the whole
  diagram commutes and such that any other object satisfying this
  requirement factors uniquely through $A$.
\end{exam}
Such limits exist in many standard situations. For example, let
$\categ{C} = R\text{-}\categ{Mod}$ be the category of left-modules
over a fixed ring $R$, and let $A_i$, $\varphi_{ji}$ be as above.
\begin{prop}{}{prop:linrep:inverse-limit-of-modules-exists}
  The limit $\varprojlim_i A_i$ exists in $R\text{-}\categ{Mod}$.
\end{prop}
\begin{proof}
  The product $\prod_i A_i$ consists of arbitrary sequences
  $(a_i)_{i > 0}$ of elements $a_i \in A_i$. Say that a sequence
  $(a_i)_{i > 0}$ is \emph{coherent} if for all $i > 0$ we have
  $a_i = \varphi_{i,i+1}(a_{i+1})$. Coherent sequences form an
  $R$-submodule $A$ of $\prod_i A_i$; the canonical projections
  restrict to $R$-module homomorphisms $\varphi_i : A \to A_i$. The
  reader will check that $A$ is a limit $\varprojlim_i A_i$
  (\Cref{exer:linrep:construction-is-inverse-limit}).
\end{proof}
This example easily generalizes to families indexed by more general
posets.  The `dual notion' to limit is the \emph{colimit} of a functor
$\mathscr{F} : \categ{I} \to \categ{C}$. The colimit is an object $C$
of $\categ{C}$, endowed with morphisms
$\gamma_I : \mathscr{F}(I) \to C$ for all objects $I$ of $\categ{I}$,
such that $\gamma_I = \gamma_J \circ \mathscr{F}(\alpha)$ for all
$\alpha : I \to J$ and that $C$ is initial with respect to this
requirement.

Again, we have already encountered several instances of this
construction: for example, the coproduct of two objects $A_1$, $A_2$
of a category is (if it exists) the colimit over the discrete category
with two objects, just as their product is the corresponding
limit. Cokernels are colimits, just as kernels are limits (cf.\
\Cref{exam:linrep:equalizer-as-limit}).

The colimit of $\mathscr{F}$ is (not surprisingly) denoted
$\varinjlim \mathscr{F}$ and is called the \emph{direct}, or
\emph{injective}, limit of $\mathscr{F}$. For a typical situation
consider again the case of a totally ordered set $\categ{I}$, for
example:
\[
  \mathbf{1} \to \mathbf{2} \to \mathbf{3} \to \mathbf{4} \to \cdots.
\]
A functor $\mathscr{F} : \categ{I} \to \categ{C}$ consists of the
choice of objects and morphisms
\[
  A_1 \xrightarrow{\psi_{12}} A_2 \xrightarrow{\psi_{23}} A_3
  \xrightarrow{\psi_{34}} A_4 \xrightarrow{\psi_{45}} \cdots
\]
and the direct limit $\varinjlim_i A_i$ will be an object $A$ with
morphisms $\psi_i : A_i \to A$ such that the relevant diagram commutes
and such that $A$ is initial with respect to this requirement.
\begin{exam}{}{exam:linrep:nested-sequence-of-sets-direct-limit}
  If $\categ{C} = \categ{Set}$ and all the $\psi_{ij}$ are injective,
  we are talking about a `nested sequence of sets':
  \[
    A_1 \subseteq A_2 \subseteq A_3 \subseteq A_4 \subseteq \cdots;
  \]
  the direct limit of this sequence would be the `infinite union'
  $\bigcup_i A_i$.
\end{exam}
I have shamelessly relied on the reader's intuition and used this
notion already, for example when I constructed the algebraic closure
of a field (in \Cref{subsec:fields:algebraic-closure}) or whenever
I have mentioned polynomial rings with `infinitely many variables'
(e.g., in \Cref{exam:rings:infinite-poly-ring-fg}). More formally,
$\bigcup_i A_i$ consists of equivalence classes of pairs $(i, a_i)$,
where $a_i \in A_i$ and $(i, a_i)$ is equivalent to $(j, a_j)$ for
$i \leq j$ if $a_j = \psi_{ij}(a_i)$.

In the context of $R$-modules, one can consider the direct sum
$D = \bigoplus_i A_i$ and the submodule $K \subseteq D$ generated by
elements of the form $a_i - \psi_{ij}(a_i)$ for all $i \leq j$ (where
each $A_i$ is identified with its image in $D$). Then the quotient
$D/K$ satisfies the needed universal property, as the reader will
check. In fact, it is no harder to perform these constructions for
more general posets
(\Cref{exer:linrep:colimits-in-set-and-mod-directed-poset}). They are
somewhat more explicit and better behaved in the case of
\emph{directed sets}: partially ordered sets $\categ{I}$ such that
$\forall i, j \in \categ{I}$ there exists $k \in \categ{I}$ such that
$i \leq k$, $j \leq k$.

Several standard constructions rely on a direct limit: for example,
\emph{germs} of functions (or, more generally, `stalks of presheaves')
are defined by means of a direct limit.

\subsection{Comparing functors}
\label{subsec:linrep:comparing-functors}
Having introduced functors as the natural notion of `morphisms between
categories', the next natural step is to consider `morphisms between
functors' and other ways to compare two given functors. A detailed
description of these notions is not needed in this book, but I want to
mention briefly a few concepts and essential remarks, again because
they offer a unifying viewpoint.
\begin{defn}{}{defn:linrep:natural-transformation}
  Let $\categ{C}$, $\categ{D}$ be categories, and let $F$, $G$ be
  (say, covariant) functors $\categ{C} \to \categ{D}$. A \emph{natural
    transformation}\footnote{The symbol $\Rightarrow$ is more standard
    than $\leadsto$, but it looks a little too much like the logical
    connective $\implies$ for my taste.} $F \leadsto G$ is the datum
  of a morphism $\nu_X : F(X) \to G(X)$ in $\categ{D}$ for every
  object $X$ in $\categ{C}$, such that $\forall \alpha : X \to Y$ in
  $\categ{C}$ the diagram
  \[
    \begin{array}{ccc}
      F(X) & \xrightarrow{F(\alpha)} & F(Y) \\
      \nu_X \downarrow & & \downarrow \nu_Y \\
      G(X) & \xrightarrow{G(\alpha)} & G(Y)
    \end{array}
  \]
  commutes. A \emph{natural isomorphism} is a natural transformation
  $\nu$ such that $\nu_X$ is an isomorphism for every $X$.
\end{defn}
Natural transformations arise in many contexts, for example when
comparing different notions of homology or cohomology in topology. The
reader is likely to have run into the Hurewicz homomorphism (from
$\pi_1$ to $H_1$); it is a natural transformation. In fact, any
statement such as `there is a natural (or canonical)
homomorphism\dots' is likely hiding a natural transformation between
two functors. This technical meaning of the word `natural' matches, as
a rule, its psychological use.

A particularly basic and important example is the notion of
\emph{adjoint functors}, which I will mention at a rather informal
level, leaving to the inextinguishable reader the task of making it
formal by suitable use of natural transformations. Let $\categ{C}$,
$\categ{D}$ be categories, and let $F : \categ{C} \to \categ{D}$,
$G : \categ{D} \to \categ{C}$ be functors. We say that $F$ and $G$ are
\emph{adjoint} (and we say that $G$ is \emph{right-adjoint} to $F$ and
$F$ is \emph{left-adjoint} to $G$) if there are natural isomorphisms
\[
  \Hom_{\categ{C}}(X, G(Y)) \xrightarrow{\;\sim\;}
  \Hom_{\categ{D}}(F(X), Y)
\]
for all objects $X$ of $\categ{C}$ and $Y$ of $\categ{D}$. (More
precisely, there should be a natural isomorphism of `bifunctors'
$\categ{C}^{\mathrm{op}} \times \categ{D} \to \categ{Set}$:
$\Hom_{\categ{C}}(\mathord{-}, G(\mathord{-})) \leadsto
\Hom_{\categ{D}}(F(\mathord{-}), \mathord{-})$.) Once again, several
constructions we have encountered along the way may be recast in these
terms, and we will run into more instances of this situation in this
chapter.
\begin{exam}{}{exam:linrep:free-group-universal-property-as-natural-iso}
  The construction of the free group on a given set
  (\Cref{sec:groups:free-groups}) is concocted so that giving a
  set-function from a set $A$ to a group $G$ is `the same as' giving a
  group homomorphism from $F(A)$ to $G$
  (\Cref{subsec:groups:universal-property}). What this really means
  is that for all sets $A$ and all groups $G$ there are natural
  identifications
  \[
    \Hom_{\categ{Set}}(A, S(G)) \cong \Hom_{\categ{Grp}}(F(A), G),
  \]
  where $S(G)$ `forgets' the group structure of $G$. That is, the
  functor $F : \categ{Set} \to \categ{Grp}$ constructing free groups
  is left-adjoint to the forgetful functor
  $S : \categ{Grp} \to \categ{Set}$. This of course applies to every
  other construction of `free' objects we have encountered: the free
  functor is, as a rule, left-adjoint to the forgetful functor.
\end{exam}
Thus, interesting functors may turn out to be adjoints of
harmless-looking functors. This has technical advantages: properties
of the interesting ones may be translated into properties of the
harmless ones, thereby giving easier proofs.

In fact, the very fact that a functor has an adjoint will endow that
functor with convenient features. We say that $F$ is a
\emph{left-adjoint functor} if it has a right-adjoint, and that $G$ is
a \emph{right-adjoint functor} if it has a left-adjoint. Some
properties of (say) right-adjoint functors may be established without
even knowing what the companion left-adjoint functor may be. Here is
the prototypical example of this phenomenon.
\begin{lem}{}{lem:linrep:right-adjoints-commute-with-limits}
  Right-adjoint functors commute with limits. That is, if
  $G : \categ{D} \to \categ{C}$ has a left-adjoint
  $F : \categ{C} \to \categ{D}$ and
  $\mathscr{A} : \categ{I} \to \categ{D}$ is another functor, then
  there is a canonical isomorphism
  \[
    G\!\left(\varprojlim \mathscr{A}\right) \xrightarrow{\;\sim\;}
    \varprojlim(G \circ \mathscr{A})
  \]
  (if the limits exist, of course).
\end{lem}
As every good calculus student should readily understand, this says
that right-adjoint functors are \emph{continuous}. (Needless to say,
left-adjoint functors turn out to commute with colimits and would
rightly be termed \emph{cocontinuous}.)  We will not prove
\Cref{lem:linrep:right-adjoints-commute-with-limits} in gory detail, but I will endeavor to convince the
reader that such statements are easier than they look and just boil
down to suitable applications of the universal properties defining the
various concepts. By contrast, proving that a specific given functor
preserves products (for instance), without appealing to abstract
nonsense, may sometimes appear to involve some `real' work.

Assume $G : \categ{D} \to \categ{C}$ is right-adjoint to
$F : \categ{C} \to \categ{D}$ and
$\mathscr{A} : \categ{I} \to \categ{D}$ is a given functor. As we have
seen, the limit of $\mathscr{A}$ is final subject to fitting in
commutative diagrams (with hopefully evident notation). Applying $G$,
we get commutative diagrams in $\categ{C}$, and hence, by the
universal property defining the limit of $G \circ \mathscr{A}$, a
morphism
$G(\varprojlim \mathscr{A}) \to \varprojlim(G \circ \mathscr{A})$.
Now, the morphisms in $\categ{C}$,
$\varprojlim(G \circ \mathscr{A}) \to G \circ \mathscr{A}(I)$,
determine via the adjunction identification
$\Hom_{\categ{C}}(X, G(Y)) \xrightarrow{\sim} \Hom_{\categ{D}}(F(X),
Y)$ morphisms in $\categ{D}$:
$F(\varprojlim(G \circ \mathscr{A})) \to \mathscr{A}(I)$.  By the
universal property defining $\varprojlim \mathscr{A}$ we get a
morphism
$F(\varprojlim(G \circ \mathscr{A})) \to \varprojlim \mathscr{A}$, and
applying adjunction again determines a morphism
$\varprojlim(G \circ \mathscr{A}) \to G(\varprojlim \mathscr{A})$.
Summarizing, we have obtained natural morphisms
\[
  G\!\left(\varprojlim \mathscr{A}\right) \to \varprojlim(G \circ
  \mathscr{A}), \quad \varprojlim(G \circ \mathscr{A}) \to
  G\!\left(\varprojlim \mathscr{A}\right).
\]
The compositions of these are easily checked to be the identity (by
virtue of the uniqueness part of various universality properties
invoked in the construction), concluding the verification of
\Cref{lem:linrep:right-adjoints-commute-with-limits}. What is missing from this outline is the explicit
verification of the fact that every needed diagram commutes, which is
necessary in order to apply the universal properties as stated. The
reader will likely agree that such verifications, while possibly
somewhat involved, must be routine.

\Cref{lem:linrep:right-adjoints-commute-with-limits} implies, for example, that right-adjoint functors
preserve products and that they must `preserve kernels' when kernels
make sense.
\begin{exam}{(Exact functors)}{exam:linrep:exact-functor-definition-example}
  A functor is \emph{exact} if it preserves exactness, that is, it
  sends exact sequences to exact sequences. So far we have only
  studied exactness in the context of modules over a ring
  (\Cref{subsec:rings:complexes-and-exact-sequences}), so let $R$,
  $S$ be rings and let
  $F : R\text{-}\categ{Mod} \to S\text{-}\categ{Mod}$ be an additive
  functor; $F$ is exact if and only if whenever
  \[
    0 \to A \xrightarrow{\phi} B \xrightarrow{\psi} C \to 0
  \]
  is an exact sequence of $R$-modules, then
  \[
    0 \to F(A) \xrightarrow{F(\phi)} F(B) \xrightarrow{F(\psi)} F(C)
    \to 0
  \]
  is an exact sequence of $S$-modules. It follows easily that the
  image of every exact complex is exact
  (\Cref{exer:linrep:exact-functor-characterization}).

  A common and very interesting situation occurs when a functor is not
  exact but preserves `some' of the exactness of sequences (we will
  encounter such examples later in this chapter, for example in
  \Cref{subsec:linrep:exactness-properties-of-tensor-flatness}). We
  say that an additive functor $F$ is \emph{left-exact} if whenever
  \[
    0 \to A \xrightarrow{\phi} B \xrightarrow{\psi} C
  \]
  is exact, then so is
  \[
    0 \to F(A) \xrightarrow{F(\phi)} F(B) \xrightarrow{F(\psi)} F(C).
  \]
  It turns out that a left-exact functor gives rise to a whole
  sequence of `new', related (`derived') functors; this hugely
  important phenomenon is studied in homological algebra, and we will
  come back to it in \Cref{chap:homol}.
\end{exam}
\begin{prop}{}{prop:linrep:right-adjoint-additive-functors-left-exact}
  Right-adjoint additive functors
  $R\text{-}\categ{Mod} \to S\text{-}\categ{Mod}$ are left-exact.
\end{prop}
As the reader will verify
(\Cref{exer:linrep:right-adjoint-left-exact-left-adjoint-right-exact}),
this follows from the fact that right-adjoint functors preserve
kernels. (Note that exactness of $0 \to A \xrightarrow{\phi} B$
amounts to the fact that $\phi : A \hookrightarrow B$ identifies $A$
with $\ker \psi$.) Of course the corresponding co-statements hold:
since left-adjoint functors preserve colimits and cokernels are
colimits, it follows that left-adjoint additive functors
$G : R\text{-}\categ{Mod} \to S\text{-}\categ{Mod}$ are necessarily
\emph{right-exact}, in the sense that if
$A \xrightarrow{\phi} B \xrightarrow{\psi} C \to 0$ is exact, then
$G(A) \xrightarrow{G(\phi)} G(B) \xrightarrow{G(\psi)} G(C) \to 0$ is
also exact.

\subsection*{Exercises}

\begin{exer}{}{exer:linrep:functor-preserves-product-map}
  Let $F : \categ{C} \to \categ{D}$ be a covariant functor, and assume
  that both $\categ{C}$ and $\categ{D}$ have products. Prove that for
  all objects $A$, $B$ of $\categ{C}$, there is a unique morphism
  $F(A \times B) \to F(A) \times F(B)$ such that the relevant diagram
  involving natural projections commutes. If $\categ{D}$ has
  coproducts (denoted $\sqcup$) and $G : \categ{C} \to \categ{D}$ is
  contravariant, prove that there is a unique morphism
  $G(A) \sqcup G(B) \to G(A \times B)$ (again, such that an
  appropriate diagram commutes).
\end{exer}

\begin{exer}{$\lnot$}{exer:linrep:fully-faithful-reflects-iso}
  Let $F : \categ{C} \to \categ{D}$ be a fully faithful functor. If
  $A$, $B$ are objects in $\categ{C}$, prove that $A \cong B$ in
  $\categ{C}$ if and only if $F(A) \cong F(B)$ in $\categ{D}$.
  [\Cref{subsec:linrep:equivalence-of-categories}]
\end{exer}

\begin{exer}{}{exer:linrep:grp-action-is-functor-from-groupoid}
  Recall (\Cref{sec:groups:definition-of-group}) that a group $G$
  may be thought of as a groupoid $\mathcal{G}$ with a single
  object. Prove that defining the action of $G$ on an object of a
  category $\categ{C}$ is equivalent to defining a functor
  $\mathcal{G} \to \categ{C}$.
\end{exer}

\begin{exer}{$\lnot$}{exer:linrep:localization-is-a-functor}
  Let $R$ be a commutative ring, and let $S \subseteq R$ be a
  multiplicative subset in the sense of
  \Cref{exer:irred:localization-construction}. Prove that
  `localization is a functor': associating with every $R$-module $M$
  the localization $S^{-1}M$
  (\Cref{exer:irred:localization-of-module}) and with every $R$-module
  homomorphism $\phi : M \to N$ the naturally induced homomorphism
  $S^{-1}M \to S^{-1}N$ defines a covariant functor from the category
  of $R$-modules to the category of $S^{-1}R$-modules.
  [\ref{exer:linrep:localization-is-exact-preserves-homology}]
\end{exer}

\begin{exer}{$*$}{exer:linrep:field-units-functor-and-norm-contravariant}
  For $F$ a field, denote by $F^*$ the group of nonzero elements of
  $F$, with multiplication. The assignment
  $\categ{Fld} \to \categ{Grp}$ mapping $F$ to $F^*$ and a
  homomorphism of fields $\phi : k \to F$ (i.e., a field
  extension\footnote{Recall that there are no `zero-homomorphisms'
    between fields; cf.\ \Cref{subsec:fields:basic-definitions}.})
  to the restriction $\phi|_{k^*} : k^* \to F^*$ is clearly a
  covariant functor. On the other hand, we can consider the category
  $\categ{Fld}^f_k$ of finite extensions of a fixed field $k$. Prove
  that the assignment $F \mapsto F^*$ on objects, together with the
  prescription associating with every $F_1 \subseteq F_2$ the norm
  $N_{F_1 \subseteq F_2} : F_2^* \to F_1^*$ (cf.\
  \Cref{exer:fields:norm-is-multiplicative}), gives a contravariant
  functor $\categ{Fld}^f_k \to \categ{Grp}$.  State and prove an
  analogous statement for the trace (cf.\
  \Cref{exer:fields:trace-is-additive}).
\end{exer}

\begin{exer}{$\lnot$}{exer:linrep:presheaf-continuous-fns-gluing}
  Formalize the notion of presheaf of abelian groups on a topological
  space $T$. If $F$ is a presheaf on $T$, elements of $F(U)$ are
  called \emph{sections} of $F$ on $U$. The homomorphism
  $\rho_{UV} : F(U) \to F(V)$ induced by an inclusion $V \subseteq U$
  is called the \emph{restriction map}. Show that an example of a
  presheaf is obtained by letting $\mathcal{C}(U)$ be the additive
  abelian group of continuous complex-valued functions on $U$, with
  restriction of sections defined by ordinary restriction of
  functions.  For this presheaf, prove that one can uniquely glue
  sections agreeing on overlapping open sets. That is, if $U$ and $V$
  are open sets and $s_U \in \mathcal{C}(U)$, $s_V \in \mathcal{C}(V)$
  agree after restriction to $U \cap V$, prove that there exists a
  unique $s \in \mathcal{C}(U \cup V)$ such that $s$ restricts to
  $s_U$ on $U$ and to $s_V$ on $V$. This is essentially the condition
  making $\mathcal{C}$ a
  sheaf. [\ref{exer:homol:exp-map-not-epi-of-presheaves}]
\end{exer}

\begin{exer}{$\lnot$}{exer:linrep:spec-zariski-topology-is-functor}
  Define a topology on $\mathrm{Spec}\,R$ by declaring the closed sets
  to be the sets $V(I)$, where $I \subseteq R$ is an ideal and $V(I)$
  denotes the set of prime ideals containing $I$.
  \begin{itemize}
  \item Verify that this indeed defines a topology on
    $\mathrm{Spec}\,R$. (This is the \emph{Zariski topology} on
    $\mathrm{Spec}\,R$.)
  \item Relate this topology to the Zariski topology defined in
    \Cref{subsec:fields:a-little-affine-algebraic-geometry}.
  \item Prove that $\mathrm{Spec}$ is then a contravariant functor
    from the category of commutative rings to the category of
    topological spaces (where morphisms are continuous functions).
  \end{itemize} [\Cref{subsec:linrep:examples-of-functors}]
\end{exer}

\begin{exer}{}{exer:linrep:hom-functor-at-point-and-poly-fns}
  Let $K$ be an algebraically closed field, and consider the category
  $K\text{-}\categ{Aff}$ defined in \Cref{exam:linrep:category-of-affine-algebraic-sets}.
  \begin{itemize}
  \item Denote by $h_S$ the functor
    $\Hom_{K\text{-}\mathrm{Aff}}(\mathord{-}, S)$ (as in
    \Cref{subsec:linrep:examples-of-functors}), and let
    $p = \A^0_K$ be a point. Show that there is a natural bijection
    between $S$ and $h_S(p)$. (Use
    \Cref{exer:fields:points-correspond-to-maximal-ideals}.)
  \item Show how every $\phi \in \Hom_{K\text{-}\mathrm{Aff}} (S, T)$
    determines a function of sets $S \to T$.
  \item If $S \subseteq \A^m_K$, $T \subseteq \A^n_K$, show that the
    function $S \to T$ determined by a morphism
    $\phi \in \Hom_{K\text{-}\mathrm{Aff}}(S, T)$ is the restriction
    of a `polynomial function' $\A^m_K \to \A^n_K$. (Part of this
    exercise is to make sense of what this means!)
  \end{itemize}
\end{exer}

\begin{exer}{$\lnot$}{exer:linrep:yoneda-embedding-is-functor}
  Let $\categ{C}$, $\categ{D}$ be categories, and assume $\categ{C}$
  to be small. Define a \emph{functor category}
  $\categ{D}^{\categ{C}}$, whose objects are covariant functors
  $\categ{C} \to \categ{D}$ and whose morphisms are natural
  transformations\footnote{The reason why $\categ{C}$ is assumed to be
    small is that this ensures that the natural transformations
    between two functors do form a set.}. Prove that the assignment
  $X \mapsto h_X := \Hom_{\categ{C}} (\mathord{-}, X)$ (cf.\
  \Cref{subsec:linrep:examples-of-functors}) defines a covariant
  functor $\categ{C} \to
  \categ{Set}^{\categ{C}^{\mathrm{op}}}$. (Define the action on
  morphisms in the natural way.)
  [\ref{exer:linrep:small-category-embeds-in-functor-category},
  \ref{exer:homol:functor-category-into-abelian-is-abelian}]
\end{exer}

\begin{exer}{$\lnot$}{exer:linrep:yoneda-lemma}
  Let $\categ{C}$ be a category, $X$ an object of $\categ{C}$, and
  consider the contravariant functor
  $h^X := \Hom_{\categ{C}} (\mathord{-}, X)$ (cf.\
  \Cref{subsec:linrep:examples-of-functors}). For every
  contravariant functor $F : \categ{C} \to \categ{Set}$, prove that
  there is a bijection between the set of natural transformations
  $h^X \leadsto F$ and $F(X)$, defined as follows. The datum of a
  natural transformation $\nu : h^X \leadsto F$ consists of a morphism
  $\nu_A$ from $h^X(A) = \Hom_{\categ{C}}(A, X)$ to $F(A)$ for every
  object $A$ of $\categ{C}$. Map $\nu$ to the image
  $\nu_X(\id_X)$ of $\id_X \in h^X(X)$ in
  $F(X)$. (Hint: Produce an inverse of the specified map. For every
  $f \in F(X)$ and every $\phi \in \Hom_{\categ{C}}(A, X)$, how do you
  construct an element of $F(A)$?) This result is called the
  \emph{Yoneda lemma}.
  [\ref{exer:linrep:small-category-embeds-in-functor-category},
  \ref{exer:homol:yoneda-embedding-left-exact-reflects-exactness}]
\end{exer}

\begin{exer}{}{exer:linrep:small-category-embeds-in-functor-category}
  (Cf.\ \Cref{exer:linrep:yoneda-embedding-is-functor}.) Let
  $\categ{C}$ be a small category. A contravariant functor
  $\categ{C} \to \categ{Set}$ is \emph{representable} if it is
  naturally isomorphic to a functor $h^X$ (cf.\
  \Cref{subsec:linrep:examples-of-functors}). In this case, $X$
  `represents' the functor. Prove that $\categ{C}$ is equivalent to
  the subcategory of representable functors in
  $\categ{Set}^{\categ{C}^{\mathrm{op}}}$. (Hint: Yoneda; see
  \Cref{exer:linrep:yoneda-lemma}.) Thus, every (small) category is
  equivalent to a subcategory of a functor category.
\end{exer}

\begin{exer}{}{exer:linrep:left-adjoint-iff-representing-functor}
  Let $\categ{C}$, $\categ{D}$ be categories, and let
  $F : \categ{C} \to \categ{D}$, $G : \categ{D} \to \categ{C}$ be
  functors. Prove that $F$ is left-adjoint to $G$ if and only if, for
  every object $Y$ in $\categ{D}$, the object $G(Y)$ represents the
  functor $h^Y \circ F$ (`naturally' in $Y$).
\end{exer}

\begin{exer}{}{exer:linrep:limit-over-empty-category}
  Let $\categ{Z}$ be the `Zen' category consisting of no objects and
  no morphisms. One can contemplate a functor $Z$ from $\categ{Z}$ to
  any category $\categ{C}$: no datum whatsoever need be
  specified. What is $\varprojlim Z$ (when such an object exists)?
\end{exer}

\begin{exer}{$\lnot$}{exer:linrep:construction-recovers-kernel}
  Verify that the construction described in \Cref{exam:linrep:equalizer-as-limit}
  indeed recovers the kernel of a homomorphism of $R$-modules, as
  claimed. [\Cref{subsec:linrep:limits-and-colimits}]
\end{exer}

\begin{exer}{$\lnot$}{exer:linrep:construction-is-inverse-limit}
  Verify that the construction given in the proof of
  \Cref{prop:linrep:inverse-limit-of-modules-exists} is an inverse limit, as
  claimed. [\Cref{subsec:linrep:limits-and-colimits}]
\end{exer}

\begin{exer}{$\lnot$}{exer:linrep:colimits-in-set-and-mod-directed-poset}
  Flesh out the sketch of the constructions of colimits in
  $\categ{Set}$ and $R\text{-}\categ{Mod}$ given in
  \Cref{subsec:linrep:limits-and-colimits}, for any indexing poset.
  In $\categ{Set}$, observe that the construction of the colimit is
  simpler if the poset $\categ{I}$ is directed; that is, if
  $\forall i, j \in \categ{I}$, there exists $k \in \categ{I}$ such
  that $i \leq k$, $j \leq
  k$. [\Cref{subsec:linrep:limits-and-colimits}]
\end{exer}

\begin{exer}{$\lnot$}{exer:linrep:i-adic-completion-exists}
  Let $R$ be a commutative ring, and let $I \subseteq R$ be an ideal.
  Note that $I^n \subseteq I^m$ if $n \geq m$, and hence we have
  natural homomorphisms $\phi_{mn} : R/I^n \to R/I^m$ for $n \geq m$.
  \begin{itemize}
  \item Prove that the inverse limit
    $\hat{R}_I := \varprojlim_n R/I^n$ exists as a commutative
    ring. This is called the \emph{$I$-adic completion} of $R$.
  \item By the universal property of inverse limits, there is a unique
    homomorphism $R \to \hat{R}_I$. Prove that the kernel of this
    homomorphism is $\bigcap_n I^n$.
  \item Let $I = (x)$ in $R[x]$. Prove that the completion
    $\widehat{R[x]}_I$ is isomorphic to the power series ring $R[[x]]$
    defined in \Cref{subsec:rings:polynomial-rings}.
  \end{itemize} [\ref{exer:linrep:intersection-of-ideal-powers-krull},
  \ref{exer:linrep:p-adic-integers-basics}]
\end{exer}

\begin{exer}{}{exer:linrep:intersection-of-ideal-powers-krull}
  Let $R$ be a commutative Noetherian ring, and let $I \subseteq R$ be
  an ideal. Then $I \cdot \bigcap_n I^n = \bigcap_n I^n$; the reader
  will prove this in
  \Cref{exer:linrep:intersection-of-ideal-powers-via-artin-rees}. Assume
  this, and prove that $\bigcap_n I^n$ equals the set of $r \in R$
  such that $(1-a)r = 0$ for some $a \in I$. (Hint: One inclusion is
  elementary. For the other, use the Nakayama lemma in the form of
  \Cref{exer:linalg:jm-eq-m-implies-1-plus-b-kills-m}. This result is
  attributed to Krull.) For example, if $I$ is proper, then
  $\bigcap_n I^n = (0)$ if $R$ is an integral domain or if it is
  local.  In these cases, the natural map $R \to \hat{R}_I$ to the
  $I$-adic completion is injective (cf.\
  \Cref{exer:linrep:i-adic-completion-exists}).
\end{exer}

\begin{exer}{$\lnot$}{exer:linrep:p-adic-integers-basics}
  An important example of the construction presented in
  \Cref{exer:linrep:i-adic-completion-exists} is the ring $\Z_p$ of
  \emph{$p$-adic integers}: this is the limit
  $\varprojlim_r \Z/p^r\Z$, for a positive prime integer $p$. The
  field of fractions of $\Z_p$ is denoted $\Q_p$; elements of $\Q_p$
  are called \emph{$p$-adic numbers}.
  \begin{itemize}
  \item Show that giving a $p$-adic integer $A$ is equivalent to
    giving a sequence of integers $A_r$, $r \geq 1$, such that
    $0 \leq A_r < p^r$, and that $A_s \equiv A_r \pmod{p^s}$ if
    $s \leq r$.
  \item Equivalently, show that every $p$-adic integer has a unique
    infinite expansion
    $A = a_0 + a_1 \cdot p + a_2 \cdot p^2 + a_3 \cdot p^3 + \dots$,
    where $0 \leq a_i \leq p - 1$. The arithmetic of $p$-adic integers
    may be carried out with these expansions in precisely the same way
    as ordinary arithmetic is carried out with ordinary decimal
    expansions.
  \item With notation as in the previous point, prove that
    $A \in \Z_p$ is invertible if and only if $a_0 \neq 0$.
  \item Prove that $\Z_p$ is a local domain, with maximal ideal
    generated by (the image in $\Z_p$ of) $p$.
  \item Prove that $\Z_p$ is a DVR (cf.\
    \Cref{exer:irred:discrete-valuation-ring-basics}).  (There is an
    evident valuation on $\Q_p$.)
  \end{itemize}
  [\Cref{subsec:groups:cyclic-groups-and-modular-arithmetic}]
\end{exer}

\begin{exer}{$\lnot$}{exer:linrep:zhat-iso-to-end-of-q-mod-z}
  If $m$, $n$ are positive integers and $m \mid n$, then
  $(n) \subseteq (m)$, and there is an onto ring homomorphism
  $\Z/n\Z \to \Z/m\Z$. The limit ring
  $\hat{\Z} := \varprojlim_{(n)} \Z/n\Z$ exists and is denoted by
  $\hat{\Z}$. Prove that
  $\hat{\Z} \cong \mathrm{End}_{\categ{Ab}}(\Q/\Z)$.  (Every
  $f \in \mathrm{End}_{\categ{Ab}}(\Q/\Z)$ is determined by
  $f(\tfrac{1}{n})$; note that since
  $n \cdot \tfrac{1}{n} = 1 \equiv 0 \pmod{\Z}$,
  $f(\tfrac{1}{n}) = \tfrac{g(n)}{n}$ for some integer $g(n)$, which
  may be chosen so that $0 \leq g(n) < n$.  Show that
  $g(m) \equiv g(n) \pmod{m}$ if $m \mid n$, and think about how
  elements of $\hat{\Z}$ may be described.)
  [\ref{exer:linrep:zhat-universal-property-product-of-zp}]
\end{exer}

\begin{exer}{}{exer:linrep:zhat-universal-property-product-of-zp}
  Let $\hat{\Z}$ be as in
  \Cref{exer:linrep:zhat-iso-to-end-of-q-mod-z}.
  \begin{itemize}
  \item If $R$ is a commutative ring endowed with homomorphisms
    $R \to \Z/p^r\Z$ for all primes $p$ and all $r$, compatible with
    all projections $\Z/p^r\Z \to \Z/p^s\Z$ for $s \leq r$, prove that
    there are ring homomorphisms $R \to \Z/n\Z$ for all $n$,
    compatible with all projections $\Z/n\Z \to \Z/m\Z$ for
    $m \mid n$.
  \item Deduce that $\hat{\Z}$ satisfies the universal property for
    the product of $\Z_p$, as $p$ ranges over all positive prime
    integers.
  \end{itemize}
  It follows that
  $\prod_p \Z_p \cong \hat{\Z} \cong
  \mathrm{End}_{\categ{Ab}}(\Q/\Z)$.
\end{exer}

\begin{exer}{}{exer:linrep:product-functor-right-adjoint-to-diagonal}
  Let $\categ{C}$ be a category, and consider the `product category'
  $\categ{C} \times \categ{C}$ (make sense of such a notion!). There
  is a `diagonal' functor associating to each object $X$ of
  $\categ{C}$ the pair $(X, X)$ as an object of
  $\categ{C} \times \categ{C}$. On the other hand, there may be a
  `product functor' $\categ{C} \times \categ{C} \to \categ{C}$,
  associating to $(X, Y)$ a product $X \times Y$; for example, this is
  the case in $\categ{Grp}$. Convince yourself that the product
  functor is right-adjoint to the diagonal functor. If there is a
  coproduct functor, verify that it is left-adjoint to the diagonal
  functor.
\end{exer}

\begin{exer}{$\lnot$}{exer:linrep:exact-functor-characterization}
  Let $R$, $S$ be rings. Prove that an additive covariant functor
  $F : R\text{-}\categ{Mod} \to S\text{-}\categ{Mod}$ is exact if and
  only if $F(A) \xrightarrow{F(\phi)} F(B) \xrightarrow{F(\psi)} F(C)$
  is exact in $S\text{-}\categ{Mod}$ whenever
  $A \xrightarrow{\phi} B \xrightarrow{\psi} C$ is exact in
  $R\text{-}\categ{Mod}$. Deduce that an exact functor sends exact
  complexes to exact
  complexes. [\Cref{subsec:linrep:comparing-functors},
  \ref{exer:homol:exact-functors-commute-with-cohomology}]
\end{exer}

\begin{exer}{}{exer:linrep:faithfully-exact-iff-kills-no-nonzero-module}
  Let $R$, $S$ be rings. An additive covariant functor
  $F : R\text{-}\categ{Mod} \to S\text{-}\categ{Mod}$ is
  \emph{faithfully exact} if
  `$F(A) \xrightarrow{F(\phi)} F(B) \xrightarrow{F(\psi)} F(C)$ is
  exact in $S\text{-}\categ{Mod}$ if and only if
  $A \xrightarrow{\phi} B \xrightarrow{\psi} C$ is exact in
  $R\text{-}\categ{Mod}$'. Prove that an exact functor
  $F : R\text{-}\categ{Mod} \to S\text{-}\categ{Mod}$ is faithfully
  exact if and only if $F(M) = 0$ for every nonzero $R$-module $M$, if
  and only if $F(\phi) = 0$ for every nonzero morphism $\phi$ in
  $R\text{-}\categ{Mod}$.
\end{exer}

\begin{exer}{$\lnot$}{exer:linrep:localization-is-exact-preserves-homology}
  Prove that localization
  (\Cref{exer:linrep:localization-is-a-functor}) is an exact
  functor. In fact, prove that localization `preserves homology': if
  \[
    M_\bullet : \cdots \to M_{i+1} \xrightarrow{d_{i+1}} M_i
    \xrightarrow{d_i} M_{i-1} \to \cdots
  \]
  is a complex of $R$-modules and $S$ is a multiplicative subset of
  $R$, then the localization $S^{-1}H_i(M_\bullet)$ of the $i$-th
  homology of $M_\bullet$ is the $i$-th homology
  $H_i(S^{-1}M_\bullet)$ of the localized complex
  \[
    S^{-1}M_\bullet : \cdots \to S^{-1}M_{i+1}
    \xrightarrow{S^{-1}d_{i+1}} S^{-1}M_i \xrightarrow{S^{-1}d_i}
    S^{-1}M_{i-1} \to \cdots.
  \] [\ref{exer:linrep:localization-is-flat},
  \ref{exer:linrep:flatness-is-local-property},
  \ref{exer:linrep:localization-commutes-with-tor}]
\end{exer}

\begin{exer}{}{exer:linrep:localization-is-faithfully-exact}
  Prove that localization is faithfully exact in the following sense:
  let $R$ be a commutative ring, and let
  \[
    (\ast)\quad 0 \to A \to B \to C \to 0
  \]
  be a sequence of $R$-modules. Then $(\ast)$ is exact if and only if
  the induced sequence of $R_{\mathfrak{p}}$-modules
  \[
    0 \to A_{\mathfrak{p}} \to B_{\mathfrak{p}} \to C_{\mathfrak{p}}
    \to 0
  \]
  is exact for every prime ideal $\mathfrak{p}$ of $R$, if and only if
  it is exact for every maximal ideal $\mathfrak{p}$.  (Cf.\
  \Cref{exer:irred:module-zero-iff-zero-locally}.)
\end{exer}

\begin{exer}{$\lnot$}{exer:linrep:right-adjoint-left-exact-left-adjoint-right-exact}
  Let $R$, $S$ be rings. Prove that right-adjoint functors
  $R\text{-}\categ{Mod} \to S\text{-}\categ{Mod}$ are left-exact and
  left-adjoint functors are
  right-exact. [\Cref{subsec:linrep:comparing-functors}]
\end{exer}

\begin{exer}{$\lnot$}{exer:linrep:center-of-mod-category-iso-to-center-of-ring}
  Let $\categ{C}$ be a category, and consider the identity functor
  $\mathrm{Id} : \categ{C} \to \categ{C}$. Prove that the set
  $\mathrm{End}(\mathrm{Id})$ of natural transformations
  $\mathrm{Id} \leadsto \mathrm{Id}$ is a commutative monoid under
  composition. This is called the \emph{center} of $\categ{C}$. If $R$
  is a ring, prove that the center of $R\text{-}\categ{Mod}$ is
  a ring in a natural way, and it is isomorphic to the center of
  $R$. [\ref{exer:linrep:morita-equivalence-and-center}]
\end{exer}

\section{Tensor products and the Tor functors}
\label{sec:linrep:tensor-products-and-the-tor-functors}

In the rest of the chapter we will work in the category
$R\text{-}\categ{Mod}$ of modules over a commutative ring
$R$. Essentially everything we will see can be upgraded to the
noncommutative case without difficulty, but a bit of structure is lost
in that case. For example, if $R$ is not commutative, then in the
category $R\text{-}\categ{Mod}$ of left-$R$-modules the Hom-sets
$\Hom_{R\text{-}\mathrm{Mod}} (M, N)$ are `only' abelian groups (cf.\
the end of \Cref{subsec:rings:the-category-r-mod}). A tensor
product $M \otimes_R N$ can only be defined if $M$ is a
right-$R$-module and $N$ is a left-$R$-module (in a sense, the two
module structures annihilate each other, and what is left is an
abelian group). By contrast, in the commutative case we will be able
to define $M \otimes_R N$ simply as an $R$-module. In general, the
theory goes through as in the commutative case if the modules carry
compatible left- and right-module structures, except in questions such
as the commutativity of tensors, where it would be unreasonable to
expect the commutativity of $R$ to have no bearing. All in all, the
commutative case is a little leaner, and (I believe) it suffices in
terms of conveying the basic intuition on the general features of the
theory.  Thus, $R$ will denote a fixed commutative ring, unless stated
otherwise.

\subsection{Bilinear maps and the definition of tensor product}
\label{subsec:linrep:bilinear-maps-and-the-definition-of-tensor-product}
If $M$ and $N$ are $R$-modules, we observed in the distant past
(\Cref{subsec:rings:products-and-coproducts}) that $M \oplus N$
serves as both the product and coproduct of $M$ and $N$: a situation
in which a limit coincides with a colimit. As a set, $M \oplus N$ is
just $M \times N$; the $R$-module structure on $M \oplus N$ is defined
by componentwise addition and multiplication by scalars. An $R$-module
homomorphism $M \oplus N \to P$ is determined by $R$-module
homomorphisms $M \to P$ and $N \to P$ (this is what makes $M \oplus N$
into a coproduct). But there is another way to map $M \times N$ to an
$R$-module $P$, compatibly with the $R$-module structures.
\begin{defn}{}{defn:linrep:bilinear-map}
  Let $M$, $N$, $P$ be $R$-modules. A function
  $\phi : M \times N \to P$ is \emph{$R$-bilinear} if
  \begin{itemize}
  \item $\forall m \in M$, the function $n \mapsto \phi(m, n)$ is an
    $R$-module homomorphism $N \to P$,
  \item $\forall n \in N$, the function $m \mapsto \phi(m, n)$ is an
    $R$-module homomorphism $M \to P$.
  \end{itemize}
\end{defn}
Thus, if $\phi : M \times N \to P$ is $R$-bilinear, then
$\forall m \in M$, $\forall n_1, n_2 \in N$, $\forall r_1, r_2 \in R$,
\[
  \phi(m, r_1 n_1 + r_2 n_2) = r_1 \phi(m, n_1) + r_2 \phi(m, n_2),
\]
and similarly for $\phi(\mathord{-}, n)$. Note that $\phi$ itself is
not linear, even if we view $M \times N$ as the $R$-module
$M \oplus N$, as recalled above. On the other hand, there ought to be
a way to deal with $R$-bilinear maps `as if' they were $R$-linear,
because such maps abound in the context of $R$-modules. For example,
the very multiplication on $R$ is itself an $R$-bilinear map
$R \times R \to R$.  Our experience with universal properties suggests
the natural way to approach this question. What we need is a new
$R$-module $M \otimes_R N$, with an $R$-bilinear map
$\otimes : M \times N \to M \otimes_R N$, such that every $R$-bilinear
map $M \times N \to P$ factors uniquely through this new module
$M \otimes_R N$, in such a way that the map $\phi$ is a usual
$R$-module homomorphism. Thus, $M \otimes_R N$ would be the `best
approximation' to $M \times N$ available in $R\text{-}\categ{Mod}$, if
we want to view $R$-bilinear maps from $M \times N$ as $R$-linear. The
module $M \otimes_R N$ is called the \emph{tensor product} of $M$ and
$N$ over $R$. The subscript $R$ is very important: if $M$ and $N$ are
modules over two rings $R$, $S$, then $S$-bilinearity is not the same
as $R$-bilinearity, so $M \otimes_R N$ and $M \otimes_S N$ may be
completely different objects. In context, it is not unusual to drop
the subscript if the base ring is understood, but I do not recommend
this practice. The prescription given above expresses the tensor
product as the solution to a universal problem; therefore we know
right away that it will be unique up to isomorphism, if it exists
(\Cref{prop:prelims:initial-final-unique} once more), and we could
proceed to study it by systematically using the universal property.
\begin{exam}{}{exam:linrep:tensor-with-base-ring-is-identity}
  For all $R$-modules $N$, $R \otimes_R N \cong N$. Indeed, every
  $R$-bilinear $R \times N \to P$ factors through $N$ (as is
  immediately verified) via $\otimes(r, n) = rn$. By the uniqueness
  property of universal objects, necessarily $N \cong R \otimes_R N$.
\end{exam}
For another example, it is easy to see that there must be a canonical
isomorphism $M \otimes_R N \xrightarrow{\sim} N \otimes_R M$. Indeed,
every $R$-bilinear $\phi : M \times N \to P$ may be
decomposed\footnote{Here is one situation in which the commutativity
  of $R$ does play a role: if $R$ is not commutative, then this
  decomposition becomes problematic, even if $M$ and $N$ carry
  bimodule structures. One can therefore not draw the conclusion
  $M \otimes_R N \cong N \otimes_R M$ in that case.} as
$M \times N \to N \times M \xrightarrow{\psi} P$, where
$\psi(n, m) = \phi(m, n)$; $\psi$ is also $R$-bilinear, so it factors
uniquely through $N \otimes_R M$. Therefore, $\phi$ factors uniquely
through $N \otimes_R M$, and this is enough to conclude that there is
a canonical isomorphism $N \otimes_R M \cong M \otimes_R N$.  However,
such considerations are a little moot unless we establish that
$M \otimes_R N$ exists to begin with. This requires a bit of work.
\begin{lem}{}{lem:linrep:tensor-products-exist}
  Tensor products exist in $R\text{-}\categ{Mod}$.
\end{lem}
\begin{proof}
  Given $R$-modules $M$ and $N$, we construct `by hand' a module
  satisfying the universal requirement. Let
  $F_R(M \times N) = R^{\oplus(M \times N)}$ be the free $R$-module on
  $M \times N$
  (\Cref{subsec:rings:free-modules-and-free-algebras}). This module
  comes equipped with a set-map $j : M \times N \to F_R(M \times N)$,
  universal with respect to all set-maps from $M \times N$ to any
  $R$-module $P$; the main task is to make this into an $R$-bilinear
  map. For example, we have to identify elements in $F_R(M \times N)$
  of the form $j(m, n_1 + n_2)$ with elements $j(m, n_1) + j(m, n_2)$,
  etc. Thus, let $K$ be the $R$-submodule of $F_R(M \times N)$
  generated by all elements
  \[
    j(m, r_1 n_1 + r_2 n_2) - r_1 j(m, n_1) - r_2 j(m, n_2)
  \]
  and
  \[
    j(r_1 m_1 + r_2 m_2, n) - r_1 j(m_1, n) - r_2 j(m_2, n)
  \]
  as $m, m_1, m_2$ range in $M$, $n, n_1, n_2$ range in $N$, and
  $r_1, r_2$ range in $R$. Let
  \[
    M \otimes_R N := \frac{F_R(M \times N)}{K},
  \]
  endowed with the map $\otimes : M \times N \to M \otimes_R N$
  obtained by composing $j$ with the natural projection. The element
  $\otimes(m, n)$ (that is, the class of $j(m, n)$ modulo $K$) is
  denoted $m \otimes n$. It is evident that
  $(m, n) \mapsto m \otimes n$ defines an $R$-bilinear map. We have to
  check that $M \otimes_R N$ satisfies the universal property, and
  this is also straightforward. If $\phi : M \times N \to P$ is any
  $R$-bilinear map, we have a unique induced $R$-linear map
  $\tilde{\phi}$ from the free $R$-module, by the universal property
  of the latter. I claim that $\tilde{\phi}$ restricts to $0$ on
  $K$. Indeed, to verify this, it suffices to verify that
  $\tilde{\phi}$ sends to zero every generator of $K$, and this
  follows from the fact that $\phi$ is $R$-bilinear. For example,
  \begin{align*}
    \tilde{\phi}(j(m, r_1 n_1 + r_2 n_2) - r_1 j(m,n_1) - r_2 j(m,n_2))
    &= \phi(m, r_1 n_1 + r_2 n_2) - r_1 \phi(m, n_1) - r_2 \phi(m, n_2)
    \\
    &= 0.
  \end{align*}
  It follows (by the universal property of quotients!) that
  $\tilde{\phi}$ factors uniquely through the quotient by $K$, and we
  are done.
\end{proof}
As is often the case with universal objects, the explicit construction
used to prove the existence of $M \otimes_R N$ is almost never
invoked.  It is however good to keep in mind that elements of
$M \otimes_R N$ arise from elements of the free $R$-module on
$M \times N$, and therefore an arbitrary element of $M \otimes_R N$ is
a finite sum
\[
  \tag{$*$} \sum_i m_i \otimes n_i
\]
with $m_i \in M$ and $n_i \in N$. The $R$-bilinearity of
$\otimes : M \times N \to M \otimes_R N$ amounts to the rules:
\begin{align*}
  m \otimes (n_1 + n_2) &= m \otimes n_1 + m \otimes n_2, \\
  (m_1 + m_2) \otimes n &= m_1 \otimes n + m_2 \otimes n, \\
  m \otimes (rn) &= (rm) \otimes n = r(m \otimes n),
\end{align*}
for all $m, m_1, m_2 \in M$, $n_1, n_2, n \in N$, and $r \in R$. In
particular, note that coefficients $r_i$ are not necessary in $(*)$,
since they can be absorbed:
$\sum_i r_i(m_i \otimes n_i) = \sum_i (r_i m_i) \otimes n_i$.

Elements of the form $m \otimes n$ (needing only one summand) are
called \emph{pure tensors}. Dear reader, please remember that pure
tensors are special: usually, not every element of the tensor product
is a pure tensor. See \Cref{exer:linrep:pure-tensor-if-n-cyclic} for
one situation in which every tensor happens to be pure, and appreciate
how special that is.

Pure tensors are nevertheless very useful, as a set of generators for
the tensor product. For example, if two homomorphisms
$\alpha, \beta : M \otimes_R N \to P$ coincide on pure tensors, then
$\alpha = \beta$. Frequently, computations involving tensor products
are reduced to simple verifications for pure tensors.

\subsection{Adjunction with Hom and explicit computations}
\label{subsec:linrep:adjunction-with-hom-and-explicit-computations}

The tensor product $\otimes_R$ is left-adjoint to Hom. Once we parse
what this rough statement means, it will be a near triviality; but as
we have found out in \Cref{subsec:linrep:comparing-functors}, the
mere fact that $\otimes_R$ is left-adjoint to any functor is enough to
draw interesting conclusions about it.

First, we note that every $R$-module $N$ defines, via $\otimes_R$, a
new covariant functor $R\text{-}\categ{Mod} \to R\text{-}\categ{Mod}$,
defined on objects by $M \mapsto M \otimes_R N$. To see how this works
on morphisms, let $\alpha : M_1 \to M_2$ be an $R$-module
homomorphism.  Crossing with $N$ and composing with $\otimes$ defines
an $R$-bilinear map $M_1 \times N \to M_2 \times N \to M_2 \otimes N$,
and hence an induced $R$-linear map
$\alpha \otimes N : M_1 \otimes N \to M_2 \otimes N$. On pure tensors,
this map is simply given by $m \otimes n \mapsto \alpha(m) \otimes n$,
and functoriality follows immediately: if $\beta : M_0 \to M_1$ is a
second homomorphism, then $(\alpha \otimes N) \circ (\beta \otimes N)$
and $(\alpha \circ \beta) \otimes N$ both map pure tensors
$m \otimes n$ to $\alpha(\beta(m)) \otimes n$, so they must agree on
all tensors.

The adjunction statement compares this functor with the covariant
functor $P \mapsto \Hom_R(N, P)$; cf.\
\Cref{subsec:linrep:examples-of-functors}. We have defined
$M \otimes_R N$ so that giving an $R$-linear map $M \otimes_R N \to P$
to an $R$-module $P$ is `the same as' giving an $R$-bilinear map
$M \times N \to P$. Now recall the definition of $R$-bilinear map:
$\phi : M \times N \to P$ is $R$-bilinear if both
$\phi(m, \mathord{-})$ and $\phi(\mathord{-}, n)$ are $R$-linear maps,
for all $m \in M$ and $n \in N$. The first part of this prescription
says that $\phi$ determines a function $M \to \Hom_R(N, P)$; the
second part says that this is an $R$-module homomorphism. Therefore,
an $R$-bilinear map is `the same as' an element of
$\Hom_R(M, \Hom_R(N, P))$. These simple considerations should be
enough to make the following seemingly complicated statement rather
natural:
\begin{lem}{}{lem:linrep:hom-tensor-adjunction}
  For all $R$-modules $M$, $N$, $P$, there is an isomorphism of
  $R$-modules
  \[
    \Hom_R(M, \Hom_R(N, P)) \cong \Hom_R(M \otimes_R N, P).
  \]
\end{lem}
\begin{proof}
  As noted before the statement, every
  $\alpha \in \Hom_R(M, \Hom_R(N, P))$ determines an $R$-bilinear map
  $\phi : M \times N \to P$ by $(m, n) \mapsto \alpha(m)(n)$. By the
  universal property, $\phi$ factors uniquely through an $R$-linear
  map $\tilde\phi : M \otimes_R N \to P$. Therefore $\alpha$
  determines a well-defined element
  $\tilde\phi \in \Hom_R(M \otimes_R N, P)$.  The reader will check
  (\Cref{exer:linrep:complete-proof-of-adjunction}) that the map
  $\alpha \mapsto \tilde\phi$ is $R$-linear and construct an inverse.
\end{proof}
\begin{coro}{}{coro:linrep:tensor-hom-adjunction}
  For every $R$-module $N$, the functor $\otimes_R N$ is left-adjoint
  to the functor $\Hom_R(N, \mathord{-})$.
\end{coro}
\begin{proof}
  The claim is that the isomorphism found in \Cref{lem:linrep:hom-tensor-adjunction} is
  natural in the sense hinted at, but not fully explained, in
  \Cref{subsec:linrep:comparing-functors}; the interested reader
  should have no problems checking this naturality.
\end{proof}
By \Cref{lem:linrep:right-adjoints-commute-with-limits} (or rather its co-version), we can conclude that
for each $R$-module $N$, the functor $\otimes_R N$ preserves colimits,
and so does $M \otimes_R \mathord{-}$, by the basic commutativity of
tensor products verified in
\Cref{subsec:linrep:bilinear-maps-and-the-definition-of-tensor-product}. In
particular, and this is good material for another Pavlovian reaction,
$M \otimes_R \mathord{-}$ and $\mathord{-} \otimes_R N$ are
right-exact functors (cf.\ \Cref{exam:linrep:exact-functor-definition-example}). These
observations have several consequences, which make `computations' with
tensor products more reasonable. Here is a sample:
\begin{coro}{}{coro:linrep:tensor-distributes-over-direct-sum}
  For all $R$-modules $M_1$, $M_2$, $N$,
  \[
    (M_1 \oplus M_2) \otimes_R N \cong (M_1 \otimes_R N) \oplus (M_2
    \otimes_R N).
  \]
\end{coro}

(Moreover, by commutativity,
$M \otimes_R (N_1 \oplus N_2) \cong (M \otimes_R N_1) \oplus (M
\otimes_R N_2)$ as well.) Indeed, coproducts are colimits. In fact,
$\otimes$ must then commute with arbitrary (possibly infinite) direct
sums:
\[
  \Bigl(\bigoplus_{\alpha \in A} M_\alpha\Bigr) \otimes_R N \cong
  \bigoplus_{\alpha \in A} (M_\alpha \otimes_R N).
\]
This computes all tensors for free $R$-modules:
\begin{coro}{}{coro:linrep:tensor-of-free-modules}
  For any two sets $A$, $B$:
  \[
    R^{\oplus A} \otimes_R R^{\oplus B} \cong R^{\oplus A \times B}.
  \]
\end{coro}

Indeed, `distributing' the direct sum identifies the left-hand side
with the direct sum $(R^{\oplus A})^{\oplus B}$, which is isomorphic
to the right-hand side
(\Cref{exer:rings:free-module-product-sets}). For finitely generated
free modules, this simply says that
$R^{\oplus m} \otimes R^{\oplus n} \cong R^{\oplus mn}$.

Note that if $e_1, \dots, e_m$ generate $M$ and $f_1, \dots, f_n$
generate $N$, then the pure tensors $e_i \otimes f_j$ must generate
$M \otimes_R N$. In the free case, if the $e_i$'s and $f_j$'s form
bases of $R^{\oplus m}$, $R^{\oplus n}$, resp., then the $mn$ elements
$e_i \otimes f_j$ must be a basis for
$R^{\oplus m} \otimes R^{\oplus n}$.  Indeed they generate it; hence
they must be linearly independent since this module is free of rank
$mn$. In particular, this is all that can happen if $R$ is a field $k$
and the modules are, therefore, just $k$-vector spaces
(\Cref{prop:linalg:lin-indep-extends-to-basis}). Tensor products are more interesting over more
general rings.
\begin{coro}{}{coro:linrep:tensor-with-quotient-by-ideal}
  For all $R$-modules $N$ and all ideals $I$ of $R$,
  \[
    \frac{R}{I} \otimes_R N \cong \frac{N}{IN}.
  \]
\end{coro}

Indeed, $\mathord{-} \otimes_R N$ is right-exact; thus, the exact
sequence $0 \to I \to R \to R/I \to 0$ induces an exact sequence
\[
  I \otimes_R N \to R \otimes_R N \to \frac{R}{I} \otimes_R N \to 0.
\]
The image of $I \otimes_R N$ in $R \otimes_R N \cong N$ is generated
by the images of pure tensors $a \otimes n$ with $a \in I$, $n \in N$;
this is $IN$. Thus the sequence identifies $N/(IN)$ with
$(R/I) \otimes_R N$, as needed.
\begin{coro}{}{coro:linrep:tensor-of-two-quotients}
  For all ideals $I$, $J$ of $R$,
  \[
    \frac{R}{I} \otimes_R \frac{R}{J} \cong \frac{R}{I + J}.
  \]
\end{coro}

This follows immediately from \Cref{coro:linrep:tensor-with-quotient-by-ideal} and the `third
isomorphism theorem', \Cref{prop:rings:third-iso-modules}. Indeed,
$I(R/J) = (I+J)/J$.

\begin{exam}{}{exam:linrep:tensor-of-cyclic-groups}
  $\Z/m\Z \otimes_{\Z} \Z/n\Z \cong \Z/\gcd(m,n)\Z$.  Indeed,
  $(m)+(n) = (\gcd(m,n))$ in $\Z$. For instance,
  $\Z/2\Z \otimes_{\Z} \Z/3\Z \cong 0$, a favorite on qualifying exams
  (cf.\ \Cref{exer:linrep:coprime-cyclic-tensor-is-zero-by-hand}).
\end{exam}

\Cref{coro:linrep:tensor-with-quotient-by-ideal} is a template example for a basic application of
$\otimes$: tensor products may be used to transfer constructions
involving $R$ (such as quotienting by an ideal $I$) to constructions
involving $R$-modules (such as quotienting by a corresponding
submodule). There are several instances of this operation; the reader
will take a look at localization in
\Cref{exer:linrep:localization-is-tensor-with-localized-ring}.

\subsection{Exactness properties of tensor; flatness}
\label{subsec:linrep:exactness-properties-of-tensor-flatness}
It is important to remember that the tensor product is not an exact
functor: left-exactness may very well fail. This can already be
observed in the sequence appearing in the discussion following
\Cref{coro:linrep:tensor-with-quotient-by-ideal}: for an ideal $I$ of $R$ and an $R$-module $N$, the
map $I \otimes_R N \to N$ induced by the inclusion $I \subseteq R$
after tensoring by $N$ may not be injective.

\begin{exam}{}{exam:linrep:tensor-of-mult-by-2-map-with-z2}
  Multiplication by $2$ gives an inclusion
  $\Z \xrightarrow{\cdot 2} \Z$ identifying the first copy of $\Z$
  with the ideal $(2)$ in the second copy. Tensoring by $\Z/2\Z$ over
  $\Z$ (and keeping in mind that $R \otimes_R N \cong N$), we get the
  homomorphism $\Z/2\Z \xrightarrow{\cdot 2} \Z/2\Z$, which sends both
  $[0]$ and $[1]$ to zero. This is the zero-morphism, and in
  particular it is not injective.
\end{exam}
On the other hand, if $N \cong R^{\oplus A}$ is free, then
$\mathord{-} \otimes_R N$ is exact. Indeed, every inclusion
$M_1 \subseteq M_2$ is mapped to
$M_1 \otimes_R R^{\oplus A} \to M_2 \otimes_R R^{\oplus A}$, which is
identified (via \Cref{coro:linrep:tensor-distributes-over-direct-sum}) with the inclusion
$M_1^{\oplus A} \subseteq M_2^{\oplus A}$.

\begin{exam}{}{exam:linrep:tensoring-exact-for-vector-spaces}
  Since vector spaces are free (\Cref{prop:linalg:lin-indep-extends-to-basis}), tensoring is exact
  in $k\text{-}\categ{Vect}$: if
  \[
    0 \to V_1 \to V_2 \to V_3 \to 0
  \]
  is an exact sequence of $k$-vector spaces and $W$ is a $k$-vector
  space, then the induced sequence
  \[
    0 \to V_1 \otimes_k W \to V_2 \otimes_k W \to V_3 \otimes_k W \to
    0
  \]
  is exact on both sides.
\end{exam}
The reader should now wonder whether it is useful to study a condition
on an $R$-module $N$, guaranteeing that the functor
$\mathord{-} \otimes_R N$ is left-exact as well as right-exact.

\begin{defn}{}{defn:linrep:flat-module}
  An $R$-module $N$ is \emph{flat} if the functor
  $\mathord{-} \otimes_R N$ is exact.
\end{defn}

In the exercises the reader will explore easy properties of this
notion and useful equivalent formulations in particular cases. We have
already checked that $\Z/2\Z$ is not a flat $\Z$-module, while free
modules are flat. Flat modules are hugely important: in algebraic
geometry, `flatness' is the condition expressing the fact that the
objects in a family vary `continuously', preserving certain key
invariants.
\begin{exam}{}{exam:linrep:coordinate-ring-of-projection-example}
  Consider the affine algebraic set $V(xy)$ in the plane $\A^2$ (over
  a fixed field $k$) and the `projection on the first coordinate'
  $V(xy) \to \A^1$, $(x,y) \mapsto x$. In terms of coordinate rings
  (cf.\ \Cref{subsec:fields:a-little-affine-algebraic-geometry}),
  this map corresponds to the homomorphism of $k$-algebras
  $k[x] \to k[x,y]/(xy)$ defined by mapping $x$ to the coset
  $x + (xy)$ (this will be completely clear to the reader who has
  worked out
  \Cref{exer:fields:poly-fns-iso-to-coordinate-ring}!). This
  homomorphism defines a $k[x]$-module structure on $k[x,y]/(xy)$, and
  we can wonder whether the latter is flat in the sense of
  \Cref{defn:linrep:flat-module}. From the geometric point of view, clearly
  something `not flat' is going on over the point $x = 0$, so we
  consider the inclusion of the ideal $(x)$ in $k[x]$:
  \[
    k[x] \xrightarrow{\cdot x} k[x].
  \]
  Tensoring by $k[x,y]/(xy)$, we obtain
  \[
    \frac{k[x,y]}{(xy)} \xrightarrow{\cdot x} \frac{k[x,y]}{(xy)},
  \]
  which is not injective, because it sends to zero the nonzero coset
  $y + (xy)$. Therefore $k[x,y]/(xy)$ is not flat as a $k[x]$-module.
  The term \emph{flat} was inspired precisely by such `geometric'
  examples.
\end{exam}

\subsection{The Tor functors}
\label{subsec:linrep:the-tor-functors}
The `failure of exactness' of the functor $\mathord{-} \otimes_R N$ is
measured by another functor
$R\text{-}\categ{Mod} \to R\text{-}\categ{Mod}$, called
$\Tor^R_1(\mathord{-}, N)$: if $N$ is flat (for example, if it is
free), then $\Tor^R_1(M, N) = 0$ for all modules $M$. In fact
(amazingly) if
\[
  0 \to A \to B \to C \to 0
\]
is an exact sequence of $R$-modules, one obtains a new exact sequence
after tensoring by any $N$:
\[
  \Tor^R_1(C, N) \to A \otimes_R N \to B \otimes_R N \to C \otimes_R N
  \to 0,
\]
so if $\Tor^R_1(C, N) = 0$, then the module on the left vanishes; thus
every short exact sequence ending in $C$ remains exact after tensoring
by $N$ in this case. In fact (astonishingly) for all $N$ one can
continue this sequence with more Tor-modules, obtaining a longer exact
complex:
\[
  \Tor^R_1(A, N) \to \Tor^R_1(B, N) \to \Tor^R_1(C, N) \to A \otimes_R
  N \to B \otimes_R N \to C \otimes_R N \to 0.
\]
This is not the end of the story: the complex may be continued even
further by invoking new functors $\Tor^R_2(\mathord{-}, N)$,
$\Tor^R_3(\mathord{-}, N)$, etc. These are the derived functors of
tensor. To `compute' these functors, one may apply the following
procedure: given an $R$-module $M$, find a free resolution
(\Cref{subsec:linalg:finitely-presented-modules-and-free-resolutions})
\[
  \cdots \to R^{\oplus S_2} \to R^{\oplus S_1} \to R^{\oplus S_0} \to
  M \to 0;
\]
throw $M$ away, and tensor the free part by $N$, obtaining a complex
$M_\bullet \otimes_R N$:
\[
  \cdots \to N^{\oplus S_2} \to N^{\oplus S_1} \to N^{\oplus S_0} \to
  0
\]
(recall again that tensor commutes with colimits, hence with direct
sums, therefore $R^{\oplus m} \otimes_R N \cong N^{\oplus m}$); then
take the homology of this complex (cf.\
\Cref{subsec:rings:homology-and-the-snake-lemma}). Astoundingly,
this will not depend (up to isomorphism) on the chosen free
resolution, so we can define
\[
  \Tor^R_i(M, N) := H_i(M_\bullet \otimes_R N).
\]
For example, according to this definition
$\Tor^R_0(M, N) \cong M \otimes_R N$
(\Cref{exer:linrep:tor0-iso-to-tensor-product}), and
$\Tor^R_i(M, N) = 0$ for all $i > 0$ and all $M$ if $N$ is flat
(because then tensoring by $N$ is an exact functor, so tensoring the
resolution of $M$ returns an exact sequence, thus with no
homology). In fact, this proves a remarkable property of the Tor
functors: if $\Tor^R_1(M, N) = 0$ for all $M$, then
$\Tor^R_i(M, N) = 0$ for all $i > 0$ for all modules $M$.  Indeed, $N$
is then flat.  At this point you may feel that something is a little
out of balance: why focus on the functor $\mathord{-} \otimes_R N$,
rather than $M \otimes_R \mathord{-}$?  Since $M \otimes_R N$ is
canonically isomorphic to $N \otimes_R M$ (in the commutative case;
cf.\ \Cref{exam:linrep:tensor-with-base-ring-is-identity}), we could expect the same to apply to
every $\Tor^R_i$: $\Tor^R_i(M, N)$ ought to be canonically isomorphic
to $\Tor^R_i(N, M)$ for all $i$. Equivalently, we should be able to
compute $\Tor^R_i(M, N)$ as the homology of $M \otimes_R N_\bullet$,
where $N_\bullet$ is a free resolution of $N$.  This is indeed the
case.

In due time (\Cref{sec:homol:derived-functors} and 8) we will prove
this and all the other wonderful facts I have stated in this
subsection. For now, I am asking the reader to believe that the Tor
functors can be defined as I have indicated, and the facts reviewed
here will suffice for simple computations (see for example
\Cref{exer:linrep:tor1-r-torsion-of-n,exer:linrep:tor-vanishes-above-degree-1-over-pid})
and applications.  In fact, we know enough
about finitely generated modules over PIDs to get a preliminary sense
of what is involved in proving such general facts. Recall that we have
been able to establish that every finitely generated module $M$ over a
PID $R$ has a free resolution of length 1:
\[
  0 \to R^{\oplus m_1} \to R^{\oplus m_0} \to M \to 0.
\]
This property characterizes PIDs (\Cref{prop:linalg:pid-characterization-via-presentations}). If
$0 \to A \to B \to C \to 0$ is an exact sequence of $R$-modules, it is
not hard to see that one can produce `compatible' resolutions, in the
sense that the rows of the following diagram will be exact as well as
the columns:
\[
  \begin{array}{ccccccc}
    0 & \to & R^{\oplus a_1} & \to & R^{\oplus b_1} & \to &
                                                            R^{\oplus c_1} \to 0 \\
      & & \downarrow\alpha & & \downarrow\beta & &
                                                   \downarrow\gamma \\
    0 & \to & R^{\oplus a_0} & \to & R^{\oplus b_0} & \to &
                                                            R^{\oplus c_0} \to 0 \\
      & & \downarrow & & \downarrow & & \downarrow \\
    0 & \to & A & \to & B & \to & C \to 0 \\
      & & \downarrow & & \downarrow & & \downarrow \\
      & & 0 & & 0 & & 0
  \end{array}
\]
(This will be proven in gory detail in
\Cref{sec:homol:derived-functors}.) Tensor the two `free' rows by
$N$; they remain exact (tensoring commutes with direct sums):
\[
  \begin{array}{ccccccc}
    0 & \to & N^{\oplus a_1} & \to & N^{\oplus b_1} & \to &
                                                            N^{\oplus c_1} \to 0 \\
      & & \downarrow{\alpha \otimes N} & & \downarrow{\beta \otimes N} &
                                                          & \downarrow{\gamma \otimes N} \\
    0 & \to & N^{\oplus a_0} & \to & N^{\oplus b_0} & \to &
                                                            N^{\oplus c_0} \to 0
  \end{array}
\]
Now the columns (preceded and followed by $0$) are precisely the
complexes $A_\bullet \otimes_R N$, $B_\bullet \otimes_R N$,
$C_\bullet \otimes_R N$ whose homology `computes' the Tor modules.
Applying the snake lemma (\Cref{lem:rings:snake-lemma}; cf.\
\Cref{rem:rings:snake-lemma-homology-restatement}) gives the exact
sequence
\[
  H_1(A_\bullet \otimes_R N) \to H_1(B_\bullet \otimes_R N) \to
  H_1(C_\bullet \otimes_R N) \xrightarrow{\delta} H_0(A_\bullet
  \otimes_R N) \to H_0(B_\bullet \otimes_R N) \to H_0(C_\bullet
  \otimes_R N) \to 0,
\]

which is precisely the sequence of Tor modules conjured up above,
\[
  \Tor^R_1(A, N) \to \Tor^R_1(B, N) \to \Tor^R_1(C, N)
  \xrightarrow{\delta} A \otimes_R N \to B \otimes_R N \to C \otimes_R
  N \to 0,
\]
with a $0$ on the left for good measure (due to the fact that
$\Tor^R_2$ vanishes if $R$ is a PID; cf.\
\Cref{exer:linrep:tor-vanishes-above-degree-1-over-pid}).

Note that $\Tor^k_i$ vanishes for $i > 0$ if $k$ is a field, as vector
spaces are flat, and $\Tor^R_i$ vanishes for $i > 1$ if $R$ is a PID
(\Cref{exer:linrep:tor-vanishes-above-degree-1-over-pid}). These facts
are not surprising, in view of the procedure described above for
computing Tor and of the considerations at the end of
\Cref{subsec:linalg:pids-and-resolutions}: a bound on the length of
free resolutions for modules over a ring $R$ will imply a bound on
nonzero Tor's. For particularly nice rings (such as the rings
corresponding to `smooth' points in algebraic geometry) this bound
agrees with the Krull dimension; but precise results of this sort are
beyond the scope of this book.

\subsection*{Exercises}

$R$ denotes a fixed commutative ring.

\begin{exer}{$?$}{exer:linrep:pure-tensor-if-n-cyclic}
  Let $M$, $N$ be $R$-modules, and assume that $N$ is cyclic. Prove
  that every element of $M \otimes_R N$ may be written as a pure
  tensor. [\Cref{subsec:linrep:bilinear-maps-and-the-definition-of-tensor-product}]
\end{exer}
\begin{exer}{$?$}{exer:linrep:coprime-cyclic-tensor-is-zero-by-hand}
  Prove `by hand' (that is, without appealing to the right-exactness
  of tensor) that $\Z/n\Z \otimes_{\Z} \Z/m\Z \cong 0$ if $m$, $n$ are
  relatively prime
  integers. [\Cref{subsec:linrep:adjunction-with-hom-and-explicit-computations}]
\end{exer}
\begin{exer}{}{exer:linrep:polynomial-ring-tensor-product}
  Prove that
  $R[x_1, \dots, x_n] \otimes_R R[y_1, \dots, y_m] \cong R[x_1, \dots,
  x_n, y_1, \dots, y_m]$.
\end{exer}
\begin{exer}{$\lnot$}{exer:linrep:tensor-of-algebras-is-coproduct}
  Let $S$, $T$ be commutative $R$-algebras. Verify the following:
  \begin{itemize}
  \item The tensor product $S \otimes_R T$ has an operation of
    multiplication, defined on pure tensors by
    $(s_1 \otimes t_1) \cdot (s_2 \otimes t_2) := s_1 s_2 \otimes t_1
    t_2$ and making it into a commutative $R$-algebra.
  \item With respect to this structure, there are $R$-algebra
    homomorphisms $i_S : S \to S \otimes T$, resp.,
    $i_T : T \to S \otimes T$, defined by $i_S(s) := s \otimes 1$,
    $i_T(t) := 1 \otimes t$.
  \item The $R$-algebra $S \otimes_R T$, with these two structure
    homomorphisms, is a coproduct of $S$ and $T$ in the category of
    commutative $R$-algebras: if $U$ is a commutative $R$-algebra and
    $f_S : S \to U$, $f_T : T \to U$ are $R$-algebra homomorphisms,
    then there exists a unique $R$-algebra homomorphism
    $f_S \otimes f_T$ making the obvious diagram commute.
  \end{itemize}
  In particular, if $S$ and $T$ are simply commutative rings, then
  $S \otimes_{\Z} T$ is a coproduct of $S$ and $T$ in the category of
  commutative rings. This settles an issue left open at the end of
  \Cref{subsec:rings:products}. [\ref{exer:linrep:f-tensor-f-separable-galois-criteria}]
\end{exer}
\begin{exer}{$?$}{exer:linrep:localization-is-tensor-with-localized-ring}
  (Cf.\
  \Cref{exer:irred:localization-construction,exer:irred:localization-of-module}.)
  Let $S$ be a multiplicative subset
  of $R$, and let $M$ be an $R$-module. Prove that
  $S^{-1}M \cong M \otimes_R S^{-1}R$ as $R$-modules. (Use the
  universal property of the tensor product.) Through this isomorphism,
  $M \otimes_R S^{-1}R$ inherits an $S^{-1}R$-module structure.
  [\Cref{subsec:linrep:adjunction-with-hom-and-explicit-computations},
  \ref{exer:linrep:dim-of-m-tensor-k-equals-rank},
  \ref{exer:linrep:localization-is-flat},
  \ref{exer:linrep:reprove-localization-tensor-via-associativity}]
\end{exer}
\begin{exer}{$\lnot$}{exer:linrep:tensor-and-localization-commute}
  (Cf.\
  \Cref{exer:irred:localization-construction,exer:irred:localization-of-module}.)
  Let $S$ be a multiplicative subset
  of $R$, and let $M$ be an $R$-module.
  \begin{itemize}
  \item Let $N$ be an $S^{-1}R$-module. Prove that
    $(S^{-1}M) \otimes_{S^{-1}R} N \cong M \otimes_R N$.
  \item Let $A$ be an $R$-module. Prove that
    $(S^{-1}A) \otimes_R M \cong S^{-1}(A \otimes_R M)$.
  \end{itemize}
  (Both can be done `by hand', by analyzing the construction in
  \Cref{lem:linrep:tensor-products-exist}. Both isomorphisms will be easy consequences of
  the associativity of tensor products; cf.\
  \Cref{exer:linrep:associativity-of-tensor-product}.)
  [\ref{exer:linrep:flatness-is-local-property},
  \ref{exer:linrep:reprove-localization-tensor-via-associativity}]
\end{exer}
\begin{exer}{}{exer:linrep:changing-base-ring-in-tensor}
  Changing the base ring in a tensor may or may not make a difference:
  \begin{itemize}
  \item Prove that $\Q \otimes_{\Z} \Q \cong \Q \otimes_{\Q} \Q$.
  \item Prove that $\C \otimes_{\R} \C \cong \C \otimes_{\C} \C$.
  \end{itemize}
\end{exer}
\begin{exer}{}{exer:linrep:dim-of-m-tensor-k-equals-rank}
  Let $R$ be an integral domain, with field of fractions $K$, and let
  $M$ be a finitely generated $R$-module. The tensor product
  $V := M \otimes_R K$ is a $K$-vector space
  (\Cref{exer:linrep:localization-is-tensor-with-localized-ring}). Prove
  that $\dim_K V$ equals the rank of $M$ as an $R$-module, in the
  sense of \Cref{defn:linalg:rank-of-fg-module}.
\end{exer}
\begin{exer}{}{exer:linrep:fin-gen-abel-grp-tensor-q-and-zpz}
  Let $G$ be a finitely generated abelian group of rank $r$. Prove
  that $G \otimes_{\Z} \Q \cong \Q^r$. Prove that for infinitely many
  primes $p$, $G \otimes_{\Z} (\Z/p\Z) \cong (\Z/p\Z)^r$.
\end{exer}
\begin{exer}{}{exer:linrep:f-tensor-f-separable-galois-criteria}
  Let $k \subseteq k(\alpha) = F$ be a finite simple field extension.
  Note that $F \otimes_k F$ has a natural ring structure; cf.\
  \Cref{exer:linrep:tensor-of-algebras-is-coproduct}.
  \begin{itemize}
  \item Prove that $\alpha$ is separable over $k$ if and only if
    $F \otimes_k F$ is reduced as a ring.
  \item Prove that $k \subseteq F$ is Galois if and only if
    $F \otimes_k F$ is isomorphic to $F^{[F:k]}$ as a ring.
  \end{itemize}
  (Use \Cref{coro:linrep:tensor-with-quotient-by-ideal} to `compute' the tensor. The CRT from
  \Cref{subsec:irred:chinese-remainder-theorem} will likely be
  helpful.)
\end{exer}
\begin{exer}{$?$}{exer:linrep:complete-proof-of-adjunction}
  Complete the proof of
  \Cref{lem:linrep:hom-tensor-adjunction}. [\Cref{subsec:linrep:adjunction-with-hom-and-explicit-computations}]
\end{exer}
\begin{exer}{}{exer:linrep:localization-is-flat}
  Let $S$ be a multiplicative subset of $R$ (cf.\
  \Cref{exer:irred:localization-construction}).  Prove that $S^{-1}R$
  is flat over $R$. (Hint:
  \Cref{exer:linrep:localization-is-tensor-with-localized-ring,exer:linrep:localization-is-exact-preserves-homology}.)
\end{exer}
\begin{exer}{}{exer:linrep:direct-sums-of-flat-are-flat}
  Prove that direct sums of flat modules are flat.
\end{exer}
\begin{exer}{$?$}{exer:linrep:tor0-iso-to-tensor-product}
  Prove that, according to the definition given in
  \Cref{subsec:linrep:the-tor-functors}, $\Tor^R_0(M, N)$ is
  isomorphic to $M \otimes_R N$.
  [\Cref{subsec:linrep:the-tor-functors}]
\end{exer}
\begin{exer}{$?$}{exer:linrep:tor1-r-torsion-of-n}
  Prove that for $r \in R$ a non-zero-divisor and $N$ an $R$-module,
  the module $\Tor^R_1(R/(r), N)$ is isomorphic to the $r$-torsion of
  $N$, that is, the submodule of elements $n \in N$ such that $rn = 0$
  (cf.\ \Cref{subsec:linalg:torsion}). (This is the reason why Tor
  is called Tor.) [\Cref{subsec:linrep:the-tor-functors},
  \ref{exer:linrep:why-ext-is-called-ext-part1}]
\end{exer}
\begin{exer}{}{exer:linrep:tor1-of-quotients-iso-to-intersection-over-product}
  Let $I$, $J$ be ideals of $R$. Prove that
  $\Tor^R_1(R/I, R/J) \cong (I \cap J)/IJ$. (For example, this
  $\Tor^R_1$ vanishes if $I + J = R$, by \Cref{lem:irred:pairwise-comaximal-product-equals-intersection}.)  Prove
  that $\Tor^R_i(R/I, R/J)$ is isomorphic to $\Tor^R_{i-1}(I, R/J)$
  for $i > 1$.
\end{exer}
\begin{exer}{$?$}{exer:linrep:tor-vanishes-above-degree-1-over-pid}
  Let $M$, $N$ be modules over a PID $R$. Prove that
  $\Tor^R_i(M, N) = 0$ for $i \geq 2$. (Assume $M$, $N$ are finitely
  generated, for simplicity.)
  [\Cref{subsec:linrep:the-tor-functors}]
\end{exer}
\begin{exer}{}{exer:linrep:cyclic-module-flat-iff-free}
  Let $R$ be an integral domain. Prove that a cyclic $R$-module is
  flat if and only if it is free.
\end{exer}
\begin{exer}{$\lnot$}{exer:linrep:flatness-criterion-via-monomorphisms}
  The following criterion is quite useful.
  \begin{itemize}
  \item Prove that an $R$-module $M$ is flat if and only if every
    monomorphism of $R$-modules $A \hookrightarrow B$ induces a
    monomorphism of $R$-modules
    $A \otimes_R M \hookrightarrow B \otimes_R M$.
  \item Prove that it suffices to verify this condition for all
    finitely generated modules $B$.
  \item Prove that it suffices to verify this condition when $B = R$
    and $A = I$ is an ideal of $R$.
  \item Deduce that an $R$-module $M$ is flat if and only if the
    natural homomorphism $I \otimes_R M \to IM$ is an isomorphism for
    every ideal $I$ of $R$.
  \end{itemize}
  If you believe in Tor's, now you can also show that an $R$-module
  $M$ is flat if and only if $\Tor^R_1(R/I, M) = 0$ for all ideals $I$
  of $R$. [\ref{exer:linrep:flat-iff-torsion-free-over-pid}]
\end{exer}
\begin{exer}{}{exer:linrep:flat-iff-torsion-free-over-pid}
  Let $R$ be a PID. Prove that an $R$-module $M$ is flat if and only
  if it is torsion-free. (If $M$ is finitely generated, the
  classification theorem of
  \Cref{subsec:linalg:the-classification-theorem} makes this
  particularly easy.  Otherwise, use
  \Cref{exer:linrep:flatness-criterion-via-monomorphisms}.)
  Geometrically, this says roughly that an algebraic set fails to be
  `flat' over a nonsingular curve if and only if some component of the
  set is contracted to a point. This phenomenon is displayed in the
  picture in \Cref{exam:linrep:coordinate-ring-of-projection-example}.
\end{exer}
\begin{exer}{$\lnot$}{exer:linrep:flatness-is-local-property}
  (Cf.\ \Cref{exer:irred:localization-at-prime-is-local}.) Prove that
  flatness is a local property: an $R$-module $M$ is flat if and only
  if $M_{\mathfrak{p}}$ is a flat $R_{\mathfrak{p}}$-module for all
  prime ideals $\mathfrak{p}$, if and only if $M_{\mathfrak{m}}$ is a
  flat $R_{\mathfrak{m}}$-module for all maximal ideals
  $\mathfrak{m}$. (Hint: Use
  \Cref{exer:linrep:localization-is-exact-preserves-homology,exer:linrep:tensor-and-localization-commute}.
  The
  $\Rightarrow$ direction will be straightforward. For the converse,
  let $A \subseteq B$ be $R$-modules, and let $K$ be the kernel of the
  induced homomorphism $A \otimes_R M \to B \otimes_R M$. Prove that
  the kernel of the localized homomorphism
  $A_{\mathfrak{m}} \otimes_{R_{\mathfrak{m}}} M_{\mathfrak{m}} \to
  B_{\mathfrak{m}} \otimes_{R_{\mathfrak{m}}} M_{\mathfrak{m}}$ is
  isomorphic to $K_{\mathfrak{m}}$, and use
  \Cref{exer:irred:module-zero-iff-zero-locally}.)
  [\ref{exer:linrep:localization-commutes-with-tor}]
\end{exer}
\begin{exer}{$\lnot$}{exer:linrep:localization-commutes-with-tor}
  Let $M$, $N$ be $R$-modules, and let $S$ be a multiplicative subset
  of $R$. Use the definition of Tor given in
  \Cref{subsec:linrep:the-tor-functors} to show
  $S^{-1}\Tor^R_i(M, N) \cong \Tor^{S^{-1}R}_i(S^{-1}M,
  S^{-1}N)$. (Use
  \Cref{exer:linrep:localization-is-exact-preserves-homology}.)  Use
  this fact to give a leaner proof that flatness is a local property
  (\Cref{exer:linrep:flatness-is-local-property}). [\ref{exer:linrep:flat-iff-tor1-vanishes-at-maximal-ideals}]
\end{exer}
\begin{exer}{$?$}{exer:linrep:flat-in-exact-sequence-with-flat-quotient}
  Let $0 \to M \to N \to P \to 0$ be an exact sequence of $R$-modules,
  and assume that $P$ is flat.
  \begin{itemize}
  \item Prove that $M$ is flat if and only if $N$ is flat.
  \item Prove that for all $R$-modules $Q$, the induced sequence
    $0 \to M \otimes_R Q \to N \otimes_R Q \to P \otimes_R Q \to 0$ is
    exact.
  \end{itemize}
  [\ref{exer:linrep:fin-gen-flat-over-noeth-local-is-free},
  \Cref{subsec:linrep:duality-and-exactness}]
\end{exer}
\begin{exer}{$\lnot$}{exer:linrep:fin-gen-flat-over-noeth-local-is-free}
  Let $R$ be a commutative Noetherian local ring with (single) maximal
  ideal $\mathfrak{m}$, and let $M$ be a finitely generated flat
  $R$-module.
  \begin{itemize}
  \item Choose elements $m_1, \dots, m_r \in M$ whose cosets mod
    $\mathfrak{m}M$ are a basis of $M/\mathfrak{m}M$ as a vector space
    over the field $R/\mathfrak{m}$. By Nakayama's lemma,
    $M = \langle m_1, \dots, m_r \rangle$
    (\Cref{exer:linalg:nakayama-gives-generators-from-quotient-basis}).
  \item Obtain an exact sequence
    $0 \to N \to R^{\oplus r} \to M \to 0$, where $N$ is finitely
    generated.
  \item Prove that this sequence induces an exact sequence
    $0 \to N/\mathfrak{m}N \to (R/\mathfrak{m})^{\oplus r} \to
    M/\mathfrak{m}M \to 0$. (Use
    \Cref{exer:linrep:flat-in-exact-sequence-with-flat-quotient}.)
  \item Deduce that $N = 0$. (Nakayama.)
  \item Conclude that $M$ is free.
  \end{itemize}
  Thus, a finitely generated module over a (Noetherian\footnote{The
    Noetherian hypothesis is actually unnecessary, but it simplifies
    the proof by allowing the use of Nakayama's lemma.}) local ring is
  flat if and only if it is free. Compare with
  \Cref{exer:linalg:direct-summand-of-free-over-local-is-free}.
  [\ref{exer:linrep:flat-iff-tor1-vanishes-at-maximal-ideals},
  \ref{exer:linrep:combine-results-reprove-rank-thm},
  \ref{exer:linrep:fin-gen-over-noeth-locally-free-iff-projective-iff-flat}]
\end{exer}

\section{Base change}
\label{sec:linrep:base-change}


\begin{exer}{}{exer:linrep:flat-iff-tor1-vanishes-at-maximal-ideals}
  Let $R$ be a commutative Noetherian ring, and let $M$ be a finitely
  generated $R$-module. Prove that $M$ is flat if and only if
  $\Tor^R_1(M, R/\mathfrak{m}) = 0$ for every maximal ideal
  $\mathfrak{m}$ of $R$. (Use
  \Cref{exer:linrep:flatness-is-local-property}, and refine the
  argument you used in
  \Cref{exer:linrep:fin-gen-flat-over-noeth-local-is-free}; remember
  that Tor localizes, by
  \Cref{exer:linrep:localization-commutes-with-tor}. The Noetherian
  hypothesis is actually unnecessary, but the proofs are harder
  without it.)
\end{exer}

I have championed several times the viewpoint that deep properties of
a ring $R$ are encoded in the category $R\text{-}\categ{Mod}$ of
$R$-modules; one extreme position is to simply replace $R$ with
$R\text{-}\categ{Mod}$ as the main object of study. The question then
arises as to how to deal with ring homomorphisms from this point of
view, or more generally how the categories $R\text{-}\categ{Mod}$,
$S\text{-}\categ{Mod}$ of modules over two (commutative) rings $R$,
$S$ may relate to each other. The reader should expect this to happen
by way of functors between the two categories and that the situation
at the categorical level will be substantially richer than at the ring
level.

\subsection{Balanced maps}
\label{subsec:linrep:balanced-maps}
Before we can survey the basic definitions, we must upgrade our
understanding of tensor products. It turns out that $M \otimes_R N$
satisfies a more encompassing universal property than the one examined
in
\Cref{subsec:linrep:bilinear-maps-and-the-definition-of-tensor-product}. Let
$M$, $N$ be modules over a commutative ring $R$, as in
\Cref{subsec:linrep:bilinear-maps-and-the-definition-of-tensor-product},
and let $G$ be an abelian group, i.e., a $\Z$-module.

\begin{defn}{}{defn:linrep:balanced-map}
  A $\Z$-bilinear map $\phi : M \times N \to G$ is \emph{$R$-balanced}
  if for all $m \in M$, $n \in N$, $r \in R$,
  \[
    \phi(rm, n) = \phi(m, rn).
  \]
\end{defn}

If $G$ is an $R$-module and $\phi : M \times N \to G$ is $R$-bilinear,
then it is $R$-balanced.\footnote{Also note that if $R$ is not
  commutative and $M$, resp., $N$, carries a right-, resp., left-,
  $R$-module structure, then the notion of `balanced map' makes sense.
  This leads to the definition of tensor (as an abelian group) in the
  noncommutative case.} But in general the notion of `$R$-balanced'
appears to be quite a bit more general, since $G$ is not even required
to be an $R$-module. This may lead the reader to suspect that a
solution to the universal problem of factoring balanced maps may be a
different gadget than the `ordinary' tensor product, but we are in
luck in this case, and the ordinary tensor product does the universal
job for balanced maps as well.

To understand this, recall that we constructed $M \otimes_R N$ as a
quotient $M \otimes_R N = R^{\oplus(M \times N)}/K$, where $K$ is
generated by the relations necessary to imply that the map
$M \times N \to R^{\oplus(M \times N)} \to R^{\oplus(M \times N)}/K$
is $R$-bilinear. We have observed that every element of
$M \otimes_R N$ may be written as a linear combination of pure tensors
$\sum_i m_i \otimes n_i$; it follows that the group homomorphism
$\Z^{\oplus(M \times N)} \twoheadrightarrow M \otimes_R N$ defined on
generators by $(m, n) \mapsto m \otimes n$ is surjective; its kernel
$K_B$ consists of the combinations
$\sum_i (m_i, n_i) \in \Z^{\oplus(M \times N)}$ such that
$\sum_i (m_i, n_i) \in K$, where the sum on the right is viewed in
$R^{\oplus(M \times N)}$. The reader will verify
(\Cref{exer:linrep:flat-iff-tor1-vanishes-at-maximal-ideals}) that
$K_B$ is generated by elements of the form
\begin{align*}
  &(m, n_1 + n_2) - (m, n_1) - (m, n_2), \\
  &(m_1 + m_2, n) - (m_1, n) - (m_2, n), \\
  &(rm, n) - (m, rn)
\end{align*}
(with $m, m_1, m_2 \in M$, $n, n_1, n_2 \in N$, $r \in R$). Therefore,
we have an induced isomorphism of abelian groups
\[
  \tag{$*$} \frac{\Z^{\oplus(M \times N)}}{K_B} \cong
  \frac{R^{\oplus(M \times N)}}{K},
\]
which amounts to an alternative description of $M \otimes_R N$. The
point of this observation is that the group on the left-hand side of
$(*)$ is manifestly a solution to the universal problem of factoring
$\Z$-bilinear, $R$-balanced maps. Therefore, we have proved:

\begin{lem}{}{lem:linrep:balanced-map-factors-through-tensor}
  Let $R$ be a commutative ring; let $M$, $N$ be $R$-modules, and let
  $G$ be an abelian group. Then every $\Z$-bilinear, $R$-balanced map
  $\phi : M \times N \to G$ factors through $M \otimes_R N$; that is,
  there exists a unique group homomorphism
  $\tilde\phi : M \otimes_R N \to G$ such that
  $\tilde\phi(m \otimes n) = \phi(m, n)$ for all $m \in M$, $n \in N$.
\end{lem}

The universal property explored in
\Cref{subsec:linrep:bilinear-maps-and-the-definition-of-tensor-product}
is recovered as the statement that if $G$ is an $R$-module and $\phi$
is $R$-bilinear, then the induced group homomorphism
$M \otimes_R N \to G$ is in fact an $R$-linear map.
\begin{rem}{}{rem:linrep:balanced-map-noncommutative-case}
  Balanced maps $\phi : M \times N \to G$ may be defined as soon as
  $M$ is a right-$R$-module and $N$ is a left-$R$-module, even if $R$
  is not commutative: require $\phi(mr, n) = \phi(m, rn)$ for all
  $m \in M$, $n \in N$, $r \in R$. The abelian group defined by the
  left-hand side of $(*)$ still makes sense and is taken as the
  definition of the tensor product $M \otimes_R N$; but note that this
  does not carry an $R$-module structure in general. This structure is
  recovered if, e.g., $M$ is a two-sided $R$-module.
\end{rem}

\subsection{Bimodules; adjunction again}
\label{subsec:linrep:bimodules-adjunction-again}
The enhanced universal property for the tensor will allow us to
upgrade the adjunction formula given in \Cref{lem:linrep:hom-tensor-adjunction}. This
requires the introduction of yet another notion.

\begin{defn}{}{defn:linrep:tensor-product-with-bimodule-scalars}
  Let $R$, $S$ be two commutative\footnote{Once more, the
    noncommutative case would be very worthwhile pursuing, but (not
    without misgivings) I have decided otherwise. In this more general
    case one requires $N$ to be both a left-$R$-module and a right-$S$
    -module with the compatibility expressed by $(rn)s = r(ns)$ for
    all choices of $n \in N$, $r \in R$, $s \in S$. This type of
    bookkeeping is precisely what is needed in order to extend the
    theory to the noncommutative case.} rings. An
  \emph{$(R,S)$-bimodule} is an abelian group $N$ endowed with
  compatible $R$-module and $S$-module structures, in the sense that
  for all $n \in N$, $r \in R$, $s \in S$,
  \[
    r(sn) = s(rn).
  \]
\end{defn}

For example, as $R$ is commutative, every $R$-module $N$ is an
$(R,R)$-bimodule: for all $r_1, r_2 \in R$ and $n \in N$,
$r_1(r_2 n) = (r_1 r_2)n = (r_2 r_1)n = r_2(r_1 n)$.

If $M$ is an $R$-module and $N$ is an $(R,S)$-bimodule, then the
tensor product $M \otimes_R N$ acquires an $S$-module structure:
define the action of $s \in S$ on pure tensors $m \otimes n$ by
$s(m \otimes n) := m \otimes (sn)$, and extend to all tensors by
linearity. In fact, this gives $M \otimes_R N$ an $(R,S)$-bimodule
structure.

Similarly, if $N$ is an $(R,S)$-bimodule and $P$ is an $S$-module,
then the abelian group $\Hom_S(N, P)$ is an $(R,S)$-bimodule: the
$R$-module structure is defined by setting $(r\alpha)(n) = \alpha(rn)$
for all $r \in R$, $n \in N$, and $\alpha \in \Hom_S(N, P)$.

This is needed to even make sense of the promised upgrade of
adjunction.  As with most such results, the proof is not difficult
once one understands what the statement says.

\begin{lem}{}{lem:linrep:hom-tensor-adjunction-bimodule}
  Suppose $M$ is an $R$-module, $N$ is an $(R,S)$-bimodule, and $P$ is
  an $S$-module. Then there is a canonical isomorphism of abelian
  groups
  \[
    \Hom_R(M, \Hom_S(N, P)) \cong \Hom_S(M \otimes_R N, P).
  \]
\end{lem}
\begin{proof}
  Every element $\alpha \in \Hom_R(M, \Hom_S(N, P))$ determines a map
  $\phi : M \times N \to P$ via $\phi(m, \mathord{-}) := \alpha(m)$;
  $\phi$ is clearly $\Z$-bilinear. Further, for all $r \in R$,
  $m \in M$, $n \in N$:
  \[
    \phi(rm, n) = \alpha(rm)(n) = r\alpha(m)(n) = \alpha(m)(rn) =
    \phi(m, rn),
  \]
  where the second equality holds by the $R$-linearity of $\alpha$ and
  the third by the definition of the $R$-module structure on
  $\Hom_S(N, P)$. Thus $\phi$ is $R$-balanced. By \Cref{lem:linrep:balanced-map-factors-through-tensor},
  such a map determines (and is determined by) a homomorphism of
  abelian groups $\tilde\phi : M \otimes_R N \to P$, with
  $\phi(m, n) = \tilde\phi(m \otimes n)$. I claim that $\tilde\phi$ is
  $S$-linear. Indeed, for all $s \in S$ and all pure tensors
  $m \otimes n$,
  \begin{align*}
    \tilde\phi(s(m \otimes n)) &= \tilde\phi(m \otimes (sn))
                                 = \phi(m, sn) = \alpha(m)(sn) \\
                               &= s\alpha(m)(n) = s\phi(m, n) = s\tilde\phi(m \otimes n),
  \end{align*}
  where I have used the $S$-linearity of $\alpha(m)$. Thus
  $\tilde\phi \in \Hom_S(M \otimes_R N, P)$. Tracing the argument
  backwards, every element of $\Hom_S(M \otimes_R N, P)$ determines an
  element of $\Hom_R(M, \Hom_S(N, P))$, and these two correspondences
  are clearly inverses of each other.
\end{proof}

If $R = S$, we recover the adjunction formula of \Cref{lem:linrep:hom-tensor-adjunction};
note that in this case the isomorphism is clearly $R$-linear.

\subsection{Restriction and extension of scalars}
\label{subsec:linrep:restriction-and-extension-of-scalars}
Coming back to the theme mentioned at the beginning of this section,
consider the case in which we have a homomorphism $f : R \to S$ of
(commutative) rings. It is natural to look for functors between the
categories $R\text{-}\categ{Mod}$ and $S\text{-}\categ{Mod}$ of
modules over $R$, $S$, respectively. There is a rather simple-minded
functor from $S\text{-}\categ{Mod}$ to $R\text{-}\categ{Mod}$
(`restriction of scalars'), while tensor products allow us to define a
functor from $R\text{-}\categ{Mod}$ to $S\text{-}\categ{Mod}$
(`extension of scalars'). A third important functor
$R\text{-}\categ{Mod} \to S\text{-}\categ{Mod}$ may be defined, also
`extending scalars', but for which I do not know a good name.

\pheader{Restriction of scalars.} Let $f : R \to S$ be a ring
homomorphism, and let $N$ be an $S$-module. Recall
(\Cref{subsec:rings:definition-of-left-r-module}) that this means
that we have chosen an action of the ring $S$ on the abelian group
$N$, that is, a ring homomorphism
$\sigma : S \to \mathrm{End}_{\categ{Ab}}(N)$.  Composing with $f$,
$\sigma \circ f : R \to S \to \mathrm{End}_{\categ{Ab}}(N)$ defines an
action of $R$ on the abelian group $N$, and hence an $R$-module
structure on $N$. More explicitly, if $r \in R$ and $n \in N$, define
the action of $r$ on $n$ by setting $rn := f(r)n$.  Since $S$ is
commutative, this defines in fact an $(R,S)$-bimodule structure on
$N$. Further, $S$-linear homomorphisms are in particular $R$-linear;
this assignment is (covariantly) functorial
$S\text{-}\categ{Mod} \to R\text{-}\categ{Mod}$.

If $f$ is injective, so that $R$ may be viewed as a subring of $S$,
then all we are doing is viewing $N$ as a module on a `restricted'
range of scalars, hence the terminology. For example, this is how we
view a complex vector space as a real vector space, in the simplest
possible way.

I will denote by $f_*$ this functor
$S\text{-}\categ{Mod} \to R\text{-}\categ{Mod}$ induced from $f$ by
restriction of scalars. Note that $f_*$ is trivially exact, because
the kernels and images of a homomorphism of modules are the same
regardless of the base ring. In view of the considerations in
\Cref{exam:linrep:exact-functor-definition-example}, this hints that $f_*$ may have both a
left-adjoint and a right-adjoint functor, and this will be precisely
the case (\Cref{prop:linrep:restriction-extension-coextension-adjunctions}).

\pheader{Extension of scalars} is defined from $R\text{-}\categ{Mod}$
to $S\text{-}\categ{Mod}$, by associating to an $R$-module $M$ the
tensor product $f^*(M) := M \otimes_R S$, which (as we have seen in
\Cref{subsec:linrep:bimodules-adjunction-again}) carries naturally
an $S$-module structure. This association is evidently covariantly
functorial. If $R^{\oplus B} \to R^{\oplus A} \to M \to 0$ is a
presentation of $M$ (cf.\
\Cref{subsec:linalg:finitely-presented-modules-and-free-resolutions}),
tensoring by $S$ gives (by the right-exactness of tensor) a
presentation of $f^*(M)$:
$S^{\oplus B} \to S^{\oplus A} \to M \otimes_R S \to 0$.  Intuitively,
this says that $f^*(M)$ is the module defined by `the same generators
and relations' as $M$, but with coefficients in $S$.

The third functor, denoted $f^!$, also acts from
$R\text{-}\categ{Mod}$ to $S\text{-}\categ{Mod}$ and is yet another
natural way to combine the ingredients we have at our disposal: if $M$
is an $R$-module, I have pointed out in
\Cref{subsec:linrep:bimodules-adjunction-again} that
$f^!(M) := \Hom_R(S, M)$ may be given a natural $S$-module structure
(by setting $s\alpha(s') := \alpha(ss')$). This is again evidently a
covariantly functorial prescription.
\begin{prop}{}{prop:linrep:restriction-extension-coextension-adjunctions}
  Let $f : R \to S$ be a homomorphism of commutative rings. Then, with
  notation as above, $f_*$ is right-adjoint to $f^*$ and left-adjoint
  to $f^!$. In particular, $f_*$ is exact, $f^*$ is right-exact, and
  $f^!$ is left-exact.
\end{prop}
\begin{proof}
  Let $M$, resp., $N$, be an $R$-module, resp., an $S$-module. Note
  that, trivially, $\Hom_S(S, N)$ is canonically isomorphic to $N$ (as
  an $S$-module) and to $f_*(N)$ (as an
  $R$-module). Thus\footnote{These are isomorphisms of abelian groups,
    and in fact isomorphisms of $R$-modules if one applies $f_*$ to
    the $\Hom_S$ terms.}
  \begin{align*}
    \Hom_R(M, f_*(N))
    &\cong \Hom_R(M, \Hom_S(S, N)) \\
    &\cong \Hom_S(M \otimes_R S, N) = \Hom_S(f^*(M), N),
  \end{align*}
  where I have used \Cref{lem:linrep:hom-tensor-adjunction-bimodule}. These bijections are canonical,
  proving that $f^*$ is left-adjoint to $f_*$.

  Similarly, there is a canonical isomorphism $N \cong N \otimes_S S$
  (\Cref{exam:linrep:tensor-with-base-ring-is-identity}); thus $N \otimes_S S \cong f_*(N)$ as
  $R$-modules, and for every $R$-module $M$,
  \begin{align*}
    \Hom_R(f_*(N), M)
    &\cong \Hom_R(N \otimes_S S, M) \\
    &\cong \Hom_S(N, \Hom_R(S, M))
      = \Hom_S(N, f^!(M))
  \end{align*}
  again by \Cref{lem:linrep:hom-tensor-adjunction-bimodule}. This shows that $f^!$ is right-adjoint
  to $f_*$, concluding the proof.
\end{proof}

\begin{rem}{(Warning)}{rem:linrep:notation-warning-pushforward}
  My choice of notation, $f_*$, etc., is somewhat nonstandard, and the
  reader should not take it too literally. It is inspired by analogs
  in the context of sheaf theory over schemes, but some of the
  properties reviewed above require crucial adjustments in that wider
  context: for example $f_*$ is not exact as a sheaf operation on
  schemes.
\end{rem}

\subsection*{Exercises}

In the following exercises, $R$, $S$ denote commutative rings.

\begin{exer}{$?$}{exer:linrep:tensor-relations-via-free-abelian-presentation}
  Verify that a combination of pure tensors $\sum_i (m_i \otimes n_i)$
  is zero in the tensor product $M \otimes_R N$ if and only if
  $\sum_i (m_i, n_i) \in \Z^{\oplus(M \times N)}$ is a combination of
  elements of the form
  \begin{align*}
    &(m, n_1 + n_2) - (m, n_1) - (m, n_2), \\
    &(m_1 + m_2, n) - (m_1, n) - (m_2, n), \\
    &(rm, n) - (m, rn),
  \end{align*}
  with $m, m_1, m_2 \in M$, $n, n_1, n_2 \in N$, $r \in R$.
  [\Cref{subsec:linrep:balanced-maps}]
\end{exer}
\begin{exer}{}{exer:linrep:canonical-map-tensor-restriction-of-scalars}
  If $f : R \to S$ is a ring homomorphism and $M$, $N$ are $S$-modules
  (hence $R$-modules by restriction of scalars), prove that there is a
  canonical homomorphism of $R$-modules
  $M \otimes_R N \to M \otimes_S N$.
\end{exer}
\begin{exer}{$?$}{exer:linrep:associativity-of-tensor-product}
  Let $R$, $S$ be commutative rings, and let $M$ be an $R$-module, $N$
  an $(R,S)$-bimodule, and $P$ an $S$-module. Prove that there is an
  isomorphism of $R$-modules
  \[
    M \otimes_R (N \otimes_S P) \cong (M \otimes_R N) \otimes_S P.
  \]
  In this sense, $\otimes$ is
  `associative'. [\ref{exer:linrep:reprove-localization-tensor-via-associativity},
  \Cref{subsec:linrep:multilinear-symmetric-alternating-maps}]
\end{exer}
\begin{exer}{$?$}{exer:linrep:reprove-localization-tensor-via-associativity}
  Use the associativity of the tensor product
  (\Cref{exer:linrep:canonical-map-tensor-restriction-of-scalars}) to
  prove again the formulas given in
  \Cref{exer:linrep:tensor-and-localization-commute}. (Use
  \Cref{exer:linrep:localization-is-tensor-with-localized-ring}.)
  [\ref{exer:linrep:tensor-and-localization-commute}]
\end{exer}
\begin{exer}{}{exer:linrep:base-change-functors-commute-with-limits}
  Let $f : R \to S$ be a ring homomorphism. Prove that $f^!$ commutes
  with limits, $f^*$ commutes with colimits, and $f_*$ commutes with
  both. In particular, deduce that these three functors all preserve
  finite direct sums.
\end{exer}
\begin{exer}{}{exer:linrep:restriction-of-scalars-conservative-faithfully-exact}
  Let $f : R \to S$ be a ring homomorphism, and let
  $\phi : N_1 \to N_2$ be a homomorphism of $S$-modules. Prove that
  $\phi$ is an isomorphism if and only if $f_*(\phi)$ is an
  isomorphism. (Functors with this property are said to be
  \emph{conservative}.) In fact, prove that $f_*$ is faithfully exact:
  a sequence of $S$-modules $0 \to L \to M \to N \to 0$ is exact if
  and only if the sequence of $R$-modules
  $0 \to f_*(L) \to f_*(M) \to f_*(N) \to 0$ is exact.  In particular,
  a sequence of $R$-modules is exact if and only if it is exact as a
  sequence of abelian groups.
\end{exer}
\begin{exer}{}{exer:linrep:dim-of-restriction-of-scalars-field-ext}
  Let $\iota : k \subseteq F$ be a finite field extension, and let $W$
  be an $F$-vector space of finite dimension $n$. Compute the
  dimension of $\iota_*(W)$ as a $k$-vector space (where $\iota_*$ is
  restriction of scalars; cf.\
  \Cref{subsec:linrep:restriction-and-extension-of-scalars}).
\end{exer}
\begin{exer}{}{exer:linrep:dim-of-extension-of-scalars-field-ext}
  Let $\iota : k \subseteq F$ be a finite field extension, and let $V$
  be a $k$-vector space of dimension $n$. Compute the dimension of
  $\iota^*(V)$ and $\iota^!(V)$ as $F$-vector spaces.
\end{exer}
\begin{exer}{}{exer:linrep:universal-property-of-extension-of-scalars}
  Let $f : R \to S$ be a ring homomorphism, and let $M$ be an
  $R$-module. Prove that the extension $f^*(M)$ satisfies the
  following universal property: if $N$ is an $S$-module and
  $\phi : M \to N$ is an $R$-linear map, then there exists a unique
  $S$-linear map $\tilde\phi : f^*(M) \to N$ making the diagram
  \[
    \begin{array}{ccc}
      M & \xrightarrow{\phi} & N \\
      \downarrow\iota & & \nearrow\tilde\phi \\
      f^*(M) & &
    \end{array}
  \]
  commute, where $\iota : M \to f^*(M) = M \otimes_R S$ is defined by
  $m \mapsto m \otimes 1$. (Thus, $f^*(M)$ is the `best approximation'
  to the $R$-module $M$ in the category of $S$-modules.)
\end{exer}
\begin{exer}{}{exer:linrep:projection-formula-base-change}
  Prove the following \emph{projection formula}: if $f : R \to S$ is a
  ring homomorphism, $M$ is an $R$-module, and $N$ is an $S$-module,
  then $f_*(f^*(M) \otimes_S N) \cong M \otimes_R f_*(N)$ as
  $R$-modules.
\end{exer}
\begin{exer}{}{exer:linrep:extension-of-scalars-preserves-flatness}
  Let $f : R \to S$ be a ring homomorphism, and let $M$ be a flat
  $R$-module. Prove that $f^*(M)$ is a flat $S$-module.
\end{exer}
\begin{exer}{}{exer:linrep:composition-rules-for-base-change-functors}
  In `geometric' contexts (such as the one hinted at in
  \Cref{rem:linrep:notation-warning-pushforward}), one would actually work with categories
  which are opposite to the category of commutative rings; cf.\
  \Cref{exam:linrep:category-of-affine-algebraic-sets}. A ring homomorphism $f : R \to S$
  corresponds to a morphism $f^\circ : S^\circ \to R^\circ$ in the
  opposite category, and we can simply define $(f^\circ)_*$, etc., to
  be $f_*$, etc. For morphisms $f^\circ : S^\circ \to R^\circ$ and
  $g^\circ : T^\circ \to S^\circ$ in the opposite category, prove that
  \begin{itemize}
  \item
    $(f^\circ \circ g^\circ)_* \cong (f^\circ)_* \circ (g^\circ)_*$,
  \item
    $(f^\circ \circ g^\circ)^* \cong (g^\circ)^* \circ (f^\circ)^*$,
  \item
    $(f^\circ \circ g^\circ)^! \cong (g^\circ)^! \circ (f^\circ)^!$,
  \end{itemize}
  where $\cong$ stands for `naturally isomorphic'.
\end{exer}
\begin{exer}{}{exer:linrep:compute-base-change-functors-examples}
  Let $p > 0$ be a prime integer, and let $\pi : \Z \to \Z/p\Z$ be the
  natural projection. Compute $\pi^*(A)$ and $\pi^!(A)$ for all
  finitely generated abelian groups $A$, as a vector space over
  $\Z/p\Z$. Compute $\iota^*(A)$ and $\iota^!(A)$ for all finitely
  generated abelian groups $A$, where $\iota : \Z \hookrightarrow \Q$
  is the natural inclusion.
\end{exer}
\begin{exer}{}{exer:linrep:onto-hom-explicit-restriction-extension-scalars}
  Let $f : R \to S$ be an onto ring homomorphism; thus, $S \cong R/I$
  for some ideal $I$ of $R$.
  \begin{itemize}
  \item Prove that, for all $R$-modules $M$,
    $f^!(M) \cong \{m \in M \mid \forall a \in I,\, am = 0\}$, while
    $f^*(M) \cong
    M/IM$. (\Cref{exer:rings:hom-contravariant-left-exact-from-ses}
    may help.)
  \item Prove that, for all $S$-modules $N$, $f^! f_*(N) \cong N$ and
    $f^* f_*(N) \cong N$.
  \item Prove that $f_*$ is fully faithful (\Cref{defn:linrep:faithful-full-functor}).
  \item Deduce that if there is an onto homomorphism $R \to S$, then
    $S\text{-}\categ{Mod}$ is equivalent to a full subcategory of
    $R\text{-}\categ{Mod}$.
  \end{itemize}
\end{exer}
\begin{exer}{}{exer:linrep:morita-equivalence-and-center}
  Let $f : R \to S$ be a homomorphism of (commutative) rings, and
  assume that the functor
  $f_* : S\text{-}\categ{Mod} \to R\text{-}\categ{Mod}$ is an
  equivalence of categories.
  \begin{itemize}
  \item Prove that there is a homomorphism of rings
    $g : S \to \mathrm{End}_{\categ{Ab}}(R)$ such that the composition
    $R \to S \to \mathrm{End}_{\categ{Ab}}(R)$ is the homomorphism
    realizing $R$ as a module over itself (that is, the homomorphism
    studied in \Cref{prop:rings:Cayley-ring}).
  \item Use the facts that $S$ is commutative and $f_*$ is fully
    faithful to deduce that $g(S)$ is isomorphic to $R$. Deduce that
    $f$ has a left-inverse $g : S \to R$.
  \item Therefore, $f_* \circ g_*$ is naturally isomorphic to the
    identity; and it follows that $g_* \circ f_*(S) \cong S$ as an
    $S$-module. Prove that this implies that $g$ is injective. (If
    $a \in \ker g$, prove that $a$ is in the annihilator of
    $g_* \circ f_*(S)$.)
  \item Conclude that $f$ is an isomorphism.
  \end{itemize}
  Two rings are \emph{Morita equivalent} if their categories of
  left-modules are equivalent. The result of this exercise is a (very)
  particular case of the fact that two commutative rings are Morita
  equivalent if and only if they are isomorphic. In fact, this more
  general statement is perhaps easier (!) to prove than the particular
  case worked out in this exercise. Recall
  (\Cref{exer:linrep:center-of-mod-category-iso-to-center-of-ring})
  that if $R$ is any ring, then the center of $R\text{-}\categ{Mod}$
  is isomorphic to the center of $R$. It is not difficult to deduce
  from this that if $R\text{-}\categ{Mod}$ and $S\text{-}\categ{Mod}$
  are equivalent, then the centers of $R$ and $S$ are isomorphic. So
  if two commutative rings $R$ and $S$ are Morita equivalent, then
  they must be isomorphic.  The commutativity is crucial in this
  statement: for example, it can be shown that any ring $R$ is Morita
  equivalent to the ring of matrices\footnote{I was once told that
    $M_{n,n}(\C)$ is `not seriously noncommutative' since it is Morita
    equivalent to $\C$, which is commutative.}  $M_{n,n}(R)$, for all
  $n > 0$.
\end{exer}

\section{Multilinear algebra}
\label{sec:linrep:multilinear-algebra}


\subsection{Multilinear, symmetric, alternating maps}
\label{subsec:linrep:multilinear-symmetric-alternating-maps}
Multilinear maps may be defined similarly to bilinear maps: if
$M_1, \dots, M_\ell$, $P$ are $R$-modules, a function
\[
  \phi : M_1 \times \cdots \times M_\ell \to P
\]
is \emph{$R$-multilinear} if it is $R$-linear in each factor, that is,
if the function obtained by arbitrarily fixing all but the $i$-th
component is $R$-linear in the $i$-th factor, for $i = 1, \dots, \ell$.

Again it is natural to ask whether $R$-multilinear maps may be turned
into $R$-linear maps: whether there exists an $R$-module
$M_1 \otimes_R \cdots \otimes_R M_\ell$ through which every
$R$-multilinear map must factor. Luckily, this module is already
available to us:

\begin{prop}{}{prop:linrep:multilinear-map-factors-through-iterated-tensor}
  Every $R$-multilinear map $M_1 \times \cdots \times M_\ell \to P$
  factors uniquely through
  $((\cdots(M_1 \otimes_R M_2) \otimes_R \cdots) \otimes_R M_{\ell-1})
  \otimes_R M_\ell$.
\end{prop}

Indeed, argue inductively: if
$\phi : M_1 \times \cdots \times M_\ell \to P$ is multilinear, then
for every $m_\ell \in M_\ell$, the function
$(m_1, \dots, m_{\ell-1}) \mapsto \phi(m_1, \dots, m_{\ell-1},
m_\ell)$ is $R$-multilinear, so it factors through
$(\cdots(M_1 \otimes_R M_2) \otimes_R \cdots) \otimes_R M_{\ell-1}$;
therefore, $\phi$ induces a unique bilinear map
$((\cdots(M_1 \otimes_R M_2) \otimes_R \cdots) \otimes_R M_{\ell-1})
\times M_\ell \to P$, and this induces a unique linear map
$((\cdots(M_1 \otimes_R M_2) \otimes_R \cdots) \otimes_R M_{\ell-1})
\otimes_R M_\ell \to P$ by the universal property of $\otimes_R$.

Of course the choice of pivoting on the last factor in this argument
is arbitrary: any other way to associate the factors would produce a
solution to the same universal problem. For $\ell = 3$, it follows in
particular that there are canonical isomorphisms
$(M_1 \otimes_R M_2) \otimes_R M_3 \cong M_1 \otimes_R (M_2 \otimes_R
M_3)$ for all $R$-modules $M_1$, $M_2$, $M_3$. In this sense, the
tensor product is associative\footnote{In fact, the tensor product is
  associative in a fancier sense; cf.\
  \Cref{exer:linrep:canonical-map-tensor-restriction-of-scalars}.} in
$R\text{-}\categ{Mod}$, and we are indeed authorized to use the
notation $M_1 \otimes_R \cdots \otimes_R M_\ell$.

Elements of this module are finite linear combinations
$\sum_i m_{1,i} \otimes \cdots \otimes m_{\ell,i}$ of pure tensors.
The structure map
$M_1 \times \cdots \times M_\ell \to M_1 \otimes_R \cdots \otimes_R
M_\ell$ acts as
$(m_1, \dots, m_\ell) \mapsto m_1 \otimes \cdots \otimes m_\ell$. By
the multilinearity of this map, we have, e.g.,
$(rm_1) \otimes \cdots \otimes m_\ell = m_1 \otimes \cdots \otimes
(rm_\ell) = r(m_1 \otimes \cdots \otimes m_\ell)$.  By convention, the
tensor product over $0$ factors is taken to be the base ring $R$.

Other important variations on the tensoring theme present themselves
when all the factors coincide. I will use the shorthand notation
\[
  \mathrm{T}^\ell_R(M) := M^{\otimes \ell} = \underbrace{M \otimes_R
    \cdots \otimes_R M}_{\ell \text{ times}}
\]
for the $\ell$-fold tensor product of a module $M$ by itself; this is
called the ($\ell$-th) \emph{tensor power} of $M$. For every
$R$-module $P$, every $R$-multilinear map $M^\ell \to P$ factors
uniquely through $\mathrm{T}^\ell_R(M)$: We set
$\mathrm{T}^1_R(M) = M$; by our convention, $\mathrm{T}^0_R(M) = R$.

Now we can impose further restrictions on $\phi$ and again study the
corresponding universal problems. Popular possibilities are:
\begin{itemize}
\item $\phi$ may be required to be \emph{symmetric}; that is, for all
  $\sigma \in S_\ell$ and all $m_1, \dots, m_\ell$, require that
  \[
    \phi(m_{\sigma(1)}, \dots, m_{\sigma(\ell)}) = \phi(m_1, \dots,
    m_\ell);
  \]
\item or $\phi$ may be required to be \emph{alternating}; that is,
  \[
    \phi(m_1, \dots, m_\ell) = 0 \quad \text{whenever } m_i = m_j
    \text{ for some } i \neq j.
  \]
\end{itemize}
The reader may have expected the second prescription to read: for all
$\sigma \in S_\ell$ and all $m_1, \dots, m_\ell$, require that
$\phi(m_{\sigma(1)}, \dots, m_{\sigma(\ell)}) = (-1)^\sigma \phi(m_1,
\dots, m_\ell)$. The problem with this version is that if the base
ring has characteristic 2, then there would be no difference between
symmetric and alternating maps. But behold the following:

\begin{lem}{}{lem:linrep:alternating-multilinear-sign-permutation}
  Let $\phi : M^\ell \to P$ be an $R$-multilinear function. If $\phi$
  is alternating, then for all $\sigma \in S_\ell$ and all
  $m_1, \dots, m_\ell$,
  \[
    \phi(m_{\sigma(1)}, \dots, m_{\sigma(\ell)}) = (-1)^\sigma
    \phi(m_1, \dots, m_\ell).
  \]
  If $2$ is a unit in $R$, the converse holds as well.
\end{lem}
\begin{proof}
  For the first statement, it suffices to show that interchanging any
  two factors switches the sign of an alternating function (since
  transpositions generate the symmetric group). Since the other
  factors have no effect on this operation, this reduces the question
  to the case $\ell = 2$. Therefore, we only have to show that if
  $\phi(m, m) = 0$ for all $m \in M$, then
  $\phi(m_2, m_1) = -\phi(m_1, m_2)$ for all $m_1, m_2 \in M$. For
  this, use bilinearity to expand $0 = \phi(m_1 + m_2, m_1 + m_2)$,
  and again use the alternating condition.

  For the second statement, again we can reduce to the $\ell = 2$
  case.  For $m_1 = m_2 = m$, $\phi(m_2, m_1) = -\phi(m_1, m_2)$ says
  $\phi(m, m) = -\phi(m, m)$; therefore $2\phi(m, m) = 0$. If $2$ is a
  unit in $R$, this implies $\phi(m, m) = 0$, as needed.
\end{proof}

Thus, alternating maps do satisfy the (perhaps) more natural
prescription, and this in fact characterizes alternating maps in most
cases.

\subsection{Symmetric and exterior powers}
\label{subsec:linrep:symmetric-and-exterior-powers}
The symmetric, resp., alternating, conditions come (of course) with
companion universal objects: modules $\mathrm{S}^\ell_R(M)$,
$\bigwedge^\ell_R(M)$ (symmetric and exterior power\footnote{The
  module $\bigwedge^\ell_R(M)$ is also called the $\ell$-th wedge
  power of $M$.}, respectively). The construction of these modules
follows patterns with which the reader is hopefully thoroughly
familiar by now.  For example, let $W \subseteq \mathrm{T}^\ell_R(M)$
be the submodule generated by all pure tensors
$m_1 \otimes \dots \otimes m_\ell$ such that $m_i = m_j$ for some
$i \neq j$; then
\[
  \bigwedge^\ell_R(M) := \frac{\mathrm{T}^\ell_R(M)}{W}
\]
satisfies the expected universal property for alternating maps, i.e.,
every multilinear, alternating map $\phi : M^\ell \to P$ induces a
unique $R$-linear map $\bar\phi$:
\[
  \begin{tikzcd}[column sep=large]
    M^\ell \ar[r, "\phi"] \ar[d, "\wedge"'] & P \\
    \bigwedge^\ell_R(M) \ar[ur, dashed, "\exists!\,\bar\phi"']
  \end{tikzcd}
\]
Here, the map $\wedge$ is of course the composition
\[
  M \times \dots \times M \to M \otimes \dots \otimes M \to \frac{M
    \otimes \dots \otimes M}{W} = \bigwedge^\ell_R(M),
\]
and $W$ is the smallest submodule making the multilinear map $\wedge$
alternating. The image of a pure tensor is denoted by using $\wedge$
rather than $\otimes$:
$m_1 \otimes \dots \otimes m_\ell \mapsto m_1 \wedge \dots \wedge
m_\ell$.  The reader will verify
(\Cref{exer:linrep:wedge-power-universal-property-and-sym-power}) that
$\bigwedge^\ell_R(M)$ may indeed be constructed as I have just
claimed, and the reader will produce an analogous construction for the
symmetric power $\mathrm{S}^\ell_R(M)$.

The module $\bigwedge^\ell_R(M)$ is generated by the pure
`alternating' tensors. By \Cref{lem:linrep:alternating-multilinear-sign-permutation},
$m_2 \wedge m_1 \wedge m_3 = -m_1 \wedge m_2 \wedge m_3$ (for
example), and, of course, $m_1 \wedge \dots \wedge m_\ell = 0$ if two
of the $m_i$'s coincide. It follows that if $e_1, \dots, e_r$ generate
$M$, then $\bigwedge^\ell_R(M)$ is generated by all
$e_{i_1} \wedge \dots \wedge e_{i_\ell}$ with
$1 \leq i_1 < \dots < i_\ell \leq r$. In particular, if $M$ is
finitely generated, then $\bigwedge^\ell_R(M) = 0$ for $\ell \gg 0$.
If $M$ is free, we can be very precise:
\begin{lem}{}{lem:linrep:exterior-power-of-free-module-rank}
  Let $R$ be a commutative ring, and let $M$ be a free $R$-module of
  rank $r$. Then $\bigwedge^\ell_R(M)$ is a free $R$-module of rank
  $\binom{r}{\ell}$.
\end{lem}
\begin{proof}
  There are $\binom{r}{\ell}$ sequences of indices
  $i_1, \dots, i_\ell$ satisfying
  $1 \leq i_1 < \dots < i_\ell \leq r$, so we just need to show that
  the generators $e_{i_1} \wedge \dots \wedge e_{i_\ell}$ are linearly
  independent.

  For a fixed $I = (i_1, \dots, i_\ell)$ with
  $1 \leq i_1 < \dots < i_\ell \leq r$, consider the map
  $\phi_I : M^\ell \to R$ obtained by setting
  \[
    \phi_I(e_{j_1}, \dots, e_{j_\ell}) = \begin{cases} (-1)^\sigma &
      \text{if } \exists \text{ a permutation } \sigma
                                                                     \text{ s.t.\ } \sigma(j_k) = i_k \text{ for all } k, \\
      0 & \text{otherwise}
    \end{cases}
  \]
  and extending by multilinearity. Note that $\sigma$ is unique if it
  exists, so the prescribed value of $\phi_I$ is well-defined. Also
  note that since $M$ is free, every element of $M$ is expressed
  uniquely as a combination of the $e_i$'s; hence $\phi_I$ is
  well-defined for all elements of $M^\ell$. Further, $\phi_I$ is
  evidently alternating
  (\Cref{exer:linrep:multilinear-map-phi-i-alternating}). Thus, it
  factors uniquely through an $R$-linear map
  $\phi_I : \bigwedge^\ell_R(M) \to R$, with the property that
  \[
    \phi_I(e_{j_1} \wedge \dots \wedge e_{j_\ell}) = \begin{cases}
      (-1)^\sigma & \text{if } \exists \text{ a permutation } \sigma
                    \text{ s.t.\ } \sigma(j_k) = i_k \text{ for all } k, \\
      0 & \text{otherwise.}
    \end{cases}
  \]
  Now assume that
  $\sum_{1 \leq j_1 < \dots < j_\ell \leq r} \lambda_{j_1 \dots
    j_\ell}\, e_{j_1} \wedge \dots \wedge e_{j_\ell} = 0$; then with
  $I = (i_1, \dots, i_\ell)$
  \[
    \lambda_{i_1 \dots i_\ell} = \phi_I\!\Bigl( \sum_{1 \leq j_1 <
      \dots < j_\ell \leq r} \lambda_{j_1 \dots j_\ell}\, e_{j_1}
    \wedge \dots \wedge e_{j_\ell} = 0 \Bigr) = \phi_I(0) = 0.
  \]
  This shows that there are no nontrivial linear dependence relations
  among the generators $e_{i_1} \wedge \dots \wedge e_{i_\ell}$, as
  claimed.
\end{proof}
Note that, in particular, it follows that if $M$ is a free $R$-module
of rank $r$, then $\bigwedge^r_R(M) \cong R$; an isomorphism is
induced by the map $\phi_{1\dots r} : M^r \to R$ obtained by setting
\[
  \phi_{1\dots r}(e_{i_1}, \dots, e_{i_r}) = \begin{cases} \pm 1 &
    \text{if the indices } i_1, \dots, i_r \text{ are
                                                                   distinct,} \\
    0 & \text{otherwise,}
  \end{cases}
\]
where the sign is determined by the permutation ordering the indices.
(This is a particular case of the maps $\phi_I$ considered in the
proof of \Cref{lem:linrep:exterior-power-of-free-module-rank}.)

\begin{prop}{}{prop:linrep:determinant-via-alternating-map}
  Let $A$ be an $r \times r$ matrix with entries in a field $k$, and
  let $a_i \in k^r$ be the $i$-th column of $A$. Then with notation as
  above, $\det(A) = \phi_{1\dots r}(a_1, \dots, a_r)$.
\end{prop}

This is immediate, since the standard properties of the determinant
(obtained in \Cref{subsec:linalg:the-determinant}) prove that it is
a multilinear, alternating map of the columns, and the determinant of
the identity matrix is $1$.  For this reason, the top exterior power
of a free module $F$ (that is, $\bigwedge^r_R(F)$ if $F$ has rank $r$)
is called the \emph{determinant} of $F$, denoted $\det(F)$.

More generally, the functions $\phi_I$ give an explicit isomorphism
between $\bigwedge^\ell_k(k^r)$ and $k^{\binom{r}{\ell}}$.
\begin{exam}{}{exam:linrep:exterior-square-of-k4}
  For $V = k^4$, $\bigwedge^2_k(V)$ has dimension $\binom{4}{2} = 6$.
  On `pure wedges' $a_1 \wedge a_2$, the isomorphism
  $\bigwedge^2_k(V) \to k^6$ works as follows. View the vectors $a_1$,
  $a_2$ as the two columns of a $4 \times 2$ matrix $A$; then send
  $a_1 \wedge a_2$ to the collection of the six $2 \times 2$ minors of
  $A$\footnote{This definition extends by linearity to the whole of
    $\bigwedge^2_k(V)$.}:
  \[
    A = \begin{pmatrix}
      a_1^1 & a_2^1 \\ a_1^2 & a_2^2 \\
      a_1^3 & a_2^3 \\ a_1^4 & a_2^4
    \end{pmatrix}
    \mapsto
    \begin{pmatrix}
      a_1^1 a_2^2 - a_1^2 a_2^1 \\
      a_1^1 a_2^3 - a_1^3 a_2^1 \\
      a_1^1 a_2^4 - a_1^4 a_2^1 \\
      a_1^2 a_2^3 - a_1^3 a_2^2 \\
      a_1^2 a_2^4 - a_1^4 a_2^2 \\
      a_1^3 a_2^4 - a_1^4 a_2^3
    \end{pmatrix}.
  \]
\end{exam}

Similarly to the alternating case, if $e_1, \dots, e_r$ generate $M$,
then the symmetric power $\mathrm{S}^\ell_R(M)$ is generated by the
(images\footnote{I am not aware of a universally accepted notation to
  denote `symmetric' tensors; a sensible choice would be a simple
  $\cdot$, as in the special case of polynomial rings.} of) all
tensors $e_{i_1} \otimes \dots \otimes e_{i_\ell}$ with
$1 \leq i_1 \leq \dots \leq i_\ell \leq r$. In this case all
$\mathrm{S}^\ell_R(M)$ are, in general, nonzero.

\subsection{Very small detour: Graded algebra}
\label{subsec:linrep:very-small-detour-graded-algebra}
We have obtained tensor, symmetric, and exterior powers of a module
$M$ for every nonneg\-ative integer $\ell$. As it happens, in each
case it is useful to consider `all at once'. The structures we obtain
by doing this are particular cases of a very important notion, which
would deserve much more space than we can allow.

A \emph{graded ring} is a (not necessarily commutative) ring $S$
endowed with a decomposition
\[
  S = \bigoplus_{i \geq 0} S_i
\]
of the abelian group $(S, +)$ into a direct sum of abelian groups
$S_i$, for nonnegative integers\footnote{Of course, much more general
  gradings may be considered, but this is the case that will occur in
  the applications in this section.} $i$, such that
$S_i \cdot S_j \subseteq S_{i+j}$. This prescription is parsed as
follows: identify $S_i$ with its image in $S$ by the natural map;
nonzero elements of $S_i$ are called \emph{homogeneous}, of degree
$i$; the condition prescribes that the product of two homogeneous
elements of degree $i$, $j$ is homogeneous, of degree $i + j$
(provided it is not $0$).
\begin{exam}{}{exam:linrep:polynomial-ring-grading}
  The polynomial ring $R[x_1, \dots, x_n]$ over any ring $R$ carries a
  natural grading, given by the (ordinary) degree of polynomials. We
  may write
  \[
    R[x_1, \dots, x_n] = R \oplus \langle x_1, \dots, x_n \rangle
    \oplus \langle x_1^2, x_1 x_2, \dots, x_n^2 \rangle \oplus \dots.
  \]
\end{exam}
The reader should have no difficulty providing the notions of graded
modules and algebras. For example, a graded ring $S = \bigoplus_i S_i$
is a \emph{graded algebra} over a graded ring $R = \bigoplus_i R_i$ if
it carries an action of $R$ (making it into a `conventional'
$R$-algebra) and further $R_i \cdot S_j \subseteq S_{i+j}$. In
particular, if $R = R_0$ (so that $R$ is a commutative ring viewed as
a graded ring by `concentrating' it in degree $0$), we are just
requiring each graded piece of $S$ to be an $R$-module.

For a concrete example, the commutative ring $R[x_1, \dots, x_n]$ is a
graded $R$-algebra, generated (as an algebra over $R$) by the piece of
degree $1$. This is an important template for many interesting
situations.

I cannot resist the temptation to mention another example from
commutative algebra:
\begin{exam}{}{exam:linrep:rees-algebra-of-an-ideal}
  Let $R$ be a commutative ring, and let $I \subseteq R$ be an ideal.
  Then the direct sum of $R$-modules
  \[
    \bigoplus_{\ell \geq 0} I^\ell = R \oplus I \oplus I^2 \oplus I^3
    \oplus \dots
  \]
  has a natural ring structure, since
  $I^i \cdot I^j \subseteq I^{i+j}$. The resulting graded $R$-algebra
  is called the \emph{Rees algebra} of $I$, $\mathrm{Rees}_R(I)$, and
  is of fundamental importance in algebraic geometry (as it translates
  into the notion of blow-up, a particularly important class of
  regular maps). The brave reader will try to explore the Rees algebra
  of the ideal $I = (x, y)$ in the polynomial ring $k[x, y]$ over a
  field $k$ (\Cref{exer:linrep:rees-algebra-example-iso-to-quotient}),
  which is the simplest interesting instance of this construction
  (well, the second simplest: the simplest example may be
  \Cref{exer:rings:rees-algebra}).
\end{exam}

The additional information carried by the grading on a ring is very
substantial, and the theory of graded rings and algebras is extremely
rich. For example, projective algebraic geometry may be developed by
using graded rings along the same lines we sketched in
\Cref{subsec:fields:a-little-affine-algebraic-geometry} for the
affine case. As I mentioned at the end of
\Cref{subsec:fields:a-little-affine-algebraic-geometry}, this
provides a `globalization' procedure which leads to important
technical advantages.

For the (very modest) purpose of this section, the following general
remarks are all we need.  Graded rings form a category under ring
homomorphisms, but they form a more interesting one if we require the
homomorphisms to preserve the grading; that is, if
$S = \bigoplus_i S_i$ and $T = \bigoplus_i T_i$ are graded rings,
consider ring homomorphisms $\phi : S \to T$ such that
$\phi(S_i) \subseteq T_i$. Such homomorphisms are called \emph{graded
  homomorphisms}. Of course, singling out a special class of
homomorphisms will then single out a special class of ideals, that is,
those ideals which appear as kernels of graded homomorphisms.

\begin{defn}{}{defn:linrep:homogeneous-ideal}
  An ideal $I$ of a graded ring $S = \bigoplus_i S_i$ is
  \emph{homogeneous} if $I = \bigoplus_i (I \cap S_i)$.
\end{defn}
\begin{lem}{}{lem:linrep:homogeneous-ideal-tfae}
  Let $S = \bigoplus_i S_i$ be a graded ring and let $I \subseteq S$
  be an ideal of $S$. Then the following are equivalent:
  \begin{enumerate}[label=(\roman*)]
  \item $I$ is homogeneous;
  \item if $s \in S$ and $s = \sum_i s_i$ is the decomposition of $s$
    into homogeneous elements $s_i \in S_i$, then
    $s \in I \iff s_i \in I$ for all $i$;
  \item $I$ admits a generating set consisting of homogeneous
    elements;
  \item $I$ is the kernel of a graded homomorphism.
  \end{enumerate}
\end{lem}
\begin{proof}
  (i)$\iff$(ii) is the very definition of homogeneous ideal;
  (ii)$\iff$(iii) is left to the reader
  (\Cref{exer:linrep:graded-algebra-lemma-equivalence}).

  (ii)$\implies$(iv): Assuming (ii) holds, define a grading on $S/I$
  by letting the piece of degree $i$ consist of ($0$ and) the cosets
  of the elements of degree $i$ in $S$. This is well-defined by (ii),
  and it is immediately checked that it induces a graded ring
  structure on $S/I$. The ideal $I$ is then the kernel of the graded
  homomorphism $S \to S/I$, verifying (iv).

  (iv)$\implies$(ii): Let $\phi : S \to T$ be a graded homomorphism,
  and assume $I = \ker\phi$. Let $s = \sum_i s_i \in \ker\phi$, with
  $s_i \in S_i$. Then $\sum_i \phi(s_i) = 0$ in $T$, and $\phi(s_i)$
  is homogeneous of degree $i$ for all $i$. It follows that
  $\phi(s_i) = 0$ for all $i$, that is, each $s_i$ is in
  $\ker\phi = I$, implying (ii).
\end{proof}

Note that a graded homomorphism $\phi : S \to T$ induces a
homomorphism of abelian groups $\phi_i : S_i \to T_i$ for each $i$.
The ideal $\ker\phi$ is then the direct sum of all ideals
$\ker\phi_i$.
\begin{exam}{}{exam:linrep:homogeneous-ideal-example}
  The ideal $I = (y - x^2)$ is not homogeneous in the ring $k[x, y]$,
  if this is given the grading by the usual degree: indeed,
  $y - x^2 \in I$ while $y \notin I$, contradicting condition (ii) of
  \Cref{lem:linrep:homogeneous-ideal-tfae}. The ideal $(y, y - x^2)$ is homogeneous,
  since it equals $(y, x^2)$ and both $y$, $x^2$ are homogeneous.

  The conventional grading of a polynomial ring is not the only
  option: we could decide to grade $k[x, y]$ by placing constants in
  degree $0$ and assigning degree $1$ to the indeterminate $x$ and
  degree $2$ to $y$. With such a grading, the ideal $(y - x^2)$ is
  homogeneous.
\end{exam}

Projective algebraic geometry (almost) follows the blueprint of affine
algebraic geometry reviewed in
\Cref{subsec:fields:a-little-affine-algebraic-geometry}, but using
only homogeneous ideals in $k[x_1, \dots, x_n]$ (taken with the usual
grading).

\subsection{Tensor algebras}
\label{subsec:linrep:tensor-algebras}
Going back to the context of multilinear algebra, consider again a
module $M$ over a commutative ring $R$. We define the \emph{tensor
  algebra} of $M$ as the graded $R$-algebra
\[
  \mathrm{T}^*_R(M) := \bigoplus_{\ell \geq 0} \mathrm{T}^\ell_R(M);
\]
the multiplication is defined on pure tensors by
\begin{align*}
  (m_1 \otimes \dots \otimes m_i) \cdot
  (n_1 \otimes \dots \otimes n_j)
  := m_1 \otimes \dots \otimes m_i \otimes n_1 \otimes \dots \otimes
  n_j
\end{align*}
and is extended by linearity. Similarly, we can define the
\emph{symmetric algebra} and the \emph{exterior algebra} of $M$:
\[
  \mathrm{S}^*_R(M) := \bigoplus_{\ell \geq 0} \mathrm{S}^\ell_R(M),
  \qquad \bigwedge^*_R(M) := \bigoplus_{\ell \geq 0}
  \bigwedge^\ell_R(M).
\]
\begin{rem}{}{rem:linrep:exterior-algebra-skew-commutative}
  The tensor algebra is not commutative; the symmetric algebra is
  commutative; and the exterior algebra is `skew-commutative' in the
  sense that if $\alpha \in \bigwedge^i_R(M)$ and
  $\beta \in \bigwedge^j_R(M)$, then
  \[
    \alpha \wedge \beta = (-1)^{ij}\, \beta \wedge \alpha
  \]
  (\Cref{exer:linrep:skew-commutativity-of-exterior-algebra}), where I
  use $\wedge$ to denote the operation in $\bigwedge^*_R(M)$. The
  reader should recognize this formula from the theory of differential
  forms; differential forms provide perhaps the most important example
  of an exterior algebra.
\end{rem}
The definitions of symmetric and exterior powers were concocted so as
to yield surjective graded homomorphisms of algebras
\[
  \mathrm{T}^*_R(M) \twoheadrightarrow \mathrm{S}^*_R(M), \qquad
  \mathrm{T}^*_R(M) \twoheadrightarrow \bigwedge^*_R(M).
\]
As we have learned in \Cref{lem:linrep:homogeneous-ideal-tfae}, the kernels of these
homomorphisms are certain homogeneous ideals of the tensor algebra;
these ideals must then be generated by homogeneous elements.

\begin{lem}{}{lem:linrep:symmetric-exterior-algebra-as-tensor-algebra-quotient}
  Let $I_S$, $I_\wedge \subseteq \mathrm{T}^*_R(M)$ be the ideals
  respectively generated by all elements of the form
  $(m \otimes n - n \otimes m)$ as $m, n \in M$ and by elements of the
  form $m \otimes m$ as $m \in M$. Then
  \[
    \mathrm{S}^*_R(M) \cong \frac{\mathrm{T}^*_R(M)}{I_S}, \qquad
    \bigwedge^*_R(M) \cong \frac{\mathrm{T}^*_R(M)}{I_\wedge}
  \]
  as graded $R$-algebras.
\end{lem}
\begin{proof}
  We have observed in
  \Cref{subsec:linrep:very-small-detour-graded-algebra} that the
  kernel of a graded homomorphism is the direct sum of the kernels of
  the induced homomorphisms in each degree; the statement of the lemma
  then follows easily from the explicit descriptions of the kernels of
  the canonical projections from the tensor powers to the symmetric
  and exterior powers, given in
  \Cref{subsec:linrep:symmetric-and-exterior-powers}.
\end{proof}
\begin{rem}{}{rem:linrep:symmetric-tensors-subalgebra-caveat}
  In some cases it is also possible to construct subalgebras of
  $\mathrm{T}^*_R(M)$ isomorphic to $\mathrm{S}^*_R(M)$ and
  $\bigwedge^*_R(M)$. For example, tensors that are invariant under
  the action of the symmetric group may be used to give a copy of
  $\mathrm{S}^*_R(M)$ inside $\mathrm{T}^*_R(M)$:
  $\frac{1}{2}(m_1 \otimes m_2 + m_2 \otimes m_1)$ would be such a
  `symmetric tensor'. However, the introduction of denominators poses
  distasteful restrictions on the base ring $R$, and for this reason I
  prefer the `quotient' constructions of \Cref{lem:linrep:symmetric-exterior-algebra-as-tensor-algebra-quotient}.
\end{rem}
It is now straightforward to establish universal properties satisfied
by the algebras defined above. Note that by construction there are
canonical isomorphisms of $M$ with $\mathrm{T}^1_R(M)$,
$\mathrm{S}^1_R(M)$, $\bigwedge^1_R(M)$, and in particular there are
canonical maps
\[
  M \to \mathrm{T}^*_R(M), \quad M \to \mathrm{S}^*_R(M), \quad M \to
  \bigwedge^*_R(M).
\]

\begin{prop}{}{prop:linrep:tensor-algebra-universal-property}
  Let $R$ be a commutative ring, and let $M$ be an $R$-module. Then
  for every $R$-algebra $T$ and every $R$-module homomorphism
  $\phi : M \to T$ there exists a unique homomorphism of $R$-algebras
  $\bar\phi : \mathrm{T}^*_R(M) \to T$ making the following diagram
  commute:
  \[
    \begin{tikzcd}
      M \ar[r, "\phi"] \ar[d] & T \\
      \mathrm{T}^*_R(M) \ar[ur, dashed, "\exists!\,\bar\phi"']
    \end{tikzcd}
  \]
\end{prop}

\begin{prop}{}{prop:linrep:symmetric-algebra-universal-property}
  Let $R$ be a commutative ring, and let $M$ be an $R$-module. Then
  for every commutative $R$-algebra $S$ and every $R$-module
  homomorphism $\psi : M \to S$ there exists a unique homomorphism of
  $R$-algebras $\bar\psi : \mathrm{S}^*_R(M) \to S$ making the
  following diagram commute:
  \[
    \begin{tikzcd}
      M \ar[r, "\psi"] \ar[d] & S \\
      \mathrm{S}^*_R(M) \ar[ur, dashed, "\exists!\,\bar\psi"']
    \end{tikzcd}
  \]
\end{prop}

\begin{prop}{}{prop:linrep:exterior-algebra-universal-property}
  Let $R$ be a commutative ring, and let $M$ be an $R$-module. Then
  for every $R$-algebra $A$ and every $R$-module homomorphism
  $\lambda : M \to A$ such that $\lambda(m)^2 = 0$ for all $m \in M$,
  there exists a unique homomorphism of $R$-algebras
  $\bar\lambda : \bigwedge^*_R(M) \to A$ making the following diagram
  commute:
  \[
    \begin{tikzcd}
      M \ar[r, "\lambda"] \ar[d] & A \\
      \bigwedge^*_R(M) \ar[ur, dashed, "\exists!\,\bar\lambda"']
    \end{tikzcd}
  \]
\end{prop}

Details may now safely be left to the reader (who may for example
establish the first proposition from the universal property of tensor
powers and then deduce the second and third from
\Cref{lem:linrep:symmetric-exterior-algebra-as-tensor-algebra-quotient}).
\begin{exam}{}{exam:linrep:symmetric-algebra-of-free-module}
  The free case is particularly easy to understand. For instance,
  $\mathrm{S}^*_R(R^{\oplus r}) \cong R[x_1, \dots, x_r]$. Indeed, the
  polynomial ring satisfies the appropriate universal property with
  respect to mapping to commutative rings (cf.\
  \Cref{subsec:rings:universal-property-of-polynomial-rings} and
  \Cref{subsec:rings:submodule-generated-by-a-subset-noetherian-modules}). Likewise,
  $\mathrm{T}^*_R(R^{\oplus r})$ should be thought of as a
  `noncommutative' polynomial ring, in which the $r$ indeterminates do
  not commute with each other (cf.\
  \Cref{subsec:rings:free-modules-and-free-algebras}).
\end{exam}

Of course the constructions $\mathrm{T}^*_R$, $\mathrm{S}^*_R$,
$\bigwedge^*_R$ are all functorial, and their behavior with respect to
exact sequences is interesting and important. For example, suppose
\[
  0 \to L \to M \to N \to 0
\]
is an exact sequence of $R$-modules. Then $L$ may be identified with a
subset of $M \cong \mathrm{S}^1_R(M)$; hence it defines an ideal
$L \cdot \mathrm{S}^*_R(M)$ of the algebra $\mathrm{S}^*_R(M)$. It is
not hard to show that the sequence\footnote{The shift
  $\mathrm{S}^{*-1}$ is introduced in order to preserve degrees in
  this sequence.}
\[
  0 \to L \cdot \mathrm{S}^{*-1}_R(M) \to \mathrm{S}^*_R(M) \to
  \mathrm{S}^*_R(N) \to 0
\]
is exact. This is often useful in computations (cf.\
\Cref{exer:linrep:rees-algebra-example-iso-to-quotient}). I will leave
any further such explorations to the more motivated readers.

\subsection*{Exercises}

In the following exercises, $R$ denotes a commutative ring and $k$
denotes a field.

\begin{exer}{$*$}{exer:linrep:wedge-power-universal-property-and-sym-power}
  Verify that the module $\bigwedge^\ell_R(M)$ constructed in
  \Cref{subsec:linrep:symmetric-and-exterior-powers} does satisfy
  the universal property for alternating multilinear maps.  Construct
  a module $\mathrm{S}^\ell_R(M)$ satisfying the universal property
  for symmetric multilinear
  maps. [\Cref{subsec:linrep:symmetric-and-exterior-powers}]
\end{exer}

\begin{exer}{}{exer:linrep:tensor-sym-wedge-powers-as-functors}
  Define the action of $\mathrm{T}^\ell_R$, $\mathrm{S}^\ell_R$,
  $\bigwedge^\ell_R$ on $R$-linear maps, making covariant functors out
  of these notions.
\end{exer}

\begin{exer}{}{exer:linrep:explicit-alternating-form-wedge2-vanishes-example}
  Let $I$ be the ideal $(x, y)$ in $k[x, y]$; so every element of $I$
  may be written (of course, not uniquely) in the form $fx + gy$ for
  some polynomials $f, g \in k[x, y]$. Define a function
  $\phi : I \times I \to k$ by prescribing
  \[
    \phi(f_1 x + g_1 y,\, f_2 x + g_2 y) := f_1(0, 0)\, g_2(0, 0) -
    f_2(0, 0)\, g_1(0, 0).
  \]
  \begin{itemize}
  \item Prove that $\phi$ is well-defined.
  \item Prove that $\phi$ is $k[x,y]$-bilinear and alternating.
  \item Prove that $\bigwedge^2_{k[x,y]}(I) \neq 0$.
  \end{itemize}
  Note that $I$ has rank $1$ as a $k[x,y]$-module; if it were free,
  its second exterior power would have to vanish by
  \Cref{lem:linrep:exterior-power-of-free-module-rank}.
\end{exer}

\begin{exer}{}{exer:linrep:det-and-wedge-of-direct-sum}
  Let $F_1$ and $F_2$ be free $R$-modules of finite rank.
  \begin{itemize}
  \item Construct a `meaningful' isomorphism
    $\det(F_1) \otimes \det(F_2) \cong \det(F_1 \oplus F_2)$.
  \item More generally, prove that
    \[
      \bigwedge^r_R(F_1 \oplus F_2) \cong \bigoplus_{i+j=r}
      \bigl(\bigwedge^i_R F_1\bigr) \otimes_R \bigl(\bigwedge^j_R
      F_2\bigr).
    \]
  \end{itemize}
\end{exer}

\begin{exer}{$*$}{exer:linrep:multilinear-map-phi-i-alternating}
  Verify that the multilinear map $\phi_I$ defined in the proof of
  \Cref{lem:linrep:exterior-power-of-free-module-rank} is
  alternating. [\Cref{subsec:linrep:symmetric-and-exterior-powers}]
\end{exer}
\begin{exer}{}{exer:linrep:linear-independence-iff-wedge-nonzero}
  Let $V$ be a vector space, and let $v_1, \dots, v_\ell \in V$. Prove
  that $v_1, \dots, v_\ell$ are linearly independent if and only if
  $v_1 \wedge \dots \wedge v_\ell \neq 0$.
\end{exer}

\begin{exer}{$\lnot$}{exer:linrep:plucker-embedding-well-defined}
  Let $V$ be a $k$-vector space, and let $\{v_1, \dots, v_\ell\}$,
  $\{w_1, \dots, w_\ell\}$ be two sets of linearly independent vectors
  in $V$. Prove that $\{v_1, \dots, v_\ell\}$ and
  $\{w_1, \dots, w_\ell\}$ span the same subspace of $V$ if and only
  if $v_1 \wedge \dots \wedge v_\ell$ and
  $w_1 \wedge \dots \wedge w_\ell$ are nonzero multiples of each other
  in $\bigwedge^\ell_k(V)$.  (For the interesting direction, if
  $\langle v_1, \dots, v_\ell \rangle = \langle w_1, \dots, w_\ell
  \rangle$, there must be a vector $u$ belonging to the first subspace
  but not the second. What can you say about
  $(v_1 \wedge \dots \wedge v_\ell) \wedge u$ and
  $(w_1 \wedge \dots \wedge w_\ell) \wedge u$ in $\bigwedge^*_k(V)$?)

  Deduce that there is an injective function from the Grassmannian of
  $\ell$-di\-men\-sional subspaces of $V$
  (\Cref{exer:linalg:grassmannian-cell-decomposition}) to the
  projective space $\mathbf{P}(\bigwedge^\ell_k V)$: the Grassmannian
  is identified with the set of `pure wedges' in the projectivization
  of the exterior power $\bigwedge^\ell_k(V)$. This is called the
  \emph{Pl\"ucker embedding} of the
  Grassmannian. [\ref{exer:linrep:grassmannian-2-4-as-projective-quadric}]
\end{exer}

\begin{exer}{}{exer:linrep:grassmannian-2-4-as-projective-quadric}
  The Pl\"ucker embedding described in
  \Cref{exer:linrep:plucker-embedding-well-defined} realizes
  the Grassmannian $\mathrm{Gr}_k(2,4)$ of $2$-dimensional subspaces
  of $k^4$ as a subset of the projectivization of
  $\bigwedge^2_k(k^4) \cong k^6$:
  $\mathrm{Gr}_k(2,4) \subseteq \mathbf{P}^5_k$. Choose projective
  coordinates (\Cref{exer:fields:affine-to-projective-space-covering})
  $(x_{12} : x_{13} : x_{14} : x_{23} : x_{24} : x_{34})$ in
  $\mathbf{P}(\bigwedge^2_k k^4)$, listed according to the
  corresponding minors, as in \Cref{exam:linrep:exterior-square-of-k4}. Prove that
  $\mathrm{Gr}_k(2,4)$ is the locus of zeros
  \[
    \mathbf{V}(x_{12}x_{34} - x_{13}x_{24} + x_{14}x_{23})
  \]
  (notation as in
  \Cref{exer:fields:homog-vanishing-well-defined-covering}). (Remember
  that you know how to write every point of $\mathrm{Gr}_k(2,4)$: look
  back at \Cref{exer:linalg:gr-2-4-union-of-six-schubert-cells}.)

  Thus, $\mathrm{Gr}_k(2,4)$ may be viewed as a projective algebraic
  set. In fact, all Grassmannians may be similarly realized as
  projective algebraic sets (in the sense of
  \Cref{exer:linrep:homog-ideal-vanishing-well-defined-proj-set})
  via the corresponding Pl\"ucker embeddings.

  Prove that $\mathrm{Gr}_k(2,4)$ may be covered with six copies of
  $\mathbf{A}^4_k$. (Recall from
  \Cref{exer:fields:affine-to-projective-space-covering} that
  $\mathbf{P}^5_k$ may be covered with six copies of $\mathbf{A}^5_k$;
  prove that the intersection of each with $\mathrm{Gr}_k(2,4)$ may be
  identified with $\mathbf{A}^4_k$ in a natural way.)
\end{exer}

\begin{exer}{}{exer:linrep:wedge2-sym2-as-kernels-in-t2}
  Assume $2$ is a unit in $R$, and let $F$ be a free $R$-module of
  finite rank.
  \begin{itemize}
  \item Define a function
    $\lambda : \bigwedge^2_R(F) \to \mathrm{T}^2_R(F)$ on a basis
    $e_i \wedge e_j$, $i < j$, by setting
    $\lambda(e_i \wedge e_j) = \frac{1}{2}(e_i \otimes e_j - e_j
    \otimes e_i)$ and extending by linearity. Prove that $\lambda$ is
    an injective $R$-module homomorphism and
    $\lambda(f_1 \wedge f_2) = \frac{1}{2}(f_1 \otimes f_2 - f_2
    \otimes f_1)$ for all $f_1, f_2 \in F$.
  \item Define a function
    $\sigma : \mathrm{S}^2_R(F) \to \mathrm{T}^2_R(F)$ on a basis
    $e_i \otimes e_j$, $i \leq j$, by setting
    $\sigma(e_i \otimes e_j) = \frac{1}{2}(e_i \otimes e_j + e_j
    \otimes e_i)$ and extending by linearity. Prove that $\sigma$ is
    an injective $R$-module homomorphism and
    $\sigma(f_1 \otimes f_2) = \frac{1}{2}(f_1 \otimes f_2 + f_2
    \otimes f_1)$ for all $f_1, f_2 \in F$.
  \item Prove that $\lambda$ identifies $\bigwedge^2_R(F)$ with the
    kernel of the map $\mathrm{T}^2_R(F) \to \mathrm{S}^2_R(F)$, and
    $\sigma$ identifies $\mathrm{S}^2_R(F)$ with the kernel of the map
    $\mathrm{T}^2_R(F) \to \bigwedge^2_R(F)$.
  \item In particular, there is a `meaningful' isomorphism
    $F \otimes_R F \cong \bigwedge^2_R(F) \oplus \mathrm{S}^2_R(F)$.
  \end{itemize}
\end{exer}
\begin{exer}{$*$}{exer:linrep:graded-algebra-lemma-equivalence}
  Prove the equivalence (ii)$\iff$(iii) in
  \Cref{lem:linrep:homogeneous-ideal-tfae}. [\Cref{subsec:linrep:very-small-detour-graded-algebra}]
\end{exer}

\begin{exer}{$\lnot$}{exer:linrep:homog-ideal-vanishing-well-defined-proj-set}
  Let $I \subseteq k[x_0, \dots, x_n]$ be a homogeneous ideal. Prove
  that the condition `$F(c_0, \dots, c_n) = 0$ for all $F \in I$' for
  a point $(c_0 : \cdots : c_n) \in \mathbf{P}^n_k$ is well-defined:
  it does not depend on the representative $(c_0, \dots, c_n)$ chosen
  for the point $(c_0 : \cdots : c_n)$. (Cf.\
  \Cref{exer:fields:homog-vanishing-well-defined-covering}, but note
  that not all $F$ in $I$ are homogeneous.)

  Thus, every homogeneous ideal $I \subseteq k[x_0, \dots, x_n]$
  determines a `projective algebraic set'
  \[
    \mathbf{V}(I) := \{(c_0 : \cdots : c_n) \in \mathbf{P}^n_k \mid
    (\forall F \in I),\, F(c_0, \dots, c_n) = 0\}.
  \]
  Note that $\mathbf{V}((x_0, \dots, x_n)) = \emptyset$. The ideal
  $(x_0, \dots, x_n) = \bigoplus_{i > 0} k[x_0, \dots, x_n]_i$ is
  irreverently called the \emph{irrelevant
    ideal}. [\ref{exer:linrep:grassmannian-2-4-as-projective-quadric},
  \ref{exer:linrep:weak-homogeneous-nullstellensatz}]
\end{exer}

\begin{exer}{}{exer:linrep:weak-homogeneous-nullstellensatz}
  (Cf.\
  \Cref{exer:linrep:homog-ideal-vanishing-well-defined-proj-set}.)
  Prove the `weak homogeneous Nullstellensatz': if $k$ is
  algebraically closed and $I \subseteq k[x_0, \dots, x_n]$ is a
  homogeneous ideal, then $\mathbf{V}(I) = \emptyset$ if and only if
  $I$ is either $k[x_0, \dots, x_n]$ or the irrelevant ideal
  $(x_0, \dots, x_n)$.  (Translate this into a question about
  $\mathbf{A}^{n+1}_k$.)
\end{exer}

\begin{exer}{}{exer:linrep:weyl-algebra-presentation}
  Let $k$ be a field of characteristic zero. A \emph{differential
    operator in one variable, with polynomial coefficients}, is a
  linear combination
  \begin{equation}
    \tag{$*$}
    a_0(x) + a_1(x)\,\partial_x + \dots + a_r(x)\,\partial_x^r
  \end{equation}
  where $a_i(x) \in k[x]$. Here $x$ acts on a polynomial
  $f(x) \in k[x]$ by multiplication by $x$, while $\partial_x$ acts by
  taking a formal derivative (as in
  \Cref{subsec:fields:separable-polynomials}). With the evident
  operations, differential operators form a ring, called the
  \emph{(first) Weyl algebra}. Note that the Weyl algebra is
  noncommutative: for example, $(x\partial_x)f(x) = xf'(x)$, while
  $(\partial_x x)f(x) = \partial_x(xf(x)) = f(x) + xf'(x)$.  Therefore
  $(\partial_x x - x\partial_x)f(x) = f(x)$, or put otherwise
  \begin{equation}
    \tag{$**$}
    \partial_x x - x\partial_x = 1
  \end{equation}
  in the Weyl algebra. Prove that the Weyl algebra is isomorphic to
  the quotient $\mathrm{T}^*_k(\langle x, y \rangle)/(yx - xy -
  1)$. (Use ($**$) to write any element of the Weyl algebra in the
  form given in ($*$). Show that this representation is unique, and
  deduce that ($**$) generates all relations among $x$ and
  $\partial_x$. Then use the first isomorphism theorem.)

  The relation ($**$) expresses (up to a factor) the basic fact that
  in quantum mechanics the position and momentum operators do not
  commute.  The Weyl algebra was introduced to study this
  phenomenon. Modules over rings of differential operators are called
  $D$-modules.
\end{exer}
\begin{exer}{}{exer:linrep:sym-power-of-free-module-is-free}
  Let $F$ be a free $R$-module of rank $r$. Prove that
  $\mathrm{S}^\ell_R(F)$ is free, and compute its rank.
\end{exer}

\begin{exer}{}{exer:linrep:sym-algebra-of-direct-sum}
  Let $F_1$, $F_2$ be free $R$-modules of finite rank. Prove that
  $\mathrm{S}^*_R(F_1 \oplus F_2) \cong \mathrm{S}^*_R(F_1) \otimes_R
  \mathrm{S}^*_R(F_2)$.
\end{exer}

\begin{exer}{$*$}{exer:linrep:skew-commutativity-of-exterior-algebra}
  Verify the skew-commutativity of the exterior algebra, stated in
  \Cref{rem:linrep:exterior-algebra-skew-commutative}. [\Cref{subsec:linrep:tensor-algebras}]
\end{exer}

\begin{exer}{}{exer:linrep:dim-of-exterior-algebra-is-2r}
  Let $V$ be a $k$-vector space of dimension $r$. Prove that, as a
  vector space, the exterior algebra $\bigwedge^*_k(V)$ has dimension
  $2^r$.
\end{exer}
\begin{exer}{$\lnot$}{exer:linrep:rees-algebra-is-noetherian}
  Prove that the Rees algebra $\mathrm{Rees}_R(I)$ of an ideal $I$ is
  Noetherian if $R$ is
  Noetherian. [\ref{exer:linrep:artin-rees-lemma}]
\end{exer}

\begin{exer}{$\lnot$}{exer:linrep:artin-rees-lemma}
  Let $R$ be a Noetherian ring, $I$ an ideal of $R$, and consider the
  Rees algebra $\mathrm{Rees}_R(I) = \bigoplus_{\ell \geq 0} I^\ell$.
  By \Cref{exer:linrep:rees-algebra-is-noetherian},
  $\mathrm{Rees}_R(I)$ is Noetherian.
  \begin{itemize}
  \item For $\ell \geq 0$, let $J_\ell \subseteq I^\ell$ be ideals of
    $R$, and view $J := \bigoplus_{\ell \geq 0} J_\ell$ as a
    sub-$R$-module of $\mathrm{Rees}_R(I)$. Prove that $J$ is an ideal
    of $\mathrm{Rees}_R(I)$ if and only if
    $I^n J_\ell \subseteq J_{\ell+n}$ for all $\ell, n \geq 0$.
  \item Assume $J := \bigoplus_{\ell \geq 0} J_\ell$ is an ideal of
    $\mathrm{Rees}_R(I)$. Prove that $J$ admits a finite set of
    homogeneous generators.
  \item Choose a finite set of homogeneous generators for $J$, and let
    $s$ be the largest degree of an element in this set. Prove that
    $J_{s+1} \subseteq I^{s+1} J_0 + I^s J_1 + \dots + I J_s$ (as
    ideals of $R$).
  \item Prove that $J_{s+\ell} = I^\ell J_s$ for all $\ell \geq 0$.
    This is a particular case of the `Artin-Rees'
    lemma. [\ref{exer:linrep:intersection-of-ideal-powers-via-artin-rees}]
  \end{itemize}
\end{exer}

\begin{exer}{$\lnot$}{exer:linrep:intersection-of-ideal-powers-via-artin-rees}
  Let $R$ be a Noetherian ring, and let $I$ be an ideal of $R$. Prove
  that $I \cdot \bigcap_{n \geq 0} I^n = \bigcap_{n \geq 0} I^n$.
  (Hint: \Cref{exer:linrep:artin-rees-lemma}.)
  [\ref{exer:linrep:intersection-of-ideal-powers-krull}]
\end{exer}

\begin{exer}{}{exer:linrep:rees-algebra-example-iso-to-quotient}
  Let $k$ be a field, and consider the ideal $I = (x, y)$ in the ring
  $R = k[x, y]$. Prove that the Rees algebra
  $\bigoplus_{\ell \geq 0} I^\ell$ is isomorphic to the quotient of a
  polynomial ring $k[x, y, s, t]$ by the ideal $(tx - sy)$. Deduce
  that, in this case, the Rees algebra of $I$ is isomorphic to the
  symmetric algebra $\mathrm{S}^*_{k[x,y]}(I)$. (For the last point,
  use the exact sequence mentioned at the end of
  \Cref{subsec:linrep:tensor-algebras}. It will be helpful to have
  an explicit presentation of $I$ as a $k[x,y]$-module; for this,
  looking back at \Cref{exer:linalg:free-resolution-of-z-over-zxy} may
  help.) [\Cref{subsec:linrep:very-small-detour-graded-algebra},
  \Cref{subsec:linrep:tensor-algebras}]
\end{exer}

\begin{exer}{$\lnot$}{exer:linrep:koszul-complex-definition}
  Let $\mathbf{a} = a_1, \dots, a_n$ be a list of elements of $R$, and
  let $F \cong R^n$. Rename $\bigwedge^r_R(F)$ by $K_r(\mathbf{a})$.
  Define $R$-module homomorphisms
  $d_r : K_r(\mathbf{a}) \to K_{r-1}(\mathbf{a})$ on bases by setting
  \[
    d_r(e_{i_1} \wedge \dots \wedge e_{i_r}) = \sum_{j=1}^r (-1)^{j-1}
    a_{i_j}\, e_{i_1} \wedge \dots \wedge \hat{e}_{i_j} \wedge \dots
    \wedge e_{i_r},
  \]
  where the hat denotes that the hatted element is omitted.
  \begin{itemize}
  \item Prove that $d_{r-1} \circ d_r = 0$.
  \end{itemize}
  Thus, a collection $a_1, \dots, a_n$ of elements of $R$ determines a
  complex of $R$-modules
  \[
    0 \to K_n(\mathbf{a}) \xrightarrow{d_n} \cdots \xrightarrow{d_2}
    K_1(\mathbf{a}) \xrightarrow{d_1} K_0(\mathbf{a}) = R \to R/I \to
    0,
  \]
  where $I = (a_1, \dots, a_n)$. This is called the \emph{Koszul
    complex} of $\mathbf{a}$.
  \begin{itemize}
  \item Check that the complexes constructed in
    \Cref{exer:linalg:koszul-complex-length-2,exer:linalg:koszul-complex-length-3}
    are Koszul complexes.
  \end{itemize}
  As proven for $n = 2, 3$ in
  \Cref{exer:linalg:koszul-complex-length-2,exer:linalg:koszul-complex-length-3},
  the
  Koszul complex is exact if $(a_1, \dots, a_n)$ is a regular sequence
  in $R$, providing a free resolution for $R/I$ in that case. Try to
  prove this in general. [\ref{exer:linalg:koszul-complex-length-3}]
\end{exer}

\section{Hom and duals}
\label{sec:linrep:hom-and-duals}

Our rapid overview of tensor products has occasionally brought us into
close proximity with the Hom functors, and I will close this chapter
by devoting some attention to these functors. As in the case of
tensors, we will be preoccupied with adjunction and exactness
properties, since concrete computations depend heavily on these
properties. We will deal with these properties more carefully in the
next (and last) section; in this section we will concentrate on one
important special case of Hom, that is, the duality functor.

As in the rest of the chapter, $R$ denotes a fixed commutative ring.

\subsection{Adjunction again}
\label{subsec:linrep:adjunction-again}
The careful reader will agree that almost everything we know about the
tensor product follows from the fact that it is a left-adjoint
functor.  This fact is spelled out in \Cref{lem:linrep:hom-tensor-adjunction}, whose
flip side gives us just as much information concerning the covariant
flavor of the Hom functor: $L \mapsto \Hom_{R\text{-Mod}}(N, L)$.
Explicitly, the following is an exact restatement of
\Cref{coro:linrep:tensor-hom-adjunction}:

\begin{coro}{}{coro:linrep:hom-tensor-adjunction-other-side}
  For every $R$-module $N$, the functor $\Hom_R(N, -)$ is
  right-adjoint to the functor $- \otimes_R N$.
\end{coro}

What about the `contravariant' flavor of Hom:
$M \mapsto h_N(M) := \Hom_{R\text{-Mod}}(M, N)$ (this functor was
discussed in general terms in
\Cref{subsec:linrep:examples-of-functors})?

\begin{prop}{}{prop:linrep:hom-contravariant-self-adjoint}
  For every $R$-module $N$, the functor $\Hom_R(-, N)$ is
  right-adjoint to itself.
\end{prop}

This statement should be parsed carefully, because $\Hom_R(-, N)$ is a
contravariant functor. The statement of \Cref{prop:linrep:hom-contravariant-self-adjoint}
may lead us astray into thinking that $\Hom_R(-, N)$ must also be its
own left-adjoint, and this is not the case (indeed this would make it
a right-exact functor, and we will soon see that $\Hom_R(-, N)$ is not
right-exact in general).

The point is that, by definition of contravariant functor,
$h_N = \Hom_R(-, N)$ should be viewed as a covariant functor from the
opposite category:
\[
  h_N : R\text{-Mod}^{\mathrm{op}} \to R\text{-Mod};
\]
it can also be viewed as a covariant functor to the opposite category,
\[
  h_N^{\mathrm{op}} : R\text{-Mod} \to R\text{-Mod}^{\mathrm{op}},
\]
by simply reversing arrows after the fact rather than before. A more
precise statement of \Cref{prop:linrep:hom-contravariant-self-adjoint} is that $h_N$ is
right-adjoint to $h_N^{\mathrm{op}}$.

\begin{proof}
  Let $L$, $M$, $N$ denote $R$-modules. Recall (cf.\ the
  considerations preceding \Cref{lem:linrep:hom-tensor-adjunction}) that $R$-bilinear
  maps $\phi : L \times M \to N$ may be identified with $R$-linear
  maps $L \to \Hom_R(M, N)$. By the same token, they may be identified
  with $R$-linear maps $M \to \Hom_R(L, N)$: for fixed $m \in M$, the
  function $\phi(-, m)$ is an $R$-linear map $L \to N$. Tracing these
  two identifications gives a canonical bijection
  \begin{equation}
    \tag{$*$}
    \Hom_R(L, \Hom_R(M, N)) \cong
    \Hom_R(M, \Hom_R(L, N)).
  \end{equation}
  We may view $\Hom_R$ just as well as
  $\Hom_{R\text{-Mod}^{\mathrm{op}}}$, provided that we reverse
  arrows; thus, the canonical bijection in $(*)$ may be rewritten as
  \[
    \Hom_{R\text{-Mod}}(L, h_N(M)) \cong
    \Hom_{R\text{-Mod}^{\mathrm{op}}}(h_N^{\mathrm{op}}(L), M).
  \]
  This says that $h_N$ is right-adjoint to $h_N^{\mathrm{op}}$ (the
  naturality requirement is also immediate, and as usual it is left to
  the reader).
\end{proof}

As both $\Hom_R(N, -)$ (by \Cref{coro:linrep:hom-tensor-adjunction-other-side}) and
$\Hom_R(-, N)$ (by \Cref{prop:linrep:hom-contravariant-self-adjoint}) are right-adjoint
functors, \Cref{lem:linrep:right-adjoints-commute-with-limits} tells us they commute with
limits. This fact must also be parsed carefully, because of the
contravariant nature of $\Hom_R(-, N)$: for example, products are
limits in $R$-Mod, but direct sums are colimits in $R$-Mod, hence
limits in $R\text{-Mod}^{\mathrm{op}}$. Thus:

\begin{coro}{}{coro:linrep:hom-commutes-with-products}
  For every $R$-module $N$ and every family $\{M_i\}_{i \in I}$ of
  $R$-modules,
  \[
    \Hom_R\!\Bigl(N, \prod_{i \in I} M_i\Bigr) \cong \prod_{i \in I}
    \Hom_R(N, M_i),
  \]
  \[
    \Hom_R\!\Bigl(\bigoplus_{i \in I} M_i, N\Bigr) \cong \prod_{i \in
      I} \Hom_R(M_i, N).
  \]
\end{coro}

Also, as pointed out in \Cref{prop:linrep:right-adjoint-additive-functors-left-exact}, the fact that both
the covariant and contravariant flavors of Hom are right-adjoints has
immediate implications for their exactness, adding yet another
Pavlovian statement to the list: $\Hom_R(M, -)$ and $\Hom_R(-, N)$ are
left-exact functors.

By the contravariant nature of $\Hom_R(-, N)$, the left-exactness of
the latter means that if
\[
  A \to B \to C \to 0
\]
is an exact sequence of $R$-modules, then the induced sequence
\[
  0 \to \Hom_R(C, N) \to \Hom_R(B, N) \to \Hom_R(A, N)
\]
is also exact. The diligent reader has verified this already `by hand'
in the distant past
(\Cref{exer:rings:hom-contravariant-left-exact-from-ses}), without
appealing to adjunction. Direct proofs of the left-exactness of both
Hom functors are straightforward and instructive.

\subsection{Dual modules}
\label{subsec:linrep:dual-modules}
We will explore further the exactness (and lack of exactness) of Hom
in
\Cref{sec:linrep:projective-and-injective-modules-and-the-ext-functors}. Before
doing that, however, I want to examine one important particular case
of the contravariant aspect of the Hom functor.

\begin{defn}{}{defn:linrep:dual-module}
  Let $M$ be an $R$-module. The \emph{dual} $M^\vee$ of $M$ is the
  $R$-module $\Hom_R(M, R)$.
\end{defn}

One use of the dual module is to translate Hom computations into
$\otimes$ computations, at least in special cases.

\begin{prop}{}{prop:linrep:hom-into-free-module-via-dual}
  Let $M$ be any $R$-module, and let $F$ be a free $R$-module of
  finite rank. Then $\Hom_R(M, F) \cong M^\vee \otimes_R F$.
\end{prop}
\begin{proof}
  By hypothesis $F \cong R^{\oplus n} = R^n$; hence
  \[
    \Hom_R(M, F) \cong \Hom_R(M, R^n) \cong \Hom_R(M, R)^n \cong
    \Hom_R(M, R) \otimes_R R^n \cong M^\vee \otimes_R F,
  \]
  by \Cref{coro:linrep:hom-commutes-with-products}.
\end{proof}

In fact, note that for all $R$-modules $M$, $N$ there is a natural
`evaluation' map
\[
  \epsilon : M^\vee \otimes_R N \to \Hom_R(M, N)
\]
defined on pure tensors by mapping $f \otimes n$ (for
$f \in \Hom_R(M, R)$ and $n \in N$) to the $R$-module homomorphism
$M \to N$ given by $m \mapsto f(m)n$. The reader will check (Exercise
5.5) that $\epsilon$ is an isomorphism if $N$ is free of finite rank;
this gives a more precise version of \Cref{prop:linrep:hom-into-free-module-via-dual}.

\subsection{Duals of free modules}
\label{subsec:linrep:duals-of-free-modules}
By definition, the `duality' functor $M \mapsto M^\vee$ is a
particular case of the contravariant flavor of Hom; hence it is itself
contravariant and commutes with limits. Here is a first, immediate
consequence:

\begin{lem}{}{lem:linrep:dual-of-direct-sum}
  For every family $\{M_i\}$ of $R$-modules,
  $\bigl(\bigoplus_i M_i\bigr)^\vee \cong \prod_i M_i^\vee$.
\end{lem}

Specializing to the case in which $M_i = R$ for all $i$ determines the
dual of a free module:

\begin{coro}{}{coro:linrep:dual-of-free-module}
  The dual of a free module is isomorphic to a product of copies of
  $R$: $(R^{\oplus S})^\vee \cong R^S$. In particular,
  $(R^n)^\vee \cong R^n$: if $F$ is a free $R$-module of finite rank,
  then $F^\vee \cong F$.
\end{coro}
\begin{proof}
  This follows from \Cref{lem:linrep:dual-of-direct-sum} and the fact that
  $R^\vee = \Hom_R(R, R)$ is isomorphic to $R$.
\end{proof}

Carefully note the magically disappearing $\oplus$ from the left-hand
side to the right-hand side in the statement of
\Cref{coro:linrep:dual-of-free-module}: direct sums become direct products through
the contravariant Hom (\Cref{coro:linrep:hom-commutes-with-products}).  However, a direct
product of finitely many modules happens to be isomorphic to their
direct sum (as we have known for a long time:
\Cref{prop:rings:direct-sum-product-coproduct}), hence the second part
of the statement.

\begin{rem}{}{rem:linrep:infinite-product-vs-direct-sum}
  If $S$ is infinite, $R^S$ is `much larger' than $R^{\oplus S}$: the
  first module consists of all functions $S \to R$, while the second
  only retains those that are zero for all but finitely many
  $s \in S$.
\end{rem}

\begin{rem}{}{rem:linrep:basis-set-not-determined-by-module}
  Note that the set $S$ is not determined by the isomorphism class of
  a free module $F = R^{\oplus S}$. We proved in \Cref{coro:linalg:free-module-iso-iff-bijection-of-sets}
  (if $R$ is an integral domain) that the cardinality of $S$ is
  determined by the isomorphism class of $F$, but of course $S$
  itself, say as a subset of $F$, is not. In fact, recall
  (\Cref{subsec:linalg:linear-independence-and-bases}) that the
  choice of a specific isomorphism $F \cong R^{\oplus S}$ is
  equivalent to the choice of a basis of $F$. Now, the isomorphism
  appearing in \Cref{coro:linrep:dual-of-free-module} requires knowledge of $S$;
  indeed, we will verify in a moment (\Cref{exam:linrep:dual-basis-depends-on-basis}) that
  this isomorphism, even in the finite case, does depend on the choice
  of a basis. It is not canonical! The reader should endeavor to
  remember this slogan: a finite-rank free module---for example, a
  finite-dimensional vector space---is isomorphic to its dual, but not
  canonically.
\end{rem}

This remark can be clarified by means of the following notion.
\begin{defn}{}{defn:linrep:dual-basis}
  Consider the standard basis $(e_1, \dots, e_n)$ of $R^n$. The
  \emph{dual basis} of $(R^n)^\vee \cong R^n$ consists of
  $(\check{e}_1, \dots, \check{e}_n)$, where
  $\check{e}_i \in (R^n)^\vee = \Hom_R(R^n, R)$ is determined by
  \[
    \check{e}_i(e_j) = \begin{cases} 1 & \text{if } j = i, \\
      0 & \text{otherwise.} \end{cases}
  \]
\end{defn}

The isomorphism $(R^n)^\vee \cong R^n$ of \Cref{coro:linrep:dual-of-free-module}
is obtained precisely by matching $e_i$ and $\check{e}_i$: this will
be crystal clear to the reader who works out
\Cref{exer:linrep:direct-computation-dual-of-free-module}. In
particular, the vectors $\check{e}_i$ do form a basis of $(R^n)^\vee$.

\begin{exam}{}{exam:linrep:dual-basis-depends-on-basis}
  To see that these isomorphisms do depend on the choice of the basis,
  consider the standard basis $(e_1, e_2)$ of $R^2$ and the
  corresponding dual basis $(\check{e}_1, \check{e}_2)$ of
  $(R^2)^\vee$, and then choose a different basis $(e_1', e_2')$ for
  $R^2$, where $e_1' = e_1$ and $e_2' = e_1 + e_2$, and let
  $(\check{e}_1', \check{e}_2')$ be the corresponding dual basis. By
  definition $\check{e}_1'(e_2') = 0$, while
  $\check{e}_1(e_2') = \check{e}_1(e_1 + e_2) = 1 + 0 = 1$.
  Therefore, $\check{e}_1' \neq \check{e}_1$ even if $e_1' = e_1$: the
  two isomorphisms $R^2 \cong (R^2)^\vee$ determined by the two bases
  are different.
\end{exam}

The fact that duality leads to noncanonical isomorphisms is somewhat
unpleasant, but something magical will happen soon
(\Cref{subsec:linrep:duals-and-matrices-biduality}): we will
recover a canonical---that is, independent of any choice---isomorphism
(on finite-rank free modules) by applying the duality functor twice.
Whatever twist is introduced by the duality functor will be untwisted
by applying duality again.

\subsection{Duality and exactness}
\label{subsec:linrep:duality-and-exactness}
Before seeing this, let us contemplate another immediate consequence
of the fact that duality is a particular case of the contravariant
Hom:

\begin{lem}{}{lem:linrep:duality-functor-left-exact}
  The duality functor is left-exact: every exact sequence
  \[
    L \to M \to N \to 0
  \]
  of $R$-modules induces an exact sequence
  \[
    0 \to N^\vee \to M^\vee \to L^\vee.
  \]
\end{lem}
\begin{proof}
  This is an immediate consequence of the left-exactness of Hom.
\end{proof}

The duality functor is not exact in general: for example, taking duals
in the exact sequence of abelian groups (a.k.a.\ $\Z$-modules)
\[
  0 \to \Z \xrightarrow{\cdot 2} \Z \to \Z/2\Z \to 0
\]
gives the sequence
\begin{equation}
  \tag{$*$}
  0 \to \Z^\vee \xrightarrow{\gamma} \Z^\vee \to 0
\end{equation}
since the dual of $\Z/2\Z$ is zero
(\Cref{exer:linrep:z2z-dual-is-zero}). The map $\gamma$ in this
sequence is the dual of the multiplication by $2$.  If $f : \Z \to \Z$
is an element of $\Z^\vee$, then $\gamma(f)$ is obtained by
composition:
\[
  \begin{tikzcd}
    \Z \ar[r, "\gamma(f)"] \ar[dr, "\cdot 2"'] &
    \Z \\
    & \Z \ar[u, "f"']
  \end{tikzcd}
\]
By linearity, $\gamma(f)(n) = f(2n) = 2f(n)$ is even for all
$n \in \Z$. In particular, the identity $\Z \to \Z$ (as an element of
$\Z^\vee$) is not in the image of $\gamma$. Thus $\gamma$ is not
surjective, and the sequence $(*)$ is not exact.

There are situations in which duality is exact; we will understand
this better after digesting
\Cref{sec:linrep:projective-and-injective-modules-and-the-ext-functors},
but the following observation will suffice for now.
\begin{prop}{}{prop:linrep:dual-sequence-exact-when-quotient-free}
  Let
  \[
    0 \to M \xrightarrow{\mu} N \xrightarrow{\nu} P \to 0
  \]
  be an exact sequence of $R$-modules, with $P$ free. Then the induced
  sequence
  \[
    0 \to P^\vee \xrightarrow{\nu^\vee} N^\vee \xrightarrow{\mu^\vee}
    M^\vee \to 0
  \]
  is exact.
\end{prop}

This is also not new to the diligent readers, as it is a particular
case of \Cref{exer:rings:exact-seq-with-free-quotient-splits} and,
further, it is a consequence of
\Cref{exer:linrep:flat-in-exact-sequence-with-flat-quotient}. But here is a direct
argument:

\begin{proof}
  \Cref{lem:linrep:duality-functor-left-exact} takes care of all but the surjectivity of
  the map $N^\vee \to M^\vee$ induced from $M \to N$:
  \[
    \begin{tikzcd}
      M \ar[r, "\mu"] \ar[dr, "f"'] &
      N \ar[d, dashed, "\exists?\, g"] \\
      & R
    \end{tikzcd}
  \]
  The question is whether every $R$-linear $f : M \to R$ can be
  extended to an $R$-linear map $g : N \to R$ so that
  $f = g \circ \mu$, that is, $f = \mu^\vee(g)$.

  As $P$ is free, $P \cong R^{\oplus S}$ for some set $S$. Choosing
  (arbitrarily!) preimages in $N$ of the standard basis vectors
  $e_s \in R^{\oplus S}$ gives a set-function $S \to N$, extending to
  an $R$-linear map $\rho : P \to N$ by the universal property of free
  modules:
  \[
    \begin{tikzcd}
      0 \ar[r] & M \ar[r, "\mu"] & N \ar[r, "\nu"] & P \ar[r] \ar[l,
      bend right, "\rho"'] & 0
    \end{tikzcd}
  \]
  By construction, $\nu \circ \rho = \id_P$. Now let $n \in N$. Then
  $\nu(\rho(\nu(n))) = \nu(n)$, giving $\nu(n - \rho(\nu(n))) = 0$. By
  the exactness of the given sequence, $\exists m \in M$ such that
  $n - \rho(\nu(n)) = \mu(m)$. The element $m$ is unique, since $\mu$
  is injective. Setting $g(n) := f(m)$ gives the necessary
  extension. Indeed, $g$ is immediately checked to be $R$-linear, and
  $g(\mu(m)) = f(m)$ by definition.
\end{proof}

\begin{coro}{}{coro:linrep:duality-exact-on-vector-spaces}
  The duality functor is exact on vector spaces.
\end{coro}

\subsection{Duals and matrices; biduality}
\label{subsec:linrep:duals-and-matrices-biduality}
Even without the benefit of full exactness, \Cref{lem:linrep:duality-functor-left-exact}
and \Cref{coro:linrep:dual-of-free-module} reduce (in principle) the computation
of the dual of any finitely presented module to matrix calculus. If
\[
  R^n \xrightarrow{\alpha} R^m \to M \to 0
\]
is a presentation of an $R$-module $M$, then the dual $M^\vee$ is
identified with the kernel of the dual of $\alpha$:
\[
  0 \to M^\vee \to R^m \xrightarrow{\alpha^\vee} R^n.
\]
Here I am applying isomorphisms $(R^m)^\vee \cong R^m$,
$(R^n)^\vee \cong R^n$ from \Cref{coro:linrep:dual-of-free-module}.

The map $\alpha^\vee$ is obtained by applying the functor
$\Hom_R(-,R)$ to the map $\alpha$. Once bases are chosen, $\alpha$ is
determined by an $m \times n$ matrix; likewise, $\alpha^\vee$
corresponds to an $n \times m$ matrix. What is the relation between
these two matrices? Of course there is only one sensible answer to
this question, and it is correct:
\begin{lem}{}{lem:linrep:dual-map-is-transpose}
  Let $A$ be the matrix representing a linear map
  $\alpha : R^n \to R^m$ with respect to the standard bases. Then the
  dual map $\alpha^\vee : (R^m)^\vee \to (R^n)^\vee$ is represented by
  the transpose of $A$ with respect to the corresponding dual bases.
\end{lem}

The (easy) verification of this fact is left to the reader (Exercise
5.10). The upshot is that the dual of the cokernel of a matrix $A$ is
the kernel of the transpose of $A$.

Left-exactness of duality implies restrictions on which modules may
appear as duals of other modules:
\begin{prop}{}{prop:linrep:dual-module-torsion-free}
  Let $R$ be an integral domain, and let $M$ be an $R$-module. Then
  $M^\vee$ is torsion-free.
\end{prop}

\begin{proof}
  There is a surjection $R^{\oplus S} \twoheadrightarrow M$, thus an
  exact sequence $R^{\oplus T} \to R^{\oplus S} \to M \to 0$.
  Dualizing, $M^\vee$ is realized as the kernel of the induced map
  $R^S \to R^T$; hence $M^\vee$ may be identified with a submodule of
  a product $R^S$. The latter has no nonzero torsion, so neither does
  $M^\vee$.
\end{proof}

The diligent reader has proved this already in
\Cref{exer:linalg:hom-torsion-free-if-target-torsion-free} and
probably by a more direct way than what we did just now, but never
mind. In any case, remember that dualizing kills torsion (over
integral domains). Practice this by working out
\Cref{exer:linrep:compute-dual-of-fin-gen-module-over-pid}.  Even if
not every module is the dual of a module, we may wonder whether an
arbitrarily given module $M$ has some relation with a suitably chosen
dual, and this is indeed the case: there is a canonical map to the
\emph{double-dual},
\[
  \omega : M \to M^{\vee\vee}.
\]
To define this, consider the $R$-bilinear map $M^\vee \times M \to R$
sending $(f, m)$ to $f(m)$. For every $m \in M$, we get\footnote{The
  reader with a categorical frame of mind can view this a little more
  formally by working out
  \Cref{exer:linrep:canonical-bijection-hom-dual-swap}.} an $R$-linear
map $M^\vee \to R$, $f \mapsto f(m)$, in other words, an element
$\omega(m)$ of $M^{\vee\vee}$. It is immediately checked that $\omega$
is $R$-linear.

By \Cref{prop:linrep:dual-module-torsion-free}, $M^{\vee\vee}$ is
torsion-free over integral domains. In
fact, the double-dual construction is a standard way to `clean up' a
module, removing its torsion.  The map $\omega$ is interesting even
when torsion is not an issue. As we have seen in
\Cref{coro:linrep:dual-of-free-module}, if $F$ is a finite-rank free $R$-module,
then a choice of basis for $F$ determines an isomorphism
$F \cong F^\vee$; but this does indeed depend on the basis, so it is
not a canonical isomorphism. Doing it twice,
$F \cong F^\vee \cong F^{\vee\vee}$, may seem a priori bound to make
the situation even worse --- surely, if we compose two noncanonical
maps, then we cannot expect to get something canonical, right? Well,
we do in this case, since this composition is nothing but the map
$\omega$
(\Cref{exer:linrep:composition-of-isos-is-canonical-biduality-map})
and $\omega$ knows nothing about choices of bases or other such
ambiguities. Therefore, our slogan admits the following addendum:
finite-rank free modules are noncanonically isomorphic to their duals
and canonically isomorphic to their double-duals.  In terms of
matrices (and keeping in mind that duality is exact on free modules)
this is perhaps not too surprising. After all, taking the transpose of
a matrix twice simply gives the matrix back; thus, even if making
sense of the homomorphism corresponding to the transpose of a matrix
may depend on the choice of a basis, surely such a choice is
inessential if the transpose is taken twice.

Incidentally, a module $M$ is \emph{reflexive} if the corresponding
map $\omega$ is an isomorphism. Thus, reflexive modules are
necessarily torsion-free and finite-rank free modules are reflexive.

\subsection{Duality on vector spaces}
\label{subsec:linrep:duality-on-vector-spaces}
In all these considerations, the case of vector spaces over a field
$k$ is only special in that vector spaces are free. Collecting the
foregoing observations (including
\Cref{exer:linrep:hom-from-free-iso-dual-tensor-n}) for this
particular case, we have the following:
\begin{itemize}
\item Duality is a contravariant, exact functor on $k\text{-Vect}$.
\item If $V, W$ are vector spaces, then
  $(V \oplus W)^\vee \cong V^\vee \oplus W^\vee$.
\item If $V, W$ are vector spaces and $\dim V < \infty$ or
  $\dim W < \infty$, then $\Hom_k(V,W)$ is isomorphic to
  $V^\vee \otimes_k W$.
\item If $\dim_k V < \infty$, a choice of a basis on $V$ determines an
  isomorphism $V \cong V^\vee$.
\item Finite-dimensional vector spaces are reflexive: if
  $\dim_k V < \infty$, then the canonical map
  $\omega : V \to V^{\vee\vee}$ is an isomorphism.
\end{itemize}

As usual, $R$ is a fixed commutative ring unless stated otherwise.

\subsection*{Exercises}
\begin{exer}{$*$}{exer:linrep:hom-from-free-iso-dual-tensor-n}
  Prove that if $F$ is a free $R$-module of finite rank and $N$ is any
  $R$-module, then $\Hom_R(F, N) \cong F^\vee \otimes_R N$.
  [\Cref{subsec:linrep:duality-on-vector-spaces}]
\end{exer}

\begin{exer}{$\lnot$}{exer:linrep:epi-mono-characterization-via-hom}
  Let $\alpha : A \to B$ be a homomorphism of $R$-modules. Prove that
  $\alpha$ is an epimorphism if the induced map\footnote{Note that the
    `set-theoretic' quality of $\alpha^*$ is all that is needed here;
    in this problem, the $R$-module structure of $\Hom_R$ is
    irrelevant. Thus, the result also holds for noncommutative $R$.}
  $\alpha^* : \Hom_R(B, N) \to \Hom_R(A, N)$ is injective for all
  $R$-modules $N$, and $\alpha$ is a monomorphism if $\alpha^*$ is
  surjective for all $N$. Prove that the converse to the first
  statement holds and the converse to the second statement does not
  hold. However, show that if $\alpha$ admits a left-inverse, then
  $\alpha^*$ is surjective for all
  $N$. [\ref{exer:linrep:exactness-criterion-via-hom-dual},
  \ref{exer:linrep:injective-iff-hom-surjective-for-injectives}]
\end{exer}

\begin{exer}{}{exer:linrep:exactness-criterion-via-hom-dual}
  Prove that a sequence
  \[
    0 \to A \xrightarrow{\alpha} B \xrightarrow{\beta} C \to 0
  \]
  of $R$-modules is exact if the induced sequence
  \[
    \tag{$*$} 0 \to \Hom_R(C,N) \xrightarrow{\beta^*} \Hom_R(B,N)
    \xrightarrow{\alpha^*} \Hom_R(A,N) \to 0
  \]
  is exact for all $R$-modules $N$. (You have done most of this
  already, in \Cref{exer:linrep:epi-mono-characterization-via-hom}. To
  show $\ker\beta \subseteq \image\alpha$, choose
  $N = B/\image(\alpha)$.) Remember that the converse does not hold,
  since in general $\Hom_R(-,N)$ is not exact.  What extra hypothesis
  on $\alpha$ would guarantee the exactness of $(*)$ for all $N$?
\end{exer}

\begin{exer}{$\lnot$}{exer:linrep:hom-i-r-mod-i-restriction-map-zero}
  Let $I$ be an ideal of $R$. As $I^2 \subseteq I$, there is a natural
  restriction map $\Hom_R(I, R/I) \to \Hom_R(I^2, R/I)$. Prove that
  the image of this map is $0$.  Prove that
  $\Hom_R(I/I^2, R/I) \cong \Hom_R(I, R/I)$. (This module is important
  in algebraic geometry, as it carries the information of a `normal
  bundle' in good situations.)
  [\ref{exer:linrep:ext1-of-quotient-iso-to-hom-normal-bundle}]
\end{exer}

\begin{exer}{$*$}{exer:linrep:evaluation-map-dual-tensor-free-iso-hom}
  Prove that the evaluation map $M^\vee \otimes_R F \to \Hom_R(M,F)$
  is an isomorphism if $F$ is free of finite rank, providing an
  alternative proof of
  \Cref{prop:linrep:hom-into-free-module-via-dual}. [\Cref{subsec:linrep:dual-modules}]
\end{exer}

\begin{exer}{$*$}{exer:linrep:z2z-dual-is-zero}
  Show that $(\Z/2\Z)^\vee = \Hom_{\Z}(\Z/2\Z, \Z) = 0$.
  [\Cref{subsec:linrep:duality-and-exactness}]
\end{exer}

\begin{exer}{$*$}{exer:linrep:direct-computation-dual-of-free-module}
  Prove `directly' that $(R^{\oplus S})^\vee \cong R^S$: how does an
  $R$-linear map $R^{\oplus S} \to R$ determine a function $S \to R$,
  and what is the inverse of this correspondence?
  [\Cref{subsec:linrep:duals-of-free-modules}]
\end{exer}

\begin{exer}{}{exer:linrep:hom-m-mvee-equiv-bilinear-map}
  Prove that the datum of an $R$-linear map $M \to M^\vee$ is
  equivalent to the datum of an $R$-bilinear map $M \times M \to R$,
  and explain why this equivalence can be set up in two ways. If
  $F = R^n$ is a free $R$-module of finite rank, determine the
  bilinear map $F \times F \to R$ corresponding to the isomorphism
  $F \cong F^\vee$ given in \Cref{coro:linrep:dual-of-free-module}.
\end{exer}

\begin{exer}{}{exer:linrep:nondegenerate-nonsingular-bilinear-forms}
  An $R$-bilinear map $\phi : M \times M \to R$ is
  \emph{nondegenerate} if the induced maps $M \to M^\vee$ are
  injective, and it is \emph{nonsingular} if they are isomorphisms.
  The notions coincide if $M$ is a finite-dimensional vector space.
  Prove that the `standard inner product' in $R^n$ (defined in
  \Cref{exer:linalg:on-r-iff-orthonormal-columns-rows}) is
  nondegenerate. If $M$ is free of rank $n$, let $(e_1, \dots, e_n)$
  be a basis of $M$, and let $A = (a_{ij})$ be the matrix with entries
  $a_{ij} = \phi(e_i, e_j)$. Prove that $\phi$ is nondegenerate if and
  only if $\det(A)$ is nonzero, and it is nonsingular if and only if
  $\det(A)$ is a unit. (Cf.\ Proposition VI.6.5.)
\end{exer}

\begin{exer}{}{exer:linrep:dual-map-is-transpose-matrix}
  Prove \Cref{lem:linrep:dual-map-is-transpose}.
\end{exer}

\begin{exer}{$*$}{exer:linrep:compute-dual-of-fin-gen-module-over-pid}
  Let $M$ be a finitely generated module over a PID, of rank $r$.
  `Compute' the dual
  $M^\vee$. [\Cref{subsec:linrep:duals-and-matrices-biduality}]
\end{exer}

\begin{exer}{$*$}{exer:linrep:canonical-bijection-hom-dual-swap}
  Let $M$, $N$ be $R$-modules. Show that there is a canonical
  bijection $\Hom_R(N, M^\vee) \cong \Hom_R(M, N^\vee)$. Choosing
  $N = M^\vee$, the left-hand side has a distinguished element, namely
  the identity $M^\vee \to M^\vee$.  Prove that the corresponding
  element on the right is the map $\omega : M \to M^{\vee\vee}$
  defined in
  \Cref{subsec:linrep:duals-and-matrices-biduality}. [\Cref{subsec:linrep:duals-and-matrices-biduality}]
\end{exer}

\begin{exer}{$*$}{exer:linrep:composition-of-isos-is-canonical-biduality-map}
  Let $F \cong R^n$ be a finite-rank free $R$-module. Verify that the
  composition of the (noncanonical) isomorphisms
  $F \cong F^\vee \cong F^{\vee\vee}$ from \Cref{coro:linrep:dual-of-free-module}
  is the (canonical) isomorphism $\omega$ defined in
  \Cref{subsec:linrep:duals-and-matrices-biduality}. [\Cref{subsec:linrep:duals-and-matrices-biduality}]
\end{exer}

\begin{exer}{}{exer:linrep:canonical-map-to-bidual-injective-free}
  Let $F$ be a free $R$-module (of any rank). Prove that the canonical
  map $F \to F^{\vee\vee}$ is injective. What is the simplest example
  you know of a module $M$ such that $M \to M^{\vee\vee}$ is not
  injective?
\end{exer}

\begin{exer}{$\lnot$}{exer:linrep:annihilator-of-subspace-basics}
  Let $V$ be a vector space, and let $W \subseteq V$ be a subspace.
  The \emph{annihilator} of $W$ is\footnote{The notation is inspired
    by
    \Cref{exer:linalg:orthogonal-complement-is-subspace-dim-n-minus-1},
    which is a particular case of this problem. Do you see why?}
  \[
    W^\perp := \{\,\check{v} \in V^\vee \mid (\forall w \in W),\,
    \check{v}(w) = 0\,\}.
  \]
  Prove that $W^\perp$ is a subspace of $V^\vee$. If $\dim V = n$ and
  $\dim W = r$, prove that $\dim W^\perp = n - r$. Assuming $V$ is
  finite dimensional, prove that, under the canonical isomorphism
  $V^{\vee\vee} \cong V$, $W^{\perp\perp}$ maps isomorphically to
  $W$. [\ref{exer:linrep:annihilator-iso-dual-via-exact-seq-free},
  \ref{exer:linrep:grassmannian-bijection-with-dual-grassmannian}]
\end{exer}

\begin{exer}{}{exer:linrep:annihilator-iso-dual-via-exact-seq-free}
  Let $0 \to F_1 \to F_2 \to F_3 \to 0$ be an exact sequence of free
  $R$-modules. Viewing $F_1$ as a submodule of $F_2$ and extrapolating
  the notation introduced for vector spaces in
  \Cref{exer:linrep:annihilator-of-subspace-basics}, prove that
  $F_1^\perp \cong F_3^\vee$.
\end{exer}

\begin{exer}{$\lnot$}{exer:linrep:grassmannian-bijection-with-dual-grassmannian}
  Let $V$ be a vector space of dimension $n$. Prove that there is a
  natural bijection between the Grassmannian $\mathrm{Gr}(r, V)$ of
  $r$-dimensional subspaces of $V$ (cf.\
  \Cref{exer:linalg:grassmannian-cell-decomposition}) and the
  Grassmannian $\mathrm{Gr}(n-r, V^\vee)$ of $(n-r)$-dimensional
  subspaces of the dual of $V$. (Use
  \Cref{exer:linrep:annihilator-of-subspace-basics}.) In particular,
  the Grassmannian $\mathrm{Gr}_k(n-1, n)$ has the same structure as
  the projective space $\mathbf{P}V = \mathrm{Gr}_k(1, n)$. We could
  in fact define the projective space associated to a vector space $V$
  of dimension $n$ to be the set of subspaces of `codimension $1$'
  (that is, dimension $n-1$) in $V$. There are reasons why this would
  be preferable\footnote{The matter is not moot: it is true that there
    is a bijection between $\mathbf{P}V$ and $\mathrm{Gr}(n-1,V)$, but
    this bijection depends on the choice of a basis. Thus we cannot
    simply identify $\mathbf{P}V$ and $\mathrm{Gr}(n-1,V)$. We can
    identify $\mathbf{P}V$ and $\mathrm{Gr}(n-1,V^\vee)$, and it is
    sometimes useful to do so.}, but established conventions are what
  they are. [\ref{exer:linalg:grassmannian-cell-decomposition}]
\end{exer}

\begin{exer}{}{exer:linrep:wedge-pairing-induces-iso-wedge-dual}
  Let $F$ be a free $R$-module of finite rank. For any $r \geq 1$,
  define a multilinear map
  \[
    \delta : \underbrace{F \times \dots \times F}_{r} \times
    \underbrace{F^\vee \times \dots \times F^\vee}_{r} \to R
  \]
  by
  $\delta(v_1, \dots, v_r, \check{w}_1, \dots, \check{w}_r) =
  \det\bigl(\check{w}_i(v_j)\bigr)_{i,j=1,\dots,r}$.
  \begin{itemize}
  \item Prove that $\delta$ is multilinear and alternating in the
    first $r$ and in the last $r$ entries.
  \item Deduce that $\delta$ induces a bilinear map
    $\tilde{\delta} : \bigwedge^r(F) \times \bigwedge^r(F^\vee) \to
    R$.
  \item Prove that $\tilde{\delta}$ induces an isomorphism
    $\bigl(\bigwedge^r(F)\bigr)^\vee \cong \bigwedge^r(F^\vee)$.
  \end{itemize}
\end{exer}

\begin{exer}{}{exer:linrep:dual-has-right-module-structure-r-iso-rop}
  Let $R$ be a ring, not necessarily commutative, and let $M$ be a
  left-$R$-module. Prove that $M^\vee$ carries a natural
  right-$R$-module structure. Prove that $R$ is isomorphic as a ring
  to the opposite ring $R^\circ$. (Cf.\
  \Cref{exer:rings:opposite-ring}.)
\end{exer}

\section{Projective and injective modules and the Ext functors}
\label{sec:linrep:projective-and-injective-modules-and-the-ext-functors}

We have seen in \Cref{subsec:linrep:adjunction-again} that the Hom
functors are left-exact, and not exact in general. Studying this
matter further leads us along a path parallel to the one traveled in
the case of tensor products. In that case, we singled out a class of
special (`flat') modules for which $\otimes$ is exact, and we
discovered that the lack of exactness of tensor products may be
precisely measured in terms of a whole collection of new (`Tor')
functors. Precisely the same situation occurs for the covariant and
contravariant flavors of Hom, and I will close the chapter by
summarizing this story.

This brief overview should suffice for standard applications the
reader may encounter. I am also hoping that it will serve as
motivation for the next chapter: the wonderful facts that I will
present here without proof (on Ext functors, analogous to the facts
about Tor that were left without proof in
\Cref{subsec:linrep:the-tor-functors}) could be proven in detail by
ad hoc methods, without too much effort; but they find their most
natural setting in the wider apparatus of homological algebra and were
in fact at the root of its development. The reader who bears with us
throughout \Cref{chap:homol} will see complete proofs of all the facts
advertised here and substantially more.

As in the rest of the chapter, $R$ will be a commutative ring. As
elsewhere, the hypothesis of commutativity is essentially unnecessary;
it has the convenient consequence that the objects defined below (such
as $\Ext^i_R(M, N)$ for two $R$-modules $M$ and $N$) are naturally
endowed with an $R$-module structure. If $R$ is not necessarily
commutative, they are `just' abelian groups.

\subsection{Projectives and injectives}
\label{subsec:linrep:projectives-and-injectives}

\begin{defn}{}{defn:linrep:projective-injective-module}
  Let $R$ be a (commutative) ring.
  \begin{itemize}
  \item An $R$-module $P$ is \emph{projective} if the functor
    $\Hom_R(P, -)$ is exact.
  \item An $R$-module $Q$ is \emph{injective} if the functor
    $\Hom_R(-, Q)$ is exact.
  \end{itemize}
\end{defn}

Since Hom is left-exact in any case, these definitions admit the
following straightforward translations:

\begin{lem}{}{lem:linrep:projective-lifting-criterion}
  An $R$-module $P$ is projective if and only if for all epimorphisms
  of $R$-modules $\mu : M \to N$, every $R$-linear map $p : P \to N$
  lifts to an $R$-linear map $\hat{p} : P \to M$:
  \[
    \begin{tikzcd}
      & P \ar[d, "p"] \ar[dl, dashed, "\exists\,\hat{p}"'] \\
      M \ar[r, "\mu"'] & N \ar[r] & 0
    \end{tikzcd}
  \]
  An $R$-module $Q$ is injective if and only if for all monomorphisms
  of $R$-modules $\lambda : L \to M$, every $R$-linear map
  $q : L \to Q$ extends to an $R$-linear map $\hat{q} : M \to Q$:
  \[
    \begin{tikzcd}
      0 \ar[r] & L \ar[r, "\lambda"] \ar[d, "q"'] &
      M \ar[dl, dashed, "\exists\,\hat{q}"] \\
      & Q
    \end{tikzcd}
  \]
\end{lem}

\begin{proof}
  This is straightforward. Since $\Hom_R(-,Q)$ is left-exact for all
  $Q$, $Q$ is injective if and only if whenever a sequence
  $0 \to L \to M$ is exact, then so is the induced sequence
  $\Hom_R(M,Q) \to \Hom_R(L,Q) \to 0$. This translates precisely into
  the given condition. The argument for projective modules is entirely
  analogous.
\end{proof}

\begin{exam}{}{exam:linrep:ring-is-projective-not-injective-over-itself}
  Taken as a module over itself, $R$ is projective (because
  $\Hom_R(R, M)$ is canonically isomorphic to $M$, for all modules
  $M$) but not injective in general (since the duality functor
  $\Hom_R(-, R)$ is not exact in general; cf.\
  \Cref{subsec:linrep:adjunction-again}).
\end{exam}

A variation on the theme of \Cref{lem:linrep:projective-lifting-criterion} rephrases the
projective/injective conditions in terms of the splitting of sequences
(cf.\ the end of
\Cref{subsec:rings:complexes-and-exact-sequences}). For example,
assume that $P$ is projective; then I claim that every exact sequence
\[
  0 \to L \xrightarrow{\lambda} M \xrightarrow{\mu} P \to 0
\]
splits, in the sense that there is a submodule $P'$ of $M$ such that
$\mu$ restricts to an isomorphism $P' \xrightarrow{\sim} P$ and
$M \cong \lambda(L) \oplus P'$. Indeed, since $P$ is projective, the
identity $P \to P$ lifts to a homomorphism $\rho : P \to M$, and the
reader can then verify that $P' = \rho(P)$ fits the requirement.
Loosely speaking, in this situation we can simply replace $M$ by
$L \oplus P$.

Similarly, if $0 \to Q \to M \to N \to 0$ is exact and $Q$ is
injective, then the sequence necessarily splits. For example, $\Z$ is
not an injective $\Z$-module, since the exact sequence
\[
  0 \to \Z \xrightarrow{\cdot 2} \Z \to \Z/(2) \to 0
\]
manifestly does not split.

It is not difficult to show that these properties in fact characterize
projective and injective modules. With the terminology introduced in
\Cref{subsec:rings:complexes-and-exact-sequences}, a module $P$ is
projective if and only if every epimorphism $M \to P$ is a split
epimorphism, and a module $Q$ is injective if and only if every
monomorphism $Q \to M$ is a split monomorphism. Dotting all the i's
and crossing all the t's is left to the enterprising reader
(\Cref{exer:linrep:projective-injective-splitting-characterization}).

\subsection{Projective modules}
\label{subsec:linrep:projective-modules}
\Cref{lem:linrep:projective-lifting-criterion} does not necessarily help in acquiring a
feeling for projective or injective modules. I am reasonably happy
with projective modules, because of the following characterization:

\begin{prop}{}{prop:linrep:projective-iff-direct-summand-of-free}
  An $R$-module $P$ is projective if and only if it is a direct
  summand of a free module, that is, if and only if there exists a
  free $R$-module $F$, an $R$-module $K$, and an isomorphism
  $K \oplus P \cong F$.
\end{prop}

\begin{proof}
  Any set $S$ of generators of $P$ determines a surjection of the free
  module $F = R^{\oplus S}$ onto $P$ and hence an exact sequence
  $0 \to K \to F \to P \to 0$. As observed above, such a sequence
  necessarily splits if $P$ is projective; thus $F \cong K \oplus P$
  in this case, as needed.

  For the converse, assume that $K \oplus P \cong F$ for some free
  $R$-module $F$. It is clear that $F$ satisfies the lifting property
  of \Cref{lem:linrep:projective-lifting-criterion}:
  \[
    \begin{tikzcd}
      & F \cong K \oplus P \ar[d, "f"] \ar[dl, dashed,
      "\exists\,\hat{f}"'] \\
      M \ar[r, "\mu"'] & N \ar[r] & 0
    \end{tikzcd}
  \]
  (if $F = R^{\oplus S}$, define $\hat{f}(s)$ to be an arbitrary
  inverse image of $f(s)$, for all $s \in S$, and extend to $F$ by the
  universal property of free modules). It follows easily that the
  lifting property holds for $P$, by judicious use of the natural maps
  $P \hookrightarrow K \oplus P \twoheadrightarrow P$. Details are
  left to the reader.
\end{proof}

For example, free modules are projective.
\Cref{prop:linrep:projective-iff-direct-summand-of-free} streamlines the verification of simple
properties of projective modules:

\begin{coro}{}{coro:linrep:projective-modules-are-flat}
  Let $P_1$, $P_2$ be projective $R$-modules. Then $P_1 \oplus P_2$
  and $P_1 \otimes_R P_2$ are projective. Projective modules are flat.
\end{coro}

\begin{proof}
  These statements follow easily from \Cref{prop:linrep:projective-iff-direct-summand-of-free},
  the fact that $\otimes$ is distributive with respect to $\oplus$,
  and the fact that free modules are flat.
\end{proof}

The categorically minded reader should look for alternate proofs of
these facts. For example, if $P_1$ and $P_2$ are both projective, the
functor $\Hom_R(P_1, \Hom_R(P_2, -))$ is exact; by adjunction, this is
the functor $\Hom_R(P_1 \otimes_R P_2, -)$; since this is exact,
$P_1 \otimes_R P_2$ is projective. A proof via adjunction of the fact
that projective modules are flat is straightforward
(\Cref{exer:linrep:projective-modules-are-flat-via-adjunction}), using
the (nontrivial) fact that $R\text{-Mod}$ has enough injectives,
discussed below. As it happens, finitely generated flat modules over,
e.g., Noetherian rings, are projective; the diligent reader has
already done all the work necessary to prove this, and will wrap it up
in
\Cref{exer:linrep:fin-gen-over-noeth-locally-free-iff-projective-iff-flat}.

The fact that free modules are projective implies that $R\text{-Mod}$
has \emph{enough projectives}: this means that every $R$-module $M$
admits a \emph{projective resolution}, that is, an exact sequence
\[
  \dots \to P_3 \to P_2 \to P_1 \to P_0 \to M.
\]
Indeed, a free resolution of $M$ is in particular a projective
resolution. Since there are in general `more' projective modules than
free modules, it can in principle be `easier' to construct a
projective resolution. The magic of homological algebra shows that
projective resolutions may be used in place of free resolutions in
(for example) computing Tor-modules. This will be clear once we go
through some homological algebra machinery, in \Cref{chap:homol},
especially \Cref{sec:homol:double-complexes}.

\subsection{Injective modules}
\label{subsec:linrep:injective-modules}
I may have inadvertently communicated to the reader the message that
every categorical notion leads in a straightforward way to a mirror
notion, by `just reversing arrows'. This would seem to be the case for
injective vs.\ projective modules: the definition of one kind is
indeed obtained from the definition of the other by simply reversing
arrows in the key diagrams. However, this is as far as this naive
perception can go. A seemingly unavoidable asymmetry of nature
manifests itself here, making injective modules look more mysterious
than projectives. For example, there is no `easy' description of
injective modules in the style of \Cref{prop:linrep:projective-iff-direct-summand-of-free}
(however, cf.\ \Cref{rem:linrep:injective-cogenerator-role}), and while it is trivially
the case that an arbitrary module is surjected upon by a projective
(e.g., free) module, the `mirror' statement for injectives is
certainly not trivial.

The problem is that it is a little hard to picture an injective
module: while $R$ itself is a trivial example of a projective
$R$-module, at this point the reader would likely find it hard to name
a single nonzero injective module over any ring whatsoever. Well, say
over any ring that is not a field: if $R$ is a field, then $R$ itself
is injective. (Why?) But if $R$ is not a field\dots? Suppose $r$ is a
non-zero-divisor in $R$, and consider the inclusion of the ideal
$(r) \cong R$ in $R$:
\[
  0 \to R \xrightarrow{\cdot r} R;
\]
extending the identity on $R$ through the multiplication by $r$,
\[
  \begin{tikzcd}
    0 \ar[r] & R \ar[r, "\cdot r"] \ar[d, "\id"'] &
    R \ar[dl, dashed, "\exists?"] \\
    & R
  \end{tikzcd}
\]
requires the existence of $s \in R$ such that $rs = 1$. It cannot be
done unless $r$ is a unit in $R$.

This tells us that in general $R$ is not injective as an $R$-module
and puts the finger on the `problem': for $Q$ to be injective, we must
in particular be able to extend any map $I \to Q$ from an ideal $I$ of
$R$ to a map $R \to Q$. A nice application of Zorn's lemma shows that
this property characterizes injective modules:

\begin{thm}{(Baer criterion)}{thm:linrep:baer-criterion}
  An $R$-module $Q$ is injective if and only if every $R$-linear map
  $f : I \to Q$, with $I$ an ideal of $R$, extends to an $R$-linear
  map $\hat{f} : R \to Q$.
\end{thm}

This criterion is attributed to \name{Reinhold Baer}, who first
introduced and studied injective modules (ca.\ 1940).

\begin{proof}
  The `only if' part of the statement is immediate from the definition
  of injective. To verify the `if' part, assume $Q$ satisfies the
  stated extension condition, let $L \subseteq M$ be any inclusion of
  $R$-modules, and let $q : L \to Q$ be a given $R$-linear map:
  \[
    \begin{tikzcd}
      0 \ar[r] & L \ar[r] \ar[d, "q"'] & M \ar[dl, dashed,
      "\exists?\,\hat{q}"] \\
      & Q
    \end{tikzcd}
  \]
  We want to construct an extension $\hat{q} : M \to Q$ of $q$.

  Consider the set $\mathscr{E}$ of pairs $(\tilde{L}, \tilde{q})$,
  where $\tilde{L}$ ranges over the submodules of $M$ containing $L$
  and $\tilde{q} : \tilde{L} \to Q$ extends $q$: $\tilde{q}|_L =
  q$. Then $\mathscr{E}$ is nonempty, since $(L, q) \in \mathscr{E}$,
  and we can define a partial order on $\mathscr{E}$ by prescribing
  that $(\tilde{L}', \tilde{q}') \preceq (\tilde{L}'', \tilde{q}'')$
  if $\tilde{L}' \subseteq \tilde{L}''$ and $\tilde{q}''$ extends
  $\tilde{q}'$. It is clear that every chain has an upper bound (take
  the union of the corresponding $\tilde{L}$), so by Zorn's lemma
  there exists a maximal extension $(\tilde{L}, \tilde{q})$, and we
  are done if we can prove that $\tilde{L} = M$.

  Arguing by contradiction, assume there exists $m \in M$,
  $m \notin \tilde{L}$, and consider the proper ideal
  $I = \{r \in R \mid rm \in \tilde{L}\}$ of $R$. We can define an
  $R$-linear map $f : I \to Q$ by putting $f(r) := \tilde{q}(rm)$; by
  hypothesis, this map extends to an $R$-linear map
  $\hat{f} : R \to Q$. But then we can use this map to extend $q$
  beyond $\tilde{L}$: let
  $L' = \tilde{L} + \langle m \rangle \supsetneq \tilde{L}$, and
  define $q' : L' \to Q$ by
  \[
    q'(\ell + rm) := \tilde{q}(\ell) + \hat{f}(r)
  \]
  for $\ell \in \tilde{L}$, $r \in R$. This is immediately checked to
  be well-defined and clearly extends $\tilde{q}$. But then
  $(L', q') \in \mathscr{E}$ and
  $(\tilde{L}, \tilde{q}) \prec (L', q')$, contradicting the
  maximality of $(\tilde{L}, \tilde{q})$ and concluding the proof.
\end{proof}

\Cref{thm:linrep:baer-criterion} completely clarifies the notion of injective
modules in the case of PIDs. An $R$-module $D$ is \emph{divisible} if
$rD = D$ for every non-zero-divisor $r \in R$, that is, if we can
`divide' every element of $D$ by every non-zero-divisor of $R$.

\begin{coro}{}{coro:linrep:injective-over-pid-iff-divisible}
  Let $R$ be a PID. Then an $R$-module $Q$ is injective if and only if
  it is divisible.
\end{coro}

\begin{proof}
  \Cref{exer:linrep:injective-over-pid-iff-divisible}.
\end{proof}

\begin{exam}{}{exam:linrep:q-and-q-mod-z-are-injective}
  Viewed as abelian groups (i.e., $\Z$-modules), $\Q$ and $\Q/\Z$ are
  injective. More generally, if $D$ is any divisible abelian group and
  $K \subseteq D$, then $D/K$ is injective. (Indeed, it is trivially
  divisible!)
\end{exam}
Going from friendly PIDs like $\Z$ to arbitrary rings may still seem
challenging, but here is where some of our previous work pays
off. Recall that for every ring homomorphism $f : S \to R$ we have
defined in
\Cref{subsec:linrep:restriction-and-extension-of-scalars} a
functor\footnote{Watch out --- in
  \Cref{subsec:linrep:restriction-and-extension-of-scalars} the
  homomorphism goes from $R$ to $S$, so the definition of $f^!$ may
  look backwards here.}  $f^! : S\text{-Mod} \to R\text{-Mod}$, by
setting $f^!(M) = \Hom_S(R, M)$. This functor preserves injectives:

\begin{lem}{}{lem:linrep:restriction-of-scalars-preserves-injectives}
  Let $f : S \to R$ be a homomorphism of commutative rings, and let
  $Q$ be an injective $S$-module. Then $f^!(Q)$ is an injective
  $R$-module.
\end{lem}

\begin{rem}{}{rem:linrep:restriction-of-scalars-preserves-projectives}
  Similarly, $f^*$ preserves projectives
  (\Cref{exer:linrep:restriction-of-scalars-of-projective-is-projective}).
\end{rem}

\begin{proof}
  By adjunction (\Cref{lem:linrep:hom-tensor-adjunction-bimodule}),
  $\Hom_R(-, f^!(Q)) \cong \Hom_S(f^*(-), Q)$ as functors
  $R\text{-Mod} \to \mathrm{Ab}$. Since $f^*$ is exact
  (\Cref{prop:linrep:restriction-extension-coextension-adjunctions}) and $\Hom_S(-, Q)$ is exact by
  hypothesis, $\Hom_R(-, f^!(Q))$ must be exact.  This is precisely
  the statement.
\end{proof}

Since every ring $R$ admits a (unique) map $\iota : \Z \to R$, we have
a good source of injective objects in $R\text{-Mod}$ for every $R$:

\begin{coro}{}{coro:linrep:hom-from-divisible-group-is-injective}
  Let $R$ be a commutative ring. If $D$ is any divisible abelian
  group, then $\iota^!(D) = \Hom_{\Z}(R, D)$ is injective in
  $R\text{-Mod}$.
\end{coro}

This result is as close as I can get to an `intuitive' feeling for the
notion of injective module. For example, at this point even I see why
`$R\text{-Mod}$ has enough injectives':

\begin{coro}{}{coro:linrep:module-embeds-in-injective}
  Let $M$ be an $R$-module. Then $M$ can be identified with a
  submodule of an injective $R$-module.
\end{coro}

\begin{proof}
  I claim that it suffices to show that $\Z\text{-Mod}$ has enough
  injectives. Indeed, this will show that there exists a divisible
  abelian group $D$ such that $M \subseteq D$ ($M$ is in particular an
  abelian group); since $R$-linear maps are in particular $\Z$-linear,
  \[
    M \cong \Hom_R(R, M) \hookrightarrow \Hom_{\Z}(R, M) \subseteq
    \Hom_{\Z}(R, D),
  \]
  and the rightmost module is injective by
  \Cref{coro:linrep:hom-from-divisible-group-is-injective}.

  Thus, we are reduced to the case $R = \Z$ and $M = A$ an abelian
  group. There is a surjective homomorphism from $\Z^{\oplus A}$ to
  $A$ and hence an identification $A \cong \Z^{\oplus A}/K$ for some
  $K \subseteq \Z^{\oplus A}$. Embedding $\Z^{\oplus A}$ in
  $\Q^{\oplus A}$, we obtain an injective homomorphism
  $A \hookrightarrow \Q^{\oplus A}/K$, and we are done since
  $\Q^{\oplus A}/K$ is divisible (\Cref{exam:linrep:q-and-q-mod-z-are-injective}).
\end{proof}

As an immediate consequence, every $R$-module admits an
\emph{injective resolution}: an exact sequence
\[
  0 \to M \to Q^0 \to Q^1 \to Q^2 \to Q^3 \to \dots
\]
in which every $Q^i$ is injective. Indeed, this is obtained by
applying \Cref{coro:linrep:module-embeds-in-injective} to $M$ to construct $Q^0$ and
then applying it again to the cokernel of $M \to Q^0$ to construct
$Q^1$, and so on.

The fact that $R\text{-Mod}$ has enough injectives has a
generalization to sheaves, which is an essential ingredient in the
definition of sheaf cohomology, an extremely important tool in modern
algebraic geometry.

\begin{rem}{}{rem:linrep:injective-cogenerator-role}
  \Cref{coro:linrep:module-embeds-in-injective} may be refined, proving that every
  module can in fact be realized as a submodule of a product
  $\Hom_{\Z}(R, \Q/\Z)^S$ (this is injective, by
  \Cref{coro:linrep:hom-from-divisible-group-is-injective} and
  \Cref{exer:linrep:direct-sum-of-injectives-iff-each-is}).  That is,
  products of $\iota^!(\Q/\Z) = \Hom_{\Z}(R, \Q/\Z)$ play for
  injectives the role played by free modules (= coproducts of
  $R = \iota^*(\Z)$) for projectives. (Not surprisingly, they are said
  to be \emph{cofree}.)  It follows that every injective $R$-module is
  a direct summand of $\Hom_{\Z}(R, \Q/\Z)^S$ for some $S$.
\end{rem}

\subsection{The Ext functors}
\label{subsec:linrep:the-ext-functors}
The next step to take now is the construction of tools to `quantify'
the lack of exactness of $\Hom_R$, in the same sense that the functors
$\Tor^R_i$ quantify the lack of exactness of $\otimes_R$. The
corresponding functors are called $\Ext^i_R$.

Since there are two functors associated with Hom (the covariant
$\Hom_R(M, -)$ and the contravariant $\Hom_R(-, N)$), it would be
natural to expect two corresponding sequences of functors.  Amazingly,
the same `bifunctors' $\Ext^i_R(-, -)$ work for both! For all
$R$-modules $M$ and $N$, we have the following:
\begin{itemize}
\item For all $i \geq 0$, $\Ext^i_R(M, -)$ and $\Ext^i_R(-, N)$ are
  covariant, resp., contravariant, functors
  $R\text{-Mod} \to R\text{-Mod}$.
\item For $i = 0$, $\Ext^0_R(M, N) = \Hom_R(M, N)$.
\item For every exact sequence $0 \to A \to B \to C \to 0$ of
  $R$-modules, there are long exact sequences
  \begin{gather*}
    0 \to \Hom_R(M,A) \to \Hom_R(M,B) \to
    \Hom_R(M,C) \xrightarrow{\delta_0}
    \Ext^1_R(M,A) \to \Ext^1_R(M,B) \to
    \Ext^1_R(M,C) \xrightarrow{\delta_1} \cdots
  \end{gather*}
  and
  \begin{gather*}
    0 \to \Hom_R(C,N) \to \Hom_R(B,N) \to
    \Hom_R(A,N) \xrightarrow{\delta^0}
    \Ext^1_R(C,N) \to \Ext^1_R(B,N) \to
    \Ext^1_R(A,N) \xrightarrow{\delta^1} \cdots
  \end{gather*}
  for suitable natural connecting morphisms $\delta_i$, $\delta^i$.
\end{itemize}
This will be proven in \Cref{chap:homol}. In any case, just knowing
that such wonderful things exist gives us enough information to be
able to play with them a little. For example,
\begin{itemize}
\item if $\Ext^1_R(P, -) = 0$, then $P$ is projective;
\item if $\Ext^1_R(-, Q) = 0$, then $Q$ is injective.
\end{itemize}
Indeed, the exactness of the corresponding $\Hom_R(P, -)$, resp.,
$\Hom_R(-, Q)$, follows immediately from the exact sequences displayed
above.  In fact, pushing the envelope a little will tell us much
more. Recall the procedure that defines $\Tor^R_i(-, N)$ from the
functor $- \otimes_R N$: to compute the value of these functors at a
module $M$,
\begin{itemize}
\item find a free resolution of $M$:
  $\dots \to R^{\oplus S_2} \to R^{\oplus S_1} \to R^{\oplus S_0} \to
  M \to 0$;
\item apply the functor $- \otimes_R N$ to the free part, obtaining a
  complex $M_\bullet \otimes_R N$:
  $\dots \to N^{\oplus S_2} \to N^{\oplus S_1} \to N^{\oplus S_0} \to
  0$;
\item define Tor by taking the homology of this complex:
  $\Tor^R_i(M, N) := H_i(M_\bullet \otimes N)$.
\end{itemize}
I also pointed out that projective resolutions may be used in place of
free resolutions.\footnote{In fact, resolutions by flat modules
  suffice; this will be proven in \Cref{chap:homol}.} This strategy is
an example of the general procedure used to `derive' functors. Can you
dream up how the same strategy may be applied in order to compute
Ext-modules? Don't read ahead until you have thought a little about
this!

$*$ Welcome back. In the case of $\Ext^i_R(M,N)$, there would seem to
be two natural and possibly different ways to go:
\begin{itemize}
\item One could start from a projective resolution of $M$:
  $\dots \to P_2 \to P_1 \to P_0 \to M \to 0$; apply the contravariant
  $\Hom_R(-, N)$ to the projective part, obtaining a new complex
  $\Hom_R(M_\bullet, N)$, i.e.,
  $\Hom_R(P_0,N) \to \Hom_R(P_1,N) \to \Hom_R(P_2,N) \to \dots$, and
  take cohomology,\footnote{Recall from
    \Cref{subsec:rings:complexes-and-exact-sequences} that the
    homology of complexes with increasing (rather than decreasing)
    indices is called cohomology; this is usually denoted $H^*$, with
    `upper' indices.} proposing the definition
  $\Ext^i_R(M,N) := H^i(\Hom_R(M_\bullet, N))$.
\item Or one could start from an injective resolution of $N$:
  $0 \to N \to Q^0 \to Q^1 \to Q^2 \to \dots$; apply the covariant
  $\Hom_R(M, -)$ to the injective part, obtaining a new complex
  $\Hom_R(M, N^\bullet)$, i.e.,
  $\Hom_R(M,Q^0) \to \Hom_R(M,Q^1) \to \Hom_R(M,Q^2) \to \dots$, and
  again take cohomology, leading to
  $\Ext^i_R(M,N) := H^i(\Hom_R(M, N^\bullet))$.
\end{itemize}
These definitions are independent of all choices and agree with each
other! Again, the reader will have to wait for \Cref{chap:homol} to
see a proof that this is the case and that the resulting modules do
satisfy the requirements spelled out at the beginning of this
subsection. As in the case of Tor, the impatient reader can get a good
feel for the needed arguments by working out the case of finitely
generated modules over PIDs, making good use of the snake lemma.

The reader who is willing to take my word for now and believe all
these beautiful results is already in a position to perform many
computations and to understand injective/projective modules (even)
better. For example,

\begin{prop}{}{prop:linrep:projective-injective-via-ext-vanishing}
  An $R$-module $P$ is projective if and only if $\Ext^1_R(P, -) = 0$,
  if and only if $\Ext^i_R(P, -) = 0$ for all $i > 0$. An $R$-module
  $Q$ is injective if and only if $\Ext^1_R(-, Q) = 0$, if and only if
  $\Ext^i_R(-, Q) = 0$ for all $i > 0$.
\end{prop}

\begin{proof}
  The second assertion: we have seen that $Q$ is injective if
  $\Ext^1_R(-, Q) = 0$, and this is trivially the case if
  $\Ext^i_R(-, Q) = 0$ for all $i > 0$. So we just have to prove that
  $\Ext^i_R(-, Q) = 0$ for all $i > 0$ if $Q$ is injective. This is
  immediate from the second definition given above: if $Q$ is
  injective, then $0 \to Q \to Q \to 0 \to 0 \to \dots$ is an
  injective resolution of $Q$; thus, for an $R$-module $M$,
  $\Ext^i_R(M, Q)$ is the cohomology of the complex
  $\Hom_R(M,Q) \to 0 \to 0 \to \dots$, giving
  $\Ext^0_R(M,Q) = \Hom_R(M,Q)$ and $\Ext^i_R(M,Q) = 0$ for $i > 0$
  and all $M$. The first assertion is proven similarly, using the
  first definition given above for Ext.
\end{proof}

\subsection{\texorpdfstring{$\Ext^*_{\Z}(G, \Z)$}{Ext*\_Z(G, Z)}}
\label{subsec:linrep:ext-Z-G-Z}
I will attempt to convince the reader that computing Ext modules is
not too unreasonable, by discussing the computation of
$\Ext^*_{\Z}(G, \Z)$ for an arbitrary finitely generated abelian group
$G$.
\begin{itemize}
\item Since Hom commutes with finite direct sums
  (\Cref{coro:linrep:hom-commutes-with-products}), so does Ext; this is an easy
  consequence of the definitions given above.
\item By the classification theorem for finitely generated abelian
  groups, we are reduced to computing $\Ext^i_{\Z}(\Z, \Z)$ and
  $\Ext^i_{\Z}(\Z/m\Z, \Z)$, for all $m > 0$ and $i \geq 0$.
\item The first module is $\Hom_{\Z}(\Z,\Z) \cong \Z$ for $i = 0$, and
  it is $0$ for $i > 0$: indeed, this follows from
  \Cref{prop:linrep:projective-injective-via-ext-vanishing}, since $\Z$ is projective.
\item The second module: use the projective resolution
  $0 \to \Z \xrightarrow{\cdot m} \Z \to \Z/m\Z \to 0$; this yields
  that $\Ext^i_{\Z}(\Z/m\Z, \Z)$ is the cohomology of the complex
  \[
    0 \to \Hom_{\Z}(\Z,\Z) \xrightarrow{\cdot m} \Hom_{\Z}(\Z,\Z) \to
    0 \to \dots,
  \]
  that is, $0 \to \Z \xrightarrow{\cdot m} \Z \to 0 \to \dots$,
  confirming that
  $\Ext^0_{\Z}(\Z/m\Z, \Z) = \Hom_{\Z}(\Z/m\Z, \Z) = 0$ ($\Z$ has no
  torsion) and computing $\Ext^1_{\Z}(\Z/m\Z, \Z) \cong \Z/m\Z$, with
  vanishing $\Ext^2$ and higher.
\item The conclusion is that $\Ext^0_{\Z}(G, \Z)$ picks up the free
  part of $G$, while $\Ext^1_{\Z}(G, \Z)$ is isomorphic to the torsion
  part, and $\Ext^i_{\Z}(G, \Z) = 0$ for $i \geq 2$ (as should have
  been expected; cf.\
  \Cref{exer:linrep:ext-vanishes-above-degree-1-over-pid}).
\end{itemize}
In particular, if $G$ is a finite abelian group, then
$G \cong \Ext^1_{\Z}(G, \Z)$; as it happens, this isomorphism is not
canonical.

Another Ext computation gives a different interpretation of this
group. The exact sequence $0 \to \Z \to \Q \to \Q/\Z \to 0$ is an
injective resolution of $\Z$; thus, $\Ext^*_{\Z}(G, \Z)$ is the
cohomology of
\[
  0 \to \Hom_{\Z}(G, \Q) \to \Hom_{\Z}(G, \Q/\Z) \to 0 \to \dots.
\]
If $G$ is torsion (for example, if $G$ is finite), it follows that
$\Ext^1_{\Z}(G, \Z)$ is isomorphic to $\Hom_{\Z}(G, \Q/\Z)$. This
group is the \emph{Pontryagin dual} of $G$, denoted $\hat{G}$. Thus,
we have verified that if $G$ is a finite abelian group, then $G$ is
isomorphic (not canonically) to its Pontryagin dual
$\hat{G} = \Hom_{\Z}(G, \Q/\Z)$. Of course this fact is not difficult
to check directly, but it is immediate from the point of view of Ext.

Finally, the notation \emph{Ext} is due to the fact that one can
construct a bijection between the group $\Ext^1_R(M, N)$ and the set
of equivalence classes of \emph{extensions} of $M$ by $N$, that is, of
exact sequences $0 \to N \to E \to M \to 0$, modulo a suitable
isomorphism relation (cf.\
\Cref{subsec:groups2:exact-sequences-of-groups-extension-problem}). For
example, the direct sum $E = M \oplus N$ corresponds to the $0$
element of this Ext group. The diligent reader will construct this
bijection in the exercises.

As usual, $R$ denotes a fixed commutative ring.

\subsection*{Exercises}
\begin{exer}{$*$}{exer:linrep:projective-injective-splitting-characterization}
  Prove that an $R$-module $P$ is projective if and only if every
  epimorphism $M \to P$ splits and $Q$ is injective if and only if
  every monomorphism $Q \to M$
  splits. [\Cref{subsec:linrep:projectives-and-injectives}]
\end{exer}

\begin{exer}{}{exer:linrep:generalize-prop-to-projective-modules}
  Prove that the result of \Cref{prop:linrep:dual-sequence-exact-when-quotient-free} holds more
  generally whenever $P$ is a projective module.
\end{exer}

\begin{exer}{$\lnot$}{exer:linrep:injective-iff-hom-surjective-for-injectives}
  Prove that an $R$-linear map $A \to B$ is injective if and only if
  the induced map $\Hom_R(B, Q) \to \Hom_R(A, Q)$ is surjective for
  all injective modules $Q$. (Hint: $R\text{-Mod}$ has enough
  injectives.) This strengthens the result of
  \Cref{exer:linrep:epi-mono-characterization-via-hom}.
  [\ref{exer:linrep:projective-modules-are-flat-via-adjunction}]
\end{exer}

\begin{exer}{$\lnot$}{exer:linrep:direct-sum-of-projectives-iff-each-is}
  Let $\{P_i\}_{i \in I}$ be a family of $R$-modules. Prove that
  $\bigoplus_i P_i$ is projective if and only if each $P_i$ is
  projective. [\ref{exer:homol:direct-sum-projectives-injectives-abelian-category}]
\end{exer}

\begin{exer}{}{exer:linrep:dual-of-fin-gen-projective-is-projective-reflexive}
  Prove that the dual of a finitely generated projective module is
  projective. Prove that finitely generated projective modules are
  reflexive (that is, isomorphic to their biduals).
\end{exer}

\begin{exer}{$\lnot$}{exer:linrep:restriction-of-scalars-of-projective-is-projective}
  Prove that if $f : S \to R$ is a ring homomorphism and $P$ is a
  projective $S$-module, then $f^*(P)$ is a projective
  $R$-module. [\Cref{subsec:linrep:injective-modules}]
\end{exer}

\begin{exer}{$*$}{exer:linrep:projective-modules-are-flat-via-adjunction}
  Give a proof of the fact that projective modules are flat, using
  adjunction. (Hint: $R\text{-Mod}$ has enough injectives, and
  \Cref{exer:linrep:injective-iff-hom-surjective-for-injectives}.)
  [\Cref{subsec:linrep:projective-modules},
  \ref{exer:linrep:combine-results-reprove-rank-thm}]
\end{exer}

\begin{exer}{}{exer:linrep:combine-results-reprove-rank-thm}
  Prove that
  \Cref{exer:linrep:fin-gen-flat-over-noeth-local-is-free,exer:linrep:projective-modules-are-flat-via-adjunction}
  together imply the result of
  \Cref{exer:linalg:direct-summand-of-free-over-local-is-free}.
\end{exer}

\begin{exer}{}{exer:linrep:vector-spaces-flat-injective-projective}
  Prove that vector spaces are flat, injective, and projective. (Try
  to do this directly, without invoking Baer's criterion.)
\end{exer}

\begin{exer}{}{exer:linrep:fin-gen-over-pid-free-iff-projective-iff-flat}
  Prove that a finitely generated module over a PID is free if and
  only if it is projective if and only if it is flat. (Do this `by
  hand', without invoking
  \Cref{exer:linrep:fin-gen-over-noeth-locally-free-iff-projective-iff-flat}.)
\end{exer}

\begin{exer}{$\lnot$}{exer:linrep:fin-gen-over-local-ring-projective-iff-free}
  Prove that a finitely generated module over a local ring is
  projective if and only if it is free. (You have already done this,
  in \Cref{exer:linalg:direct-summand-of-free-over-local-is-free}.)
  Besides PIDs and local rings, there are other classes of rings for
  which projective and free modules coincide.  Topological
  considerations suggested to \name{Serre} that projective modules
  over a polynomial ring over a field should necessarily be free, but
  it took two decades to prove that this is indeed the case.
  [\ref{exer:linrep:fin-gen-over-noeth-locally-free-iff-projective-iff-flat}]
\end{exer}

\begin{exer}{$\lnot$}{exer:linrep:fin-gen-over-noeth-locally-free-iff-projective-iff-flat}
  An $R$-module $M$ is `locally free' if $M_{\mathfrak{p}}$ is free as
  an $R_{\mathfrak{p}}$-module for every prime ideal $\mathfrak{p}$ of
  $R$. Prove that a finitely generated module over a Noetherian ring
  is locally free if and only if it is projective if and only if it is
  flat. (Use the results of
  \Cref{exer:linrep:fin-gen-flat-over-noeth-local-is-free,exer:linrep:fin-gen-over-local-ring-projective-iff-free}.)
  The hypothesis
  of finite generation is necessary (cf.\
  \Cref{exer:linrep:q-is-flat-not-projective-as-z-module}).
  [\ref{exer:linrep:fin-gen-over-pid-free-iff-projective-iff-flat}]
\end{exer}

\begin{exer}{$\lnot$}{exer:linrep:q-is-flat-not-projective-as-z-module}
  Prove that $\Q$ is flat, but not projective, as a
  $\Z$-module. [\ref{exer:linrep:fin-gen-over-noeth-locally-free-iff-projective-iff-flat}]
\end{exer}

\begin{exer}{$*$}{exer:linrep:injective-over-pid-iff-divisible}
  Prove that a module over a PID is injective if and only if it is
  divisible. (Use Baer's criterion.)
  [\Cref{subsec:linrep:injective-modules}]
\end{exer}

\begin{exer}{$*$}{exer:linrep:direct-sum-of-injectives-iff-each-is}
  Prove that $Q_1 \oplus Q_2$ is injective if and only if $Q_1$, $Q_2$
  are both injective. More generally, prove that if
  $\{Q_i\}_{i \in I}$ is a family of $R$-modules, then $\prod_i Q_i$
  is injective if and only if each $Q_i$ is injective.
  [\Cref{subsec:linrep:injective-modules},
  \ref{exer:homol:direct-sum-projectives-injectives-abelian-category}]
\end{exer}

\begin{exer}{}{exer:linrep:z-plus-q-flat-not-projective-not-injective}
  Prove that $\Z \oplus \Q$ is flat as a $\Z$-module but neither
  projective nor injective.
\end{exer}

\begin{exer}{}{exer:linrep:hom-iso-to-ext0}
  Prove that $\Hom_R(M, N) \cong \Ext^0_R(M, N)$, using any of the
  definitions provided for $\Ext^i_R(M, N)$.
\end{exer}

\begin{exer}{$*$}{exer:linrep:ext-vanishes-above-degree-1-over-pid}
  Prove that if $R$ is a PID, then $\Ext^i_R(M, N) = 0$ for all
  $i \geq 2$ and all $R$-modules $M$,
  $N$. [\Cref{subsec:linrep:ext-Z-G-Z}]
\end{exer}

\begin{exer}{}{exer:linrep:compute-ext-of-r-mod-r-with-m}
  Let $r$ be a non-zero-divisor in $R$, and let $M$ be an
  $R$-module. Compute all $\Ext^i_R(R/(r), M)$.
\end{exer}

\begin{exer}{}{exer:linrep:ext1-of-quotient-iso-to-hom-normal-bundle}
  Let $I$ be an ideal of $R$. Prove that
  $\Ext^1_R(R/I, R/I) \cong \Hom_R(I/I^2, R/I)$.  (Cf.\
  \Cref{exer:linrep:hom-i-r-mod-i-restriction-map-zero}. This says
  that this Ext module essentially computes the normal bundle of an
  embedding.)
\end{exer}

\begin{exer}{$\lnot$}{exer:linrep:why-ext-is-called-ext-part1}
  In \Cref{exer:linrep:tor1-r-torsion-of-n} we have seen why Tor is
  called Tor. Why is Ext called Ext? Let $M$, $N$, $E$ be
  $R$-modules. An \emph{extension} of $M$ by $N$ is an exact sequence
  \[
    \tag{$*$} 0 \to N \to E \to M \to 0.
  \]
  Two extensions are `isomorphic' if there is a commutative diagram
  \[
    \begin{tikzcd}
      0 \ar[r] & N \ar[r] \ar[d, equal] & E \ar[r] \ar[d] &
      M \ar[r] \ar[d, equal] & 0 \\
      0 \ar[r] & N \ar[r] & E' \ar[r] & M \ar[r] & 0
    \end{tikzcd}
  \]
  linking them. (Note that, by the snake lemma, the middle arrow must
  then be an isomorphism: cf.\
  \Cref{exer:rings:short-five-lemma-via-snake-lemma}.) Extensions that
  are isomorphic to the standard sequence
  $0 \to N \to N \oplus M \to M \to 0$ are `trivial'.

  Every extension $E$ as above determines an element
  $\xi = \xi(E) \in \Ext^1_R(M, N)$ as follows. The sequence $(*)$
  induces a long exact sequence of Ext, i.e.,
  \[
    \dots \to \Hom_R(E, N) \to \Hom_R(N, N) \to \Ext^1_R(M, N) \to
    \dots,
  \]
  and we let $\xi(E)$ be the image in $\Ext^1_R(M, N)$ of the
  distinguished element $\id_N \in \Hom_R(N, N)$.

  Prove that if $E$ and $E'$ are isomorphic extensions, then
  $\xi(E) = \xi(E')$. Prove that if $E$ is a trivial extension, then
  $\xi(E) = 0$. [\ref{exer:linrep:ext1-bijection-with-extensions}]
\end{exer}

\begin{exer}{}{exer:linrep:ext1-bijection-with-extensions}
  \Cref{exer:linrep:why-ext-is-called-ext-part1} teaches us that
  extensions determine elements of $\Ext^1_R$. We get an even sharper
  statement by constructing an inverse to the map $\xi$: for every
  element $\kappa \in \Ext^1_R(M, N)$, we will construct an extension
  $e(\kappa)$ such that $\xi(e(\kappa)) = \kappa$ and such that
  $e(\xi(E))$ is isomorphic to $E$.
  \begin{itemize}
  \item Let $F$ be any free module surjecting onto $M$, and let
    $\iota : K \hookrightarrow F$ be the kernel of
    $\pi : F \twoheadrightarrow M$. Since $F$ is free (hence
    projective), a piece of the long exact sequence of Ext, i.e.,
    \[
      \dots \to \Hom_R(K, N) \to \Ext^1_R(M, N) \to \Ext^1_R(F, N) = 0
      \to \dots,
    \]
    tells us that there exists a homomorphism $k : K \to N$ mapping to
    the element $\kappa \in \Ext^1_R(M, N)$.
  \item We now have a monomorphism $(\iota, k) : K \to F \oplus N$.
    Let $E$ be the cokernel $(F \oplus N)/K$. Prove that the
    epimorphism $(\pi, 0) : F \oplus N \to M$ factors through this
    cokernel, defining an epimorphism $E \twoheadrightarrow M$.
  \item Prove that the natural monomorphism
    $N \cong 0 \oplus N \hookrightarrow F \oplus N$ defines a
    monomorphism $N \hookrightarrow E$, identifying $N$ with the
    kernel of $E \to M$.
  \item We let $e(\kappa)$ be the extension
    $0 \to N \to E \to M \to 0$ that we have obtained. Prove that
    different choices in the procedure lead to isomorphic extensions.
  \item Prove that $e(\xi(E))$ is isomorphic to $E$ and that
    $\xi(e(\kappa)) = \kappa$.
  \end{itemize}
  The upshot is that there is a natural bijection between
  $\Ext^1_R(M, N)$ and the set of isomorphism classes of extensions of
  $M$ by $N$. Hence the name. `Higher' $\Ext^i_R$ ($i > 1$) may also
  be treated similarly: for example, $\Ext^2_R(M, N)$ can be analyzed
  in terms of two-step extensions consisting of exact sequences
  $0 \to N \to E_1 \to E_2 \to M \to 0$, where two such extensions are
  `isomorphic' if there is a diagram
  \[
    \begin{tikzcd}
      0 \ar[r] & N \ar[r] \ar[d, equal] & E_1 \ar[r] \ar[d] &
      E_2 \ar[r] \ar[d] & M \ar[r] \ar[d, equal] & 0 \\
      0 \ar[r] & N \ar[r] & E_1' \ar[r] & E_2' \ar[r] & M \ar[r] & 0
    \end{tikzcd}
  \]
  This approach to Ext is attributed to \name{Nobuo Yoneda}.
\end{exer}

%%% Local Variables:
%%% mode: LaTeX
%%% TeX-engine: luatex
%%% TeX-master: "../master"
%%% End:
